\documentclass[12pt]{article}

%==================== PACKAGES ====================
\usepackage[a4paper,margin=0.8in]{geometry}
\usepackage{amsmath,amssymb,amsthm}
\usepackage{enumitem}
\usepackage{bm}
\usepackage{hyperref}
\usepackage{fancyhdr}
\usepackage{xcolor}
\usepackage{array}
\usepackage{booktabs}
\usepackage{tikz}
\usepackage{pgfplots}
\pgfplotsset{compat=1.17}
\usepackage[most]{tcolorbox}
\tcbuselibrary{theorems,skins,breakable}
\usepackage{graphicx}
\usepackage{multicol}
\usepackage{draftwatermark}

%==================================================
% COURSE METADATA (EDIT THESE FOR EACH MODULE)
%==================================================
\newcommand{\CourseCode}{HE2001}
\newcommand{\CourseName}{Microeconomics II}
\newcommand{\DocType}{Revision Notes}
\newcommand{\AcademicYear}{Academic Year 2025--2026}

\title{
\Huge \CourseCode\ \CourseName \\[0.3em]
\Large \AcademicYear
}
\author{
  \textit{\href{https://qrsntu.org}{Quantitative Research Society @NTU}}
}
\date{\today}


%==================================================
% TCOLORBOX THEOREM STYLES
%==================================================

% ---------- DEFINITION: dark green, title in darker band ----------
\newtcbtheorem[number within=section]{definition}{Definition}%
{enhanced,
 breakable,
 colback=green!5!white,        % light interior
 colframe=green!40!black,      % border
 colbacktitle=green!50!black,  % dark title band
 coltitle=white,               % title text colour
 fonttitle=\bfseries,
 title filled,                 % puts title inside the band
 boxed title style={
   size=small,
   top=2pt, bottom=2pt,
   left=6pt, right=6pt,
   arc=1pt, outer arc=1pt
 },
 boxrule=0.9pt,
 arc=1pt,
 left=8pt, right=8pt, top=10pt, bottom=8pt,
 before skip=12pt, after skip=12pt
}{def}

% ---------- THEOREM: blue, title in darker band ----------
\newtcbtheorem[number within=section]{theorem}{Theorem}%
{enhanced,
 breakable,
 colback=blue!3!white,
 colframe=blue!40!black,
 colbacktitle=blue!55!black,
 coltitle=white,
 fonttitle=\bfseries,
 title filled,
 boxed title style={
   size=small,
   top=2pt, bottom=2pt,
   left=6pt, right=6pt,
   arc=1pt, outer arc=1pt
 },
 boxrule=0.9pt,
 arc=1pt,
 left=8pt, right=8pt, top=10pt, bottom=8pt,
 before skip=12pt, after skip=12pt
}{thm}

% ---------- ENVIRONMENTS USING TCOLORBOXENVIRONMENT ----------
\tcolorboxenvironment{identity}{
  enhanced,
  breakable,
  colback=green!5!white,
  colframe=green!60!black,
  fonttitle=\bfseries,
  boxrule=1pt,
  arc=3pt,
  left=8pt, right=8pt, top=8pt, bottom=8pt,
  before skip=12pt, after skip=12pt
}

\tcolorboxenvironment{principle}{
  enhanced,
  breakable,
  colback=red!5!white,
  colframe=red!60!black,
  fonttitle=\bfseries,
  boxrule=1pt,
  arc=3pt,
  left=8pt, right=8pt, top=8pt, bottom=8pt,
  before skip=12pt, after skip=12pt
}

%==================================================
% CUSTOM TCOLORBOX ENVIRONMENTS (WITH NAMES)
%==================================================
% Optional argument [1][] lets you override the default title if needed:
%   \begin{keyformula}[title={Key Formula: Demand Curve}] ... \end{keyformula}

\newtcolorbox{keyformula}[1][]{
    enhanced,
    breakable,
    colback=blue!5!white,
    colframe=blue!75!black,
    title=Key Formula,
    fonttitle=\bfseries,
    boxrule=0.9pt,
    arc=2pt,
    left=8pt, right=8pt, top=8pt, bottom=8pt,
    #1
}

% *** CHANGED TO BROWN AS REQUESTED ***
\newtcolorbox{examplebox}[1][]{
    enhanced,
    breakable,
    colback=brown!5!white,
    colframe=brown!75!black,
    title=Worked Example,
    fonttitle=\bfseries,
    boxrule=0.9pt,
    arc=2pt,
    left=8pt, right=8pt, top=8pt, bottom=8pt,
    #1
}

\newtcolorbox{commonmistake}[1][]{
    enhanced,
    breakable,
    colback=red!5!white,
    colframe=red!75!black,
    title=Common Mistake,
    fonttitle=\bfseries,
    boxrule=0.9pt,
    arc=2pt,
    left=8pt, right=8pt, top=8pt, bottom=8pt,
    #1
}

\newtcolorbox{policynote}[1][]{
    enhanced,
    breakable,
    colback=yellow!10!white,
    colframe=orange!75!black,
    title=Policy Application,
    fonttitle=\bfseries,
    boxrule=0.9pt,
    arc=2pt,
    left=8pt, right=8pt, top=8pt, bottom=8pt,
    #1
}

\newtcolorbox{examtip}[1][]{
    enhanced,
    breakable,
    colback=purple!5!white,
    colframe=purple!75!black,
    title=Exam Focus,
    fonttitle=\bfseries,
    boxrule=0.9pt,
    arc=2pt,
    left=8pt, right=8pt, top=8pt, bottom=8pt,
    #1
}

%==================================================
% HEADER / FOOTER / WATERMARK
%==================================================



\pagestyle{fancy}
\fancyhf{}
\fancyhead[L]{\CourseCode\ \CourseName{} -- \DocType}
\fancyhead[R]{\href{https://qrsntu.org}{QRS@NTU}}
\fancyfoot[C]{\thepage}
\renewcommand{\headrulewidth}{0.4pt}
\renewcommand{\footrulewidth}{0pt}

% Watermark
\SetWatermarkText{QRS@NTU — qrsntu.org}
\SetWatermarkScale{1}
\SetWatermarkLightness{0.99}

%------------- HEADER & FOOTER (for page 2 onward) -------------
\fancypagestyle{main}{
  \fancyhf{}
  \fancyhead[L]{\CourseCode\ \CourseName{} -- \DocType}
  \fancyhead[R]{\href{https://qrsntu.org}{QRS@NTU}}
  \fancyfoot[C]{\thepage}
  \renewcommand{\headrulewidth}{0.4pt}
  \renewcommand{\footrulewidth}{0pt}
}

%------------- SIMPLE FOOTER (page 1 only) -------------
\fancypagestyle{plainfirst}{
  \fancyhf{}
  \renewcommand{\headrulewidth}{0pt}
  \renewcommand{\footrulewidth}{0pt}
  \fancyfoot[C]{\scriptsize\textcolor{gray}{\textit{For academic reference only.
 This document is provided for educational purposes and does not constitute an 
 official question paper, answer key, or any formally authorised academic 
 material. Its content is not guaranteed to be complete or accurate.}}}
}

%------------- HYPERREF METADATA -------------
\hypersetup{
    hidelinks,
    pdftitle={\CourseCode\ \CourseName{} -- \DocType},
    pdfauthor={Quantitative Research Society, NTU},
    pdfsubject={\CourseName},
    pdfkeywords={microeconomics, QRS@NTU, revision notes}
}

%==================================================
% DOCUMENT
%==================================================
\begin{document}
\maketitle
\thispagestyle{plainfirst}

\begin{abstract}
\noindent
Comprehensive revision notes for \CourseCode\ \CourseName, covering consumer choice, Slutsky decomposition, endowments, intertemporal choice, uncertainty, welfare measurement, competitive equilibrium, social choice, and utilitarian welfare. The chapters combine core theory, key formulas, worked examples, and exam-focused synthesis.
\end{abstract}

\vspace{0.5cm}
\tableofcontents
\newpage

\newgeometry{left=0.7in,right=0.7in,top=1in,bottom=1in}
\fontsize{10.5pt}{11.5pt}\selectfont
\pagestyle{main}

%====================================================================
\section*{Course Overview \& Topic Map}
\addcontentsline{toc}{section}{Course Overview \& Topic Map}
%====================================================================

\begin{center}
\renewcommand{\arraystretch}{1.2}
\begin{tabular}{|p{3.0cm}|p{2.0cm}|p{2.8cm}|p{2.0cm}|p{4.8cm}|}
\hline
\textbf{Topic Area} & \textbf{Course Part} & \textbf{Related Lectures} & \textbf{Tutorials} & \textbf{Role in Notes} \\
\hline
Chapter 1: Consumer Theory Foundations & Part I & Lecture 1 & Tutorial 1 & Full theory chapter on mathematical setup, consumption, preferences, utility, and revealed preference. \\
\hline
Chapter 2 & Part I & Lecture 2 & Tutorial 2 & Full theory chapter. \\
\hline
Chapter 3 & Part I & Lecture 3 & Tutorial 3 & Full theory chapter. \\
\hline
Chapter 4 & Part I & Lecture 4 & Tutorial 4 & Full theory chapter. \\
\hline
Chapter 5 & Part I & Lecture 5 & Tutorial 5 & Full theory chapter. \\
\hline
Chapter 6: Part I Summary & Part I & Lecture 6 & Tutorial 6 & Synthesis chapter summarising and connecting Part I material rather than introducing a normal new-content chapter. \\
\hline
Chapter 7 & Part II & Lecture 7 & Tutorial 7 & Full theory chapter. \\
\hline
Chapter 8 & Part II & Lecture 8 & Tutorial 8 & Full theory chapter. \\
\hline
Chapter 9 & Part II & Lecture 9 & Tutorial 9 & Full theory chapter. \\
\hline
Chapter 10 & Part II & Lecture 10 & Tutorial 10 & Full theory chapter. \\
\hline
Chapter 11 & Part II & Lecture 11 & --- & Lecture-only chapter. \\
\hline
Chapter 12: Part II Summary & Part II & Lectures 7--11 & Tutorials 7--10 & Synthesis chapter consolidating the second half of the course, including Lecture 11 as lecture-only material. \\
\hline
\end{tabular}
\end{center}

\noindent
``This revision guide is organized into chapter notes plus part-level synthesis chapters, preserving lecture traceability while improving revision efficiency.''

\newpage

%====================================================================
\section*{How to Use This Document}
\addcontentsline{toc}{section}{How to Use This Document}
%====================================================================

Lectures provide the core theory, definitions, models, and examinable derivations. Tutorials are used selectively as worked examples only when they sharpen technique or clarify a central idea. The summary chapters are designed to compress, compare, and connect topics within each course part rather than to function as ordinary theory chapters.

\subsection*{Box Legend}
\begin{itemize}
    \item \textbf{Definition / Theorem} boxes state formal concepts and results.
    \item \textbf{Key Formula} boxes collect expressions, identities, and compact derivations worth memorising.
    \item \textbf{Worked Example} boxes show representative tutorial-style applications.
    \item \textbf{Common Mistake} boxes flag errors that often appear under exam pressure.
    \item \textbf{Policy Application} boxes connect the theory to interpretation and welfare reasoning.
    \item \textbf{Exam Focus} boxes identify the highest-yield revision targets.
\end{itemize}

\subsection*{Suggested Revision Flow}
\begin{enumerate}
    \item Read the chapter prose and Definition/Theorem boxes to fix the conceptual structure.
    \item Memorise the Key Formula boxes and reconstruct the algebra without notes.
    \item Attempt each Worked Example before reading the solution.
    \item Use the summary chapters for cross-topic comparison and final consolidation.
    \item Review Common Mistake and Exam Focus boxes immediately before quizzes and tests.
\end{enumerate}

\newpage
%==================================================
\part{Main Notes}
%==================================================

\section{Consumer Theory Foundations}

This chapter establishes the formal language used throughout \textsc{HE2001}. The lecture begins with minimal mathematical setup, then builds the consumption model, formalises preferences, and finally shows how observed choices can be used to recover information about welfare through revealed preference. Conceptually, the chapter answers three basic questions:
\begin{enumerate}
    \item What objects describe a consumer's choice problem?
    \item What properties make a binary relation behave like a rational preference?
    \item What restrictions must observed choices satisfy if they come from maximisation of an increasing preference?
\end{enumerate}

\subsection{Mathematical Setup}

The course treats a consumer as choosing a vector of goods from a feasible set. That requires only a small amount of notation.

\begin{definition}{Basic sets, vectors, and functions}{ch1-basic-math}
For an integer $L \ge 2$:
\begin{itemize}
    \item $\mathbb{R}$ is the set of real numbers.
    \item $\mathbb{R}_+ = \{x \in \mathbb{R} : x \ge 0\}$ and $\mathbb{R}_{++} = \{x \in \mathbb{R} : x > 0\}$.
    \item $\mathbb{R}^L_+$ is the set of all $L$-dimensional nonnegative vectors.
    \item $\mathbb{R}^L_{++}$ is the set of all $L$-dimensional strictly positive vectors.
    \item A function $f : X \to Y$ maps each element of a set $X$ to an element of a set $Y$.
\end{itemize}
In this course, typical functions include utility functions, demand functions, and cost functions.
\end{definition}

A set records the objects under discussion; a vector records several quantities at once; a function translates one kind of economic object into another. These are the only preliminaries needed for Lecture 1.

\begin{keyformula}[title={Key Formula: Dot Product and Expenditure}]
If $p = (p_1,\dots,p_L) \in \mathbb{R}^L_{++}$ and $x = (x_1,\dots,x_L) \in \mathbb{R}^L_+$, then the dot product is
\[
p \cdot x = \sum_{\ell=1}^{L} p_\ell x_\ell.
\]
Economically, $p \cdot x$ is the total expenditure required to purchase bundle $x$ at prices $p$.
\end{keyformula}

\subsection{Consumption Bundles, Prices, and Budget Sets}

The primitive objects of consumer theory are the consumption bundle, the price vector, and the consumer's expenditure level.

\begin{definition}{Consumption bundle and price vector}{ch1-consumption-bundle}
Suppose there are $L$ goods.
\begin{itemize}
    \item A \emph{consumption bundle} is a vector $x = (x_1,\dots,x_L) \in \mathbb{R}^L_+$, where $x_\ell$ is the amount of good $\ell$.
    \item A \emph{price vector} is a vector $p = (p_1,\dots,p_L) \in \mathbb{R}^L_{++}$, where $p_\ell$ is the price of good $\ell$.
\end{itemize}
\end{definition}

The restriction $x \in \mathbb{R}^L_+$ says negative consumption is excluded. The restriction $p \in \mathbb{R}^L_{++}$ says every good has a strictly positive price, which matters later when we convert unused expenditure into more of some good.

\begin{definition}{Budget set and budget line}{ch1-budget-set}
If the consumer has expenditure level $m > 0$ and faces prices $p \in \mathbb{R}^L_{++}$, the \emph{budget set} is
\[
B(p,m) = \{x \in \mathbb{R}^L_+ : p \cdot x \le m\}.
\]
In the two-good case, the \emph{budget line} is the boundary
\[
p_1 x_1 + p_2 x_2 = m.
\]
\end{definition}

The budget set contains every affordable bundle. The budget line contains bundles that exactly exhaust expenditure. In two dimensions, the budget set is the region on or below the line in the nonnegative quadrant.

\begin{keyformula}[title={Key Formula: Two-Good Budget Geometry}]
For $L=2$, the budget line is
\[
p_1x_1 + p_2x_2 = m.
\]
Its intercepts are
\[
x_1 = \frac{m}{p_1} \quad \text{when } x_2=0,
\qquad
x_2 = \frac{m}{p_2} \quad \text{when } x_1=0.
\]
The slope is
\[
-\frac{p_1}{p_2}.
\]
Thus relative prices determine the steepness of the line, while $m$ shifts the feasible set outward or inward.
\end{keyformula}

\begin{theorem}{Why increasing preferences push choice to the boundary}{ch1-boundary-choice}
Suppose preferences are increasing and the consumer chooses $x \in B(p,m)$. Then any optimal bundle must satisfy
\[
p \cdot x = m.
\]
\end{theorem}

\begin{proof}
If instead $p \cdot x < m$, then because each price is strictly positive, there exists some good $i$ and some $\varepsilon > 0$ such that $p_i\varepsilon \le m - p \cdot x$. The bundle $y = x + \varepsilon e_i$ is still affordable and satisfies $y > x$. Increasing preferences imply $y \succ x$, contradicting optimality of $x$.
\end{proof}

\begin{examplebox}
\textbf{Budget affordability and the budget line.}

Take two goods, expenditure $m=2$, and consider the two price vectors
\[
p^{(1)}=(1,2), \qquad p^{(2)}=(2,1).
\]
The corresponding budget lines are
\[
x_1 + 2x_2 = 2, \qquad 2x_1 + x_2 = 2.
\]
Now test the bundle $x=(1,1)$.
\[
p^{(1)} \cdot x = 1(1)+2(1)=3 > 2,
\qquad
p^{(2)} \cdot x = 2(1)+1(1)=3 > 2.
\]
Therefore $x=(1,1)$ is not affordable under either price system.

This is the fastest algebraic way to check feasibility in an exam: compute $p \cdot x$ and compare it with $m$.
\end{examplebox}

\begin{commonmistake}
Do not confuse the \emph{budget set} with the \emph{budget line}. The line is only the boundary $p\cdot x = m$; the set is the full collection of affordable bundles $p\cdot x \le m$ together with nonnegativity.
\end{commonmistake}

\subsection{Preferences}

A consumer theory model needs a way to rank bundles. Lecture 1 uses a binary relation to represent this ranking.

\begin{definition}{Preference relation}{ch1-preference-relation}
A binary relation $\succeq$ on $\mathbb{R}^L_+$ assigns to each pair $(x,x')$ a statement of whether $x$ is weakly preferred to $x'$. The notation means:
\[
x \succeq x' \quad \text{``$x$ is weakly preferred to $x'$.''}
\]
From $\succeq$ we derive:
\begin{itemize}
    \item \emph{strict preference}: $x \succ x'$ means $x \succeq x'$ and not $x' \succeq x$;
    \item \emph{indifference}: $x \sim x'$ means $x \succeq x'$ and $x' \succeq x$.
\end{itemize}
A binary relation is called a \emph{preference} if it is complete and transitive.
\end{definition}

\begin{definition}{Completeness and transitivity}{ch1-complete-transitive}
A binary relation $\succeq$ on $\mathbb{R}^L_+$ is:
\begin{itemize}
    \item \emph{complete} if for every $x,x'$, either $x \succeq x'$ or $x' \succeq x$ (or both);
    \item \emph{transitive} if $x \succeq x'$ and $x' \succeq x''$ imply $x \succeq x''$.
\end{itemize}
\end{definition}

Completeness says the consumer can compare any two bundles. Transitivity says the ranking is logically consistent across chains of comparisons. Without completeness, some bundles may be incomparable; without transitivity, circular rankings can arise.

\begin{definition}{Increasing preferences}{ch1-increasing}
For bundles $x,x' \in \mathbb{R}^L_+$:
\begin{align*}
x \ge x' &\iff x_\ell \ge x'_\ell \text{ for all } \ell, \\
x > x' &\iff x \ge x' \text{ and } x_\ell > x'_\ell \text{ for some } \ell.
\end{align*}
A preference relation $\succeq$ is \emph{increasing} if
\[
x > x' \implies x \succ x'.
\]
\end{definition}

Increasingness is the formal version of ``more is better.'' If bundle $x$ gives at least as much of every good and strictly more of at least one good, then it must be strictly preferred.

\begin{examplebox}
\textbf{Checking completeness, transitivity, and increasingness.}

Consider the relation on $\mathbb{R}^2_+$ defined by
\[
(x_1,x_2) \succeq (y_1,y_2)
\quad \Longleftrightarrow \quad
\text{either } x_1 > y_1, \text{ or } x_1 = y_1 \text{ and } x_2 \ge y_2.
\]
This is a lexicographic-type ranking.

\begin{itemize}
    \item \textbf{Complete:} for any two bundles, either the first coordinate differs, or the first coordinates are equal and the second coordinate decides the comparison.
    \item \textbf{Transitive:} if $x$ beats $y$ by the first coordinate, and $y$ beats $z$, then $x$ also beats $z$; if first coordinates tie, the second coordinate preserves the ordering.
    \item \textbf{Increasing:} whenever $x \ge y$ and $x \ne y$, either $x_1 > y_1$, or $x_1 = y_1$ and $x_2 > y_2$, so $x \succ y$.
\end{itemize}

Hence this relation is complete, transitive, and increasing.
\end{examplebox}

\begin{commonmistake}
Increasing does \emph{not} mean ``strictly more of every good.'' The condition $x > x'$ only requires weakly more of every good and strictly more of at least one good.
\end{commonmistake}

\subsection{Utility Representation}

Economists often work with utility functions because optimisation and algebra are easier to express with numbers than with a raw binary relation.

\begin{definition}{Utility representation}{ch1-utility-representation}
A utility function is a map $U : \mathbb{R}^L_+ \to \mathbb{R}$. It \emph{represents} a preference relation $\succeq$ if
\[
x \succeq x' \quad \Longleftrightarrow \quad U(x) \ge U(x').
\]
If $U$ is increasing in the sense that $x > x'$ implies $U(x) > U(x')$, then the induced preference relation is increasing.
\end{definition}

Utility is therefore a numerical encoding of a ranking. In Lecture 1, the important point is representational: utility is useful because maximising $U(x)$ is equivalent to choosing the most preferred bundle whenever $U$ represents the underlying preference.

\begin{examplebox}
\textbf{A utility function that generates an increasing preference.}

Let
\[
U(x_1,x_2)=\sqrt{x_1}+\sqrt{x_2}.
\]
Then for any two bundles $x$ and $y$,
\[
x \succeq y \quad \Longleftrightarrow \quad U(x) \ge U(y).
\]
Because the square-root function is increasing on $\mathbb{R}_+$, if $x > y$ then
\[
\sqrt{x_1} \ge \sqrt{y_1}, \qquad \sqrt{x_2} \ge \sqrt{y_2},
\]
with strict inequality in at least one coordinate. Hence
\[
U(x)=\sqrt{x_1}+\sqrt{x_2} > \sqrt{y_1}+\sqrt{y_2}=U(y).
\]
So the represented preference is increasing, complete, and transitive.
\end{examplebox}

\begin{examtip}
When asked whether a utility function represents preferences, the core test is always the ranking statement
\[
x \succeq y \iff U(x) \ge U(y).
\]
Do not interpret utility levels cardinally unless the question explicitly introduces that structure. In this chapter, utility is used only as a ranking device.
\end{examtip}

\subsection{Revealed Preference}

The lecture then asks a more empirical question: if we only observe prices and chosen bundles, what can we infer about preferences?

\begin{examtip}[title={Core Principles Behind Revealed Preference}]
Classic consumer theory in Lecture 1 relies on three connected ideas:
\begin{enumerate}
    \item What is chosen is preferred to what could have been chosen.
    \item More is preferred.
    \item What is preferred is best from the consumer's own point of view.
\end{enumerate}
Together, these principles allow observed choices to reveal welfare information.
\end{examtip}

Suppose an observation consists of prices $p^t \in \mathbb{R}^L_{++}$ and a chosen bundle $x^t \in \mathbb{R}^L_+$. Let expenditure in that observation be $m^t = p^t \cdot x^t$.

\begin{definition}{Directly and strictly directly revealed preference}{ch1-revealed-pref}
Given an observation $(p^t,x^t)$:
\begin{itemize}
    \item $x^t$ is \emph{directly revealed preferred} to $x$, written $x^t R x$, if
    \[
    p^t \cdot x \le p^t \cdot x^t.
    \]
    That is, bundle $x$ was affordable when $x^t$ was chosen.
    \item $x^t$ is \emph{strictly directly revealed preferred} to $x$, written $x^t P x$, if
    \[
    p^t \cdot x < p^t \cdot x^t.
    \]
    That is, bundle $x$ was strictly cheaper than the chosen bundle.
\end{itemize}
\end{definition}

Economically, $R$ is an observed revealed ranking and $P$ is its strict version. The distinction matters because strict affordability combined with increasingness is strong enough to generate strict welfare conclusions.

\begin{theorem}{From revealed preference to actual preference}{ch1-rp-actual-pref}
Suppose the consumer chooses each observed bundle by maximising an increasing preference relation $\succeq$.
\begin{enumerate}
    \item If $x R y$, then $x \succeq y$.
    \item If $x P y$, then $x \succ y$.
\end{enumerate}
\end{theorem}

\begin{proof}
If $x R y$, then $y$ was affordable when $x$ was chosen. By optimality of $x$, we must have $x \succeq y$.

If $x P y$, then $y$ was strictly cheaper than $x$ at the prices where $x$ was chosen. Because prices are strictly positive, some bundle $z$ exists with $z > y$ and still affordable at those prices. Increasingness implies $z \succ y$. Since $x$ was chosen when $z$ was affordable, optimality gives $x \succeq z$. By transitivity, $x \succ y$.
\end{proof}

\begin{examplebox}
\textbf{Comparing welfare across two periods.}

Suppose a consumer chooses
\[
x^1=(2,2) \text{ at } p^1=(1,1),
\qquad
x^2=(0,3) \text{ at } p^2=(2,1).
\]
Compute expenditures:
\[
p^1 \cdot x^1 = 4,
\qquad
p^2 \cdot x^2 = 3.
\]
Now check whether period-2 bundle $x^2$ was affordable in period 1:
\[
p^1 \cdot x^2 = 1(0)+1(3)=3 \le 4.
\]
So $x^1 R x^2$, which implies
\[
x^1 \succeq x^2.
\]
Next check whether $x^1$ was affordable in period 2:
\[
p^2 \cdot x^1 = 2(2)+1(2)=6 > 3.
\]
So period 2 does \emph{not} reveal $x^2 \succeq x^1$.

Hence the data support the welfare conclusion
\[
\boxed{x^1 \succeq x^2.}
\]
The consumer is weakly better off in period 1.
\end{examplebox}

\begin{policynote}
Revealed preference gives a practical welfare test after policy changes. If an old bundle remains affordable after a tax or price change but the consumer switches away from it, the new choice reveals at least weak improvement relative to the old affordable alternatives. Conversely, if the new bundle was affordable before the change but the old bundle was chosen instead, the pre-change situation is revealed to be weakly better. This is the logic behind welfare comparisons from observed demand rather than from directly measured happiness.
\end{policynote}

\subsection{WARP}

Revealed preference yields testable restrictions. The main one in Lecture 1 is WARP.

\begin{definition}{Weak Axiom of Revealed Preference (WARP)}{ch1-warp-def}
A dataset satisfies \emph{WARP} if
\[
x R y \implies \text{not } (y P x).
\]
Equivalently: once $x$ is chosen when $y$ is affordable, it cannot happen that $y$ is later chosen from a budget where $x$ is strictly cheaper.
\end{definition}

WARP is a consistency condition on observable choices. It does not require knowledge of the utility function itself. It only checks whether observed decisions could have come from maximisation of an increasing preference.

\begin{theorem}{Increasing preference maximisation implies WARP}{ch1-warp-thm}
If each observed choice maximises an increasing preference relation $\succeq$, then the resulting data satisfy WARP.
\end{theorem}

\begin{proof}
Assume $x R y$. By Theorem \ref{thm:ch1-rp-actual-pref}, this implies $x \succeq y$.
Suppose for contradiction that $y P x$ also holds. Then Theorem \ref{thm:ch1-rp-actual-pref} implies $y \succ x$. But $y \succ x$ means it is not true that $x \succeq y$, contradicting $x \succeq y$. Therefore $y P x$ cannot occur, so WARP must hold.
\end{proof}

\begin{examplebox}
\textbf{A clean WARP violation.}

Observation 1: prices are $p=(2,1)$ and the consumer chooses
\[
x=(1,0).
\]
Observation 2: prices are $p'=(1,2)$ and the consumer chooses
\[
x'=(0,1).
\]
Now compare costs crosswise:
\[
p \cdot x = 2,
\qquad
p \cdot x' = 1.
\]
So under prices $p$, bundle $x'$ is strictly cheaper than the chosen bundle $x$, hence
\[
x P x'.
\]
Similarly,
\[
p' \cdot x' = 2,
\qquad
p' \cdot x = 1,
\]
so under prices $p'$, bundle $x$ is strictly cheaper than the chosen bundle $x'$, hence
\[
x' P x.
\]
Thus both strict revealed rankings hold in opposite directions, which violates WARP.

Conclusion: these two observations cannot be rationalised by maximisation of an increasing preference.
\end{examplebox}

\begin{commonmistake}
Do not say that $x R y$ and $y R x$ automatically violate WARP. Mutual \emph{weak} direct revealed preference can occur when each bundle is merely affordable at the other observation. WARP is violated only when the reverse comparison is \emph{strict}, i.e. when $y P x$ also holds.
\end{commonmistake}

\begin{examtip}
For any revealed preference question, use the same checklist:
\begin{enumerate}
    \item Write down each observed expenditure $m^t = p^t \cdot x^t$.
    \item For every pair of bundles, test affordability by computing cross-costs $p^t \cdot x^s$.
    \item Convert affordability into $R$ or $P$ statements.
    \item Use $R \Rightarrow \succeq$ and $P \Rightarrow \succ$ only when the choice is interpreted as maximisation of an increasing preference.
    \item Check WARP by ruling out the pattern $x R y$ together with $y P x$.
\end{enumerate}
\end{examtip}

\subsection{Chapter Summary}

\begin{keyformula}[title={Key Formula: Chapter 1 Summary}]
\begin{align*}
B(p,m) &= \{x \in \mathbb{R}^L_+ : p \cdot x \le m\}, \\
x R y &\iff y \text{ was affordable when } x \text{ was chosen}, \\
x P y &\iff y \text{ was strictly cheaper when } x \text{ was chosen}, \\
\text{WARP:} \quad x R y &\implies \text{not } (y P x).
\end{align*}
\end{keyformula}

\begin{examtip}
High-yield takeaways from Lecture 1:
\begin{itemize}
    \item Know the definitions of completeness, transitivity, increasingness, strict preference, and indifference precisely.
    \item Be able to move fluently between verbal statements and formal notation.
    \item Always interpret $p \cdot x$ as expenditure.
    \item For welfare comparisons across periods, affordability is the central test.
    \item WARP is the first observable implication of rational choice with increasing preferences.
\end{itemize}
\end{examtip}
\newpage
% chapters/ch02_slutsky_equation_and_demand_response.tex

\section{Slutsky Equation and Demand Response}

Lecture 2 extends the revealed-preference framework from Chapter 1 and then turns to comparative statics of consumer demand. The chapter has two main aims. First, it strengthens WARP into a multi-observation consistency condition, GARP, which uses the transitivity of preferences. Second, it studies how Marshallian demand responds to changes in prices and expenditure, culminating in the Slutsky decomposition of a price effect into substitution and income effects.

\subsection{Generalized Revealed Preference}

Recall from Chapter 1 that if a bundle $x$ is chosen at prices $p$, then for any bundle $x'$:
\[
x R x' \iff p \cdot x \ge p \cdot x',
\qquad
x P x' \iff p \cdot x > p \cdot x'.
\]
If choices come from maximization of an increasing preference relation $\succeq$, then
\[
x R x' \implies x \succeq x',
\qquad
x P x' \implies x \succ x'.
\]

Because preferences are transitive, revealed preference should also propagate through chains of observations.

\begin{definition}{Revealed and strictly revealed preference through chains}{ch2-rstar}
Let $R$ and $P$ denote direct and strict direct revealed preference.

\begin{itemize}
    \item $x$ is \emph{revealed preferred} to $x'$, written $x R^* x'$, if there exist bundles
    \[
    x^1,\dots,x^N
    \]
    such that
    \[
    x R x^1 R x^2 R \cdots R x^N R x'.
    \]
    \item $x$ is \emph{strictly revealed preferred} to $x'$, written $x P^* x'$, if such a chain exists and at least one link is strict, i.e. at least one $R$ can be replaced by $P$.
\end{itemize}
\end{definition}

\begin{theorem}{Behavioral meaning of $R^*$ and $P^*$}{ch2-rstar-meaning}
Suppose observed choices maximize an increasing preference relation $\succeq$. Then
\[
x R^* x' \implies x \succeq x',
\qquad
x P^* x' \implies x \succ x'.
\]
\end{theorem}

\begin{proof}
If
\[
x R x^1 R x^2 R \cdots R x^N R x',
\]
then each link implies a weak preference comparison under utility maximization:
\[
x \succeq x^1,\quad x^1 \succeq x^2,\quad \dots,\quad x^N \succeq x'.
\]
By transitivity, $x \succeq x'$.

If at least one link is strict, then at least one strict preference comparison appears in the chain. Combining weak and strict comparisons by transitivity yields $x \succ x'$.
\end{proof}

\begin{definition}{Generalized Axiom of Revealed Preference (GARP)}{ch2-garp}
A dataset satisfies \emph{GARP} if
\[
x R^* x' \implies \text{not }(x' P x).
\]
Equivalently: once $x$ is revealed preferred to $x'$ through a chain of affordable choices, it cannot be that $x'$ is later chosen when $x$ was strictly cheaper.
\end{definition}

GARP is stronger than WARP because it checks for contradictions generated not only by one-step comparisons, but also by multi-step revealed-preference chains.

\begin{theorem}{Afriat's theorem, as used in this course}{ch2-afriat}
For a finite dataset of consumer choices:
\begin{enumerate}
    \item If the data are generated by maximization of an increasing utility function, then the data satisfy GARP.
    \item If the data satisfy GARP, then there exists an increasing utility function that rationalizes the data.
\end{enumerate}
Hence GARP is the correct revealed-preference test for rationalizability by an increasing utility function.
\end{theorem}

\begin{examtip}
In two-good settings, Lecture 2 notes that WARP and GARP are equivalent. In three or more goods, GARP is strictly stronger because cycles can arise through multiple observations even when no strict two-observation cycle exists.
\end{examtip}

\begin{examplebox}
\textbf{WARP can hold while GARP fails in three goods.}

Consider three observations:
\[
x^1=(1,0,0),\qquad x^2=(0,1,0),\qquad x^3=(0,0,1),
\]
at prices
\[
p^1=(2,4,1),\qquad p^2=(1,2,4),\qquad p^3=(4,1,2).
\]
Since each chosen bundle is a unit vector,
\[
p^1 \cdot x^1=2,\qquad p^2 \cdot x^2=2,\qquad p^3 \cdot x^3=2.
\]
Now compute cross-costs:
\[
p^1 \cdot x^3 = 1 \le 2 \implies x^1 R x^3,
\]
\[
p^3 \cdot x^2 = 1 \le 2 \implies x^3 R x^2,
\]
\[
p^2 \cdot x^1 = 1 \le 2 \implies x^2 R x^1.
\]
So
\[
x^1 R x^3 R x^2 R x^1,
\]
which gives a revealed-preference cycle. Moreover,
\[
p^2 \cdot x^1 = 1 < 2 = p^2 \cdot x^2,
\]
so in fact
\[
x^2 P x^1.
\]
Thus GARP is violated because $x^1 R^* x^1$ through a chain that ends with a strict reversal.

However, there is no strict \emph{two-observation} cycle, so WARP is not violated. This shows that GARP is strictly stronger than WARP in dimensions $L \ge 3$.
\end{examplebox}

\begin{commonmistake}
Do not test GARP by looking only for pairs $x R y$ and $y P x$. That checks WARP. GARP requires checking whether a contradiction appears \emph{after a chain} of revealed-preference comparisons.
\end{commonmistake}

\subsection{Budgets, Utility Maximization, and Demand}

The second part of the lecture returns to standard consumer theory. The two core objects are the utility function and the budget set.

\begin{definition}{Indifference curve and demand function}{ch2-demand-def}
Let $U : \mathbb{R}^L_+ \to \mathbb{R}$ be a utility function.

\begin{itemize}
    \item An \emph{indifference curve} is a set of bundles with the same utility level:
    \[
    \{x \in \mathbb{R}^L_+ : U(x)=\bar u\}.
    \]
    \item A \emph{demand function} is a map
    \[
    x : \mathbb{R}^L_{++} \times \mathbb{R}_{++} \to \mathbb{R}^L_+
    \]
    such that $x(p,m)$ is affordable and utility-maximizing at prices $p$ and expenditure $m$:
    \[
    p \cdot x(p,m) \le m,
    \qquad
    U(x(p,m)) \ge U(\tilde x)\ \ \forall \tilde x \in B(p,m).
    \]
\end{itemize}
\end{definition}

The demand function gives the consumer's optimal bundle as a function of prices and expenditure. Its $\ell$-th component, $x_\ell(p,m)$, is demand for good $\ell$.

\begin{theorem}{Exhausting the budget under increasing utility}{ch2-budget-exhaustion}
If $U$ is increasing and $x(p,m)$ maximizes $U$ over $B(p,m)$, then
\[
p \cdot x(p,m)=m.
\]
\end{theorem}

\begin{proof}
If $p \cdot x(p,m) < m$, then some strictly positive amount of some good can be added while remaining in the budget set because each price is strictly positive. The resulting bundle is strictly larger componentwise, so increasing utility implies it yields higher utility, contradicting optimality.
\end{proof}

\begin{examtip}
When the question gives a Marshallian demand function $x(p,m)$, always verify two things first:
\begin{enumerate}
    \item affordability: $p \cdot x(p,m) \le m$;
    \item if utility is increasing, budget exhaustion: $p \cdot x(p,m)=m$.
\end{enumerate}
These are the fastest internal consistency checks.
\end{examtip}

\subsection{Demand Response to Price Changes}

Let $x(p,m)=(x_1(p,m),x_2(p,m))$ denote Marshallian demand in the two-good case. The lecture studies what happens when the price of good 1 increases from $p_1$ to $p_1' > p_1$, holding $p_2$ and $m$ fixed.

\begin{definition}{Gross substitutes, gross complements, ordinary goods, and Giffen goods}{ch2-gross-def}
Suppose the price of good 1 rises from $p_1$ to $p_1'$.

\begin{itemize}
    \item If
    \[
    x_2(p_1',p_2,m) > x_2(p_1,p_2,m),
    \]
    then goods 1 and 2 are \emph{gross substitutes}.
    \item If
    \[
    x_2(p_1',p_2,m) < x_2(p_1,p_2,m),
    \]
    then goods 1 and 2 are \emph{gross complements}.
    \item If
    \[
    x_1(p_1',p_2,m) < x_1(p_1,p_2,m),
    \]
    then good 1 is \emph{ordinary}.
    \item If
    \[
    x_1(p_1',p_2,m) > x_1(p_1,p_2,m),
    \]
    then good 1 is \emph{Giffen}.
\end{itemize}
\end{definition}

These are purely Marshallian concepts: they describe the \emph{total} demand response to a price change before any decomposition into substitution and income effects.

\begin{examplebox}
\textbf{An explicit demand system: ordinary own-price effect and positive cross-price effect.}

Suppose
\[
U(x_1,x_2)=\sqrt{x_1}+\sqrt{x_2},
\]
which generates the Marshallian demand
\[
x(p_1,p_2,m)
=
\frac{m}{p_1+p_2}
\left(
\frac{p_2}{p_1},
\frac{p_1}{p_2}
\right).
\]
Hence
\[
x_1(p_1,p_2,m)=\frac{m p_2}{p_1(p_1+p_2)},
\qquad
x_2(p_1,p_2,m)=\frac{m p_1}{p_2(p_1+p_2)}.
\]

For good 2, holding $p_2$ and $m$ fixed,
\[
\frac{\partial x_2}{\partial p_1}
=
\frac{m}{p_2}\cdot \frac{p_2}{(p_1+p_2)^2}
=
\frac{m}{(p_1+p_2)^2}
>0.
\]
So when the price of good 1 rises, demand for good 2 rises: the goods are gross substitutes.

For good 1,
\[
\frac{\partial x_1}{\partial p_1}
=
-\frac{m p_2(2p_1+p_2)}{p_1^2(p_1+p_2)^2}
<0,
\]
so good 1 is ordinary, not Giffen.
\end{examplebox}

\begin{policynote}
The lecture's ``meat and potatoes'' example shows why Giffen behavior is theoretically possible. When a staple good becomes more expensive, the loss of real purchasing power may be so severe that the consumer buys \emph{more} of the staple and cuts back on the more expensive alternative. This is unusual, but it clarifies that the total Marshallian price effect need not always be negative.
\end{policynote}

\subsection{Demand Response to Expenditure Changes}

Now hold prices fixed and increase expenditure from $m$ to $m' > m$.

\begin{definition}{Normal and inferior goods}{ch2-normal-inferior}
A good $\ell$ is:
\begin{itemize}
    \item \emph{normal} if
    \[
    x_\ell(p,m') > x_\ell(p,m)
    \quad \text{whenever } m' > m;
    \]
    \item \emph{inferior} if
    \[
    x_\ell(p,m') < x_\ell(p,m)
    \quad \text{for some increase in expenditure } m' > m.
    \]
\end{itemize}
\end{definition}

Normality and inferiority concern income or expenditure effects, not substitution effects. A good can be ordinary or Giffen in price-response language, and normal or inferior in expenditure-response language.

\begin{examplebox}
\textbf{Normal goods in the square-root demand system.}

Continuing the example
\[
x_1(p_1,p_2,m)=\frac{m p_2}{p_1(p_1+p_2)},
\qquad
x_2(p_1,p_2,m)=\frac{m p_1}{p_2(p_1+p_2)},
\]
we compute
\[
\frac{\partial x_1}{\partial m}=\frac{p_2}{p_1(p_1+p_2)}>0,
\qquad
\frac{\partial x_2}{\partial m}=\frac{p_1}{p_2(p_1+p_2)}>0.
\]
Therefore both goods are normal.
\end{examplebox}

\begin{commonmistake}
Do not infer that a good is normal merely because its demand curve slopes downward. Downward-sloping Marshallian demand describes a price effect; normality describes the sign of the expenditure effect.
\end{commonmistake}

\subsection{Compensated Price Changes and the Slutsky Identity}

The central technical result of Lecture 2 is the decomposition of a price change into substitution and income effects.

Suppose prices change from
\[
p=(p_1,p_2)
\qquad \text{to} \qquad
p'=(p_1',p_2),
\quad p_1' > p_1,
\]
while original expenditure is $m$.

\begin{definition}{Slutsky compensated expenditure}{ch2-ms}
Let $x(p,m)$ be the original Marshallian demand. The \emph{Slutsky compensated expenditure} is
\[
m^s = p' \cdot x(p,m).
\]
This is the amount of money needed at the new prices to just afford the original bundle.
\end{definition}

So the consumer is ``compensated'' enough to remain able to purchase the original bundle after the price increase. This isolates a pure substitution comparison between the old chosen bundle and the new chosen bundle on a budget that passes through the old choice.

\begin{theorem}{The Slutsky substitution effect is nonpositive for the good whose price rises}{ch2-sub-neg}
If the price of good 1 rises from $p_1$ to $p_1' > p_1$, then
\[
x_1(p',m^s) - x_1(p,m) \le 0.
\]
\end{theorem}

\begin{proof}
By construction, the old bundle $x(p,m)$ is affordable at the new prices with compensated expenditure $m^s$ because
\[
p' \cdot x(p,m) = m^s.
\]
Hence $x(p',m^s)$ is chosen from a budget set that contains $x(p,m)$.

Suppose, for contradiction, that
\[
x_1(p',m^s) > x_1(p,m).
\]
Since the price of good 2 is unchanged, compare the old prices:
\begin{align*}
p \cdot x(p',m^s)
&= p_1 x_1(p',m^s) + p_2 x_2(p',m^s) \\
&< p_1' x_1(p',m^s) + p_2 x_2(p',m^s) \\
&= p' \cdot x(p',m^s) \\
&\le m^s \\
&= p' \cdot x(p,m) \\
&= p_1' x_1(p,m) + p_2 x_2(p,m).
\end{align*}
Because $x_1(p',m^s) > x_1(p,m)$ and $p_1' > p_1$, the only way this can hold is if the new choice is cheaper at old prices than the old choice, implying
\[
x(p,m) P x(p',m^s).
\]
But at the new compensated prices, the new choice is chosen while the old one is affordable, so
\[
x(p',m^s) R x(p,m).
\]
This violates WARP, and therefore also GARP, for utility-maximizing behavior. Hence the supposition was false, so
\[
x_1(p',m^s) \le x_1(p,m).
\]
\end{proof}

The lecture packages this into three objects.

\begin{keyformula}[title={Key Formula: Slutsky Decomposition}]
Let
\[
\Delta x_1 = x_1(p',m) - x_1(p,m)
\]
denote the total Marshallian change in demand for good 1 after the price increase. Define
\[
\Delta x_1^s = x_1(p',m^s) - x_1(p,m)
\]
as the \emph{Slutsky substitution effect}, and
\[
\Delta x_1^n = x_1(p',m) - x_1(p',m^s)
\]
as the \emph{Slutsky income effect}. Then
\[
\Delta x_1 = \Delta x_1^s + \Delta x_1^n.
\]
This is the \emph{Slutsky identity}.
\end{keyformula}

The identity is purely algebraic: add and subtract $x_1(p',m^s)$. Its power comes from the fact that the substitution effect has a definite sign, while the income effect depends on whether the good is normal or inferior.

\begin{theorem}{Law of demand for normal goods}{ch2-law-demand}
Suppose good 1 is normal. If its price rises from $p_1$ to $p_1' > p_1$, holding $p_2$ and $m$ fixed, then
\[
x_1(p_1',p_2,m) - x_1(p_1,p_2,m) \le 0.
\]
\end{theorem}

\begin{proof}
First,
\[
m^s = p' \cdot x(p,m)
= p_1' x_1(p,m) + p_2 x_2(p,m)
\ge
p_1 x_1(p,m) + p_2 x_2(p,m)
= m.
\]
So the compensated expenditure is at least as large as the original expenditure.

If good 1 is normal, then at the new prices $p'$ a higher expenditure raises demand for good 1:
\[
m^s \ge m
\implies
x_1(p',m^s) \ge x_1(p',m),
\]
which is equivalent to
\[
x_1(p',m) - x_1(p',m^s) \le 0.
\]
Hence the income effect is nonpositive:
\[
\Delta x_1^n \le 0.
\]
From Theorem \ref{thm:ch2-sub-neg}, the substitution effect also satisfies
\[
\Delta x_1^s \le 0.
\]
Therefore, by the Slutsky identity,
\[
\Delta x_1 = \Delta x_1^s + \Delta x_1^n \le 0.
\]
\end{proof}

\begin{examplebox}
\textbf{Computing the Slutsky effects explicitly.}

Use the demand system
\[
x_1(p_1,p_2,m)=\frac{m p_2}{p_1(p_1+p_2)}.
\]
Take
\[
p=(1,1),\qquad m=2,\qquad p'=(2,1).
\]
Original demand:
\[
x(p,m)=\left(1,1\right).
\]
Compensated expenditure:
\[
m^s = p' \cdot x(p,m)=2(1)+1(1)=3.
\]
Compensated demand for good 1:
\[
x_1(p',m^s)=\frac{3\cdot 1}{2(2+1)}=\frac12.
\]
Uncompensated new demand for good 1:
\[
x_1(p',m)=\frac{2\cdot 1}{2(2+1)}=\frac13.
\]

So:
\[
\Delta x_1^s = \frac12 - 1 = -\frac12,
\]
\[
\Delta x_1^n = \frac13 - \frac12 = -\frac16,
\]
and hence
\[
\Delta x_1
=
\Delta x_1^s + \Delta x_1^n
=
-\frac12 - \frac16
=
-\frac23.
\]

Interpretation:
\begin{itemize}
    \item the substitution effect is negative, as theory predicts;
    \item the income effect is also negative because good 1 is normal;
    \item therefore the total demand response is negative.
\end{itemize}
\end{examplebox}

\begin{examtip}
For price-change questions, keep the order of objects clear:
\begin{enumerate}
    \item old choice $x(p,m)$;
    \item new prices $p'$;
    \item compensated expenditure $m^s = p' \cdot x(p,m)$;
    \item compensated demand $x(p',m^s)$;
    \item uncompensated new demand $x(p',m)$.
\end{enumerate}
Most sign errors come from mixing up $m$ and $m^s$.
\end{examtip}

\subsection{Chapter Summary}

\begin{keyformula}[title={Key Formula: Lecture 2 Summary}]
\begin{align*}
x R^* x' &\iff \text{$x$ is revealed preferred to $x'$ through a chain of direct revealed preferences}, \\
x P^* x' &\iff \text{the chain contains at least one strict link}, \\
\text{GARP:}\quad x R^* x' &\implies \text{not }(x' P x), \\
m^s &= p' \cdot x(p,m), \\
\Delta x_1 &= x_1(p',m)-x_1(p,m), \\
\Delta x_1^s &= x_1(p',m^s)-x_1(p,m), \\
\Delta x_1^n &= x_1(p',m)-x_1(p',m^s), \\
\Delta x_1 &= \Delta x_1^s + \Delta x_1^n.
\end{align*}
If good 1 is normal and its price rises, then both $\Delta x_1^s \le 0$ and $\Delta x_1^n \le 0$, hence
\[
\Delta x_1 \le 0.
\]
\end{keyformula}

\begin{examtip}
The highest-yield ideas from Lecture 2 are:
\begin{itemize}
    \item GARP is the multi-observation revealed-preference test and rationalizability condition.
    \item Marshallian demand describes the \emph{total} response to price and expenditure changes.
    \item Gross substitutes/complements refer to cross-price effects; ordinary/Giffen refer to own-price effects.
    \item Normal/inferior refer to expenditure effects.
    \item The Slutsky identity decomposes a price effect into substitution and income effects.
    \item For normal goods, the law of demand follows immediately because both components are nonpositive.
\end{itemize}
\end{examtip}
\newpage

% chapters/ch03_hicksian_substitution_mrs_and_endowments.tex

\section{Hicksian Substitution, MRS, and Endowments}

Chapter 2 decomposed a Marshallian price effect using Slutsky compensation: adjust income just enough so that the original bundle remains affordable at the new prices. Lecture 3 introduces a second compensation concept, Hicksian compensation, which instead keeps utility constant. The chapter then develops marginal utility and the marginal rate of substitution, and finally studies consumer choice when expenditure comes from an endowment rather than exogenous income.

\subsection{Hicksian Substitution}

Let the initial price vector be $p=(p_1,p_2)$, expenditure be $m$, and initial demand be
\[
x^0 := x(p,m).
\]
Suppose the price of good 1 rises to
\[
p'=(p_1',p_2), \qquad p_1' > p_1.
\]

\begin{definition}{Slutsky and Hicksian compensation}{ch3-compensation}
Let $u_0 := U(x^0)$ denote the utility level at the original optimum.

\begin{itemize}
    \item The \emph{Slutsky compensated expenditure} is
    \[
    m^s := p' \cdot x^0.
    \]
    This is the income needed at the new prices to still afford the original bundle.
    \item The \emph{Hicksian compensated expenditure} is
    \[
    m^h := \min \left\{\, p' \cdot y : y \in \mathbb{R}_+^2,\ U(y) \ge u_0 \,\right\}.
    \]
    This is the minimum expenditure needed at the new prices to reach the original utility level.
\end{itemize}
\end{definition}

The difference is conceptual:
\begin{itemize}
    \item Slutsky compensation keeps the \emph{old bundle} feasible.
    \item Hicksian compensation keeps the \emph{old utility level} feasible at minimum cost.
\end{itemize}

\begin{theorem}{Hicksian compensation is weakly below Slutsky compensation}{ch3-ms-mh}
For a price increase from $p$ to $p'$, the compensation levels satisfy
\[
m^h \le m^s.
\]
Equivalently, Slutsky compensation is weakly larger than Hicksian compensation.
\end{theorem}

\begin{proof}
The original bundle $x^0$ delivers utility $u_0$, so it is feasible in the expenditure-minimization problem defining $m^h$. Therefore
\[
m^h \le p' \cdot x^0 = m^s.
\]
\end{proof}

\begin{examtip}
The intuition is simple: the old bundle is one way to attain the old utility level at the new prices, but it need not be the \emph{cheapest} way. Hence Hicksian compensation is weakly smaller.
\end{examtip}

To isolate substitution under constant utility, let $h(p',u_0)$ denote Hicksian demand at the new prices and old utility. In the lecture's notation, this is often written as $x(p',m^h)$.

\begin{definition}{Hicksian substitution and income effects}{ch3-hicks-effects}
For a price increase in good 1, define
\[
\Delta x_1 := x_1(p',m) - x_1(p,m)
\]
as the total Marshallian change in demand for good 1. The Hicksian decomposition is
\[
\Delta x_1 = \Delta x_1^h + \Delta x_1^q,
\]
where
\[
\Delta x_1^h := h_1(p',u_0) - x_1(p,m)
\]
is the \emph{Hicksian substitution effect}, and
\[
\Delta x_1^q := x_1(p',m) - h_1(p',u_0)
\]
is the \emph{Hicksian income effect}.

\end{definition}

\begin{theorem}{Hicksian compensated demand for the good whose price rises is nonincreasing}{ch3-hicks-negative}
If the price of good 1 rises from $p_1$ to $p_1' > p_1$, then
\[
\Delta x_1^h = h_1(p',u_0) - x_1(p,m) \le 0.
\]
\end{theorem}

\begin{proof}
Let
\[
x^h := h(p',u_0).
\]
Because $x^0$ and $x^h$ both deliver utility level $u_0$, each solves an expenditure-minimization problem at its own price vector:
\[
p \cdot x^0 \le p \cdot x^h,
\qquad
p' \cdot x^h \le p' \cdot x^0.
\]
Subtracting the first inequality from the second gives
\[
(p'-p)\cdot(x^h-x^0)\le 0.
\]
Only the price of good 1 changes, so
\[
(p_1'-p_1)\bigl(x_1^h-x_1^0\bigr)\le 0.
\]
Since $p_1'-p_1>0$, it follows that
\[
x_1^h-x_1^0\le 0.
\]
Hence
\[
\Delta x_1^h\le 0.
\]
\end{proof}

\begin{examplebox}[title={Worked Example: Slutsky versus Hicksian compensation}]
Suppose the original choice is $x^0=x(p,m)$ and utility is $u_0=U(x^0)$. Let new prices be $p'=(p_1',p_2)$ with $p_1'>p_1$.

The two compensation levels are:
\[
m^s = p' \cdot x^0,
\qquad
m^h = \min\{\,p' \cdot y : U(y)\ge u_0\,\}.
\]
Since $U(x^0)=u_0$, the original bundle $x^0$ is feasible for the Hicksian minimization problem. Therefore
\[
m^h \le p' \cdot x^0 = m^s.
\]

So the correct ranking is
\[
\boxed{m^h \le m^s.}
\]

Interpretation:
\begin{itemize}
    \item Slutsky gives enough money to buy the old bundle.
    \item Hicks only gives enough money to reach the old utility level as cheaply as possible.
\end{itemize}
Hence Hicksian compensation is more economical.
\end{examplebox}

\begin{commonmistake}
Do not reverse the inequality. After a price increase, Slutsky compensation is \emph{not} smaller than Hicksian compensation. Keeping the exact old bundle affordable is generally more expensive than keeping only the old utility level attainable.
\end{commonmistake}

\subsection{Marginal Utility and the Marginal Rate of Substitution}

Lecture 3 then formalizes local trade-offs using derivatives.

\begin{definition}{Derivative and partial derivatives}{ch3-derivatives}
If $f:\mathbb{R}\to\mathbb{R}$, the derivative of $f$ at $x$ is
\[
f'(x)=\lim_{\alpha\to 0}\frac{f(x+\alpha)-f(x)}{\alpha},
\]
when the limit exists.

If $f:\mathbb{R}_+^2\to\mathbb{R}$, the partial derivative with respect to the first input is
\[
\frac{\partial f(x_1,x_2)}{\partial x_1}
=
\lim_{\alpha\to 0}
\frac{f(x_1+\alpha,x_2)-f(x_1,x_2)}{\alpha},
\]
and similarly for $\partial f/\partial x_2$.
\end{definition}

In consumer theory, the function of interest is utility.

\begin{definition}{Marginal utility}{ch3-mu}
For a differentiable utility function $U(x_1,\dots,x_L)$, the marginal utility of good $k$ is
\[
MU_k := \frac{\partial U(x_1,\dots,x_L)}{\partial x_k}.
\]
It measures the rate at which utility changes when the quantity of good $k$ increases by a small amount, holding other goods fixed.
\end{definition}

\begin{definition}{Marginal rate of substitution}{ch3-mrs}
In the two-good case, fix a utility level $u$ and write the indifference curve locally as
\[
x_2 = f(x_1,u)
\quad \text{so that} \quad
U(x_1,f(x_1,u))=u.
\]
The \emph{marginal rate of substitution of good 1 for good 2} is
\[
MRS_{12} := -\frac{\partial f(x_1,u)}{\partial x_1}.
\]
It is the amount of good 2 the consumer is willing to give up for a small increase in good 1 while holding utility constant.
\end{definition}

\begin{theorem}{MRS equals the ratio of marginal utilities}{ch3-mrs-mu}
Suppose $U$ is differentiable and $MU_2>0$. Then along an indifference curve,
\[
MRS_{12}=\frac{MU_1}{MU_2}.
\]
\end{theorem}

\begin{proof}
Since utility is constant along the indifference curve,
\[
U(x_1,f(x_1,u))=u.
\]
Differentiate both sides with respect to $x_1$:
\[
\frac{\partial U}{\partial x_1}
+
\frac{\partial U}{\partial x_2}\cdot \frac{\partial f(x_1,u)}{\partial x_1}
=0.
\]
Rearranging gives
\[
-\frac{\partial f(x_1,u)}{\partial x_1}
=
\frac{\frac{\partial U}{\partial x_1}}{\frac{\partial U}{\partial x_2}}
=
\frac{MU_1}{MU_2}.
\]
By definition of $MRS_{12}$, this proves the result.
\end{proof}

\begin{theorem}{Interior tangency condition}{ch3-tangency}
If utility is differentiable and the consumer chooses an interior optimum from the budget set
\[
p_1x_1+p_2x_2=m,
\]
then
\[
MRS_{12}=\frac{MU_1}{MU_2}=\frac{p_1}{p_2}.
\]
\end{theorem}

\begin{proof}
At an interior optimum, the indifference curve is tangent to the budget line. The slope of the budget line is $-p_1/p_2$, while the slope of the indifference curve is $-MRS_{12}$. Equality of slopes gives
\[
MRS_{12}=\frac{p_1}{p_2}.
\]
Combining with Theorem \ref{thm:ch3-mrs-mu} yields
\[
\frac{MU_1}{MU_2}=\frac{p_1}{p_2}.
\]
\end{proof}

\begin{examplebox}[title={Worked Example: Cobb--Douglas utility}]
Let
\[
U(x_1,x_2)=x_1x_2.
\]
Then
\[
MU_1=\frac{\partial U}{\partial x_1}=x_2,
\qquad
MU_2=\frac{\partial U}{\partial x_2}=x_1.
\]
Hence
\[
MRS_{12}=\frac{MU_1}{MU_2}=\frac{x_2}{x_1}.
\]

At an interior optimum,
\[
\frac{x_2}{x_1}=\frac{p_1}{p_2}
\quad \Longrightarrow \quad
x_2=\frac{p_1}{p_2}x_1.
\]
Substitute into the budget constraint:
\[
p_1x_1+p_2x_2
=
p_1x_1+p_2\left(\frac{p_1}{p_2}x_1\right)
=
2p_1x_1
=
m.
\]
So
\[
x_1=\frac{m}{2p_1},
\qquad
x_2=\frac{m}{2p_2}.
\]
Therefore the Marshallian demand is
\[
\boxed{
x(p_1,p_2,m)
=
\left(
\frac{m}{2p_1},
\frac{m}{2p_2}
\right).
}
\]

If the endowment is $\omega=(3,1)$ and prices are $(p_1,p_2)=(1,3)$, then
\[
m=p\cdot \omega = 1\cdot 3 + 3\cdot 1 = 6,
\]
so
\[
x(1,3,6)=\left(\frac{6}{2},\frac{6}{6}\right)=(3,1).
\]
Thus the consumer chooses exactly the endowment.
\end{examplebox}

\begin{examplebox}[title={Worked Example: quasilinear utility and corner logic}]
Let
\[
U(x_1,x_2)=x_1+2\sqrt{x_2}.
\]
Then
\[
MU_1=1,
\qquad
MU_2=\frac{1}{\sqrt{x_2}},
\qquad
MRS_{12}=\frac{MU_1}{MU_2}=\sqrt{x_2}.
\]

Consider prices $(p_1,p_2)=(2,1)$ and expenditure $m=1$. An interior tangency would require
\[
\sqrt{x_2}=\frac{p_1}{p_2}=2
\quad \Longrightarrow \quad
x_2=4,
\]
but this is impossible because the budget set satisfies
\[
2x_1+x_2\le 1.
\]
So the tangency candidate is infeasible, and the optimum must be a corner.

Because utility is increasing in both goods, the budget binds:
\[
2x_1+x_2=1
\quad \Longrightarrow \quad
x_1=\frac{1-x_2}{2}.
\]
Hence maximize
\[
f(x_2)=\frac{1-x_2}{2}+2\sqrt{x_2},
\qquad x_2\in[0,1].
\]
Then
\[
f'(x_2)=-\frac12+\frac{1}{\sqrt{x_2}}>0
\quad \text{for } x_2\in(0,1].
\]
So $f$ is increasing on $[0,1]$, and the optimum is
\[
x_2^*=1,\qquad x_1^*=0.
\]

Thus
\[
\boxed{x(2,1,1)=(0,1).}
\]

This is a standard exam pattern: the tangency condition gives a candidate, but feasibility and corner checks still matter.
\end{examplebox}

\begin{commonmistake}
The tangency condition
\[
\frac{MU_1}{MU_2}=\frac{p_1}{p_2}
\]
applies only at an \emph{interior} optimum. If the tangency candidate is infeasible or gives a negative quantity, the true optimum is a corner and must be checked separately.
\end{commonmistake}

\subsection{Endowments and Net Demand}

Up to this point, expenditure $m$ was taken as given. Lecture 3 now studies the case where expenditure comes from selling an initial endowment.

\begin{definition}{Endowment economy budget set}{ch3-endowment}
Let the consumer's endowment be
\[
\omega=(\omega_1,\dots,\omega_L)\in \mathbb{R}_+^L.
\]
At prices $p$, the value of the endowment is
\[
m=p\cdot \omega.
\]
The consumer solves
\[
\max_{x\in \mathbb{R}_+^L} U(x)
\quad \text{subject to} \quad
p\cdot x \le p\cdot \omega.
\]
If preferences are increasing, the budget binds:
\[
p\cdot x(p,p\cdot \omega)=p\cdot \omega.
\]
\end{definition}

\begin{definition}{Net demand}{ch3-net-demand}
Given endowment $\omega$ and Marshallian demand $x(p,m)$, the \emph{net demand function} is
\[
z(p):=x(p,p\cdot \omega)-\omega.
\]
For each good $\ell$,
\[
z_\ell(p)=x_\ell(p,p\cdot \omega)-\omega_\ell.
\]
Interpretation:
\begin{itemize}
    \item $z_\ell(p)>0$: the consumer is a net buyer of good $\ell$;
    \item $z_\ell(p)<0$: the consumer is a net seller of good $\ell$.
\end{itemize}
\end{definition}

\begin{theorem}{Walras' law for net demand}{ch3-walras}
For any endowment $\omega$ and increasing utility,
\[
p\cdot z(p)=0.
\]
\end{theorem}

\begin{proof}
By definition,
\[
z(p)=x(p,p\cdot \omega)-\omega.
\]
Taking the dot product with $p$ gives
\[
p\cdot z(p)
=
p\cdot x(p,p\cdot \omega)-p\cdot \omega.
\]
Because the chosen bundle exhausts the budget,
\[
p\cdot x(p,p\cdot \omega)=p\cdot \omega.
\]
Hence
\[
p\cdot z(p)=0.
\]
\end{proof}

In the two-good case, positive prices imply the two net demands must have opposite signs whenever one is nonzero. Thus a consumer cannot be a net buyer of both goods simultaneously.

\begin{examplebox}[title={Worked Example: net demand and Walras' law}]
Suppose
\[
U(x_1,x_2)=\sqrt{x_1}+\sqrt{x_2},
\]
which generates the Marshallian demand
\[
x(p_1,p_2,m)
=
\frac{m}{p_1+p_2}
\left(
\frac{p_2}{p_1},
\frac{p_1}{p_2}
\right).
\]
Let the endowment be $\omega=(3,1)$.

At prices $p=(1,1)$, endowment income is
\[
m=p\cdot \omega = 1\cdot 3 + 1\cdot 1 = 4.
\]
So
\[
x(1,1,4)=\frac{4}{2}(1,1)=(2,2).
\]
Hence net demand is
\[
z(1,1)=(2,2)-(3,1)=(-1,1).
\]
Therefore:
\begin{itemize}
    \item the consumer is a net seller of good 1;
    \item the consumer is a net buyer of good 2.
\end{itemize}

Walras' law holds:
\[
(1,1)\cdot(-1,1)=0.
\]
So
\[
\boxed{z(1,1)=(-1,1),\qquad p\cdot z(p)=0.}
\]
\end{examplebox}

\subsection{Price Changes with Endowments}

Price changes have a new feature when wealth depends on prices: the consumer's budget changes not only because the slope of the budget line changes, but also because the value of the endowment changes.

Suppose the price of good 1 rises from $p_1$ to $p_1' > p_1$, with $p_2$ unchanged. Let
\[
\tilde x := x(p,p\cdot \omega)
\]
denote the original chosen bundle, and define the Slutsky compensation relative to that bundle as
\[
m^s := p'\cdot \tilde x.
\]

\begin{definition}{Slutsky decomposition with endowments}{ch3-endowment-slutsky}
Define
\[
\Delta x_1 := x_1(p',p'\cdot \omega)-x_1(p,p\cdot \omega),
\]
\[
\Delta x_1^s := x_1(p',m^s)-x_1(p,p\cdot \omega),
\]
\[
\Delta x_1^n := x_1(p',p'\cdot \omega)-x_1(p',m^s).
\]
Then
\[
\Delta x_1 = \Delta x_1^s + \Delta x_1^n.
\]
\end{definition}

As in Chapter 2, the substitution effect $\Delta x_1^s$ is always nonpositive when the price of good 1 rises. The new issue is the sign of the income effect $\Delta x_1^n$.

\begin{theorem}{Endowment-value identity}{ch3-income-identity}
Let $z_1(p)=x_1(p,p\cdot \omega)-\omega_1$ denote net demand for good 1 at the original prices. Then
\[
p'\cdot \omega = m^s - (p_1'-p_1)z_1(p).
\]
\end{theorem}

\begin{proof}
Let $\tilde x = x(p,p\cdot \omega)$. Since the original bundle exhausts the original budget,
\[
p\cdot \tilde x = p\cdot \omega.
\]
Now
\begin{align*}
p'\cdot \omega
&= p'\cdot \omega + (p\cdot \tilde x - p\cdot \omega) + (p'\cdot \tilde x - p'\cdot \tilde x) \\
&= p'\cdot (\omega-\tilde x) + p\cdot (\tilde x-\omega) + p'\cdot \tilde x \\
&= p'\cdot \tilde x - (p'-p)\cdot(\tilde x-\omega).
\end{align*}
Only the price of good 1 changes, so
\[
(p'-p)\cdot(\tilde x-\omega)=(p_1'-p_1)(\tilde x_1-\omega_1)=(p_1'-p_1)z_1(p).
\]
Since $m^s=p'\cdot \tilde x$, the result follows:
\[
p'\cdot \omega = m^s-(p_1'-p_1)z_1(p).
\]
\end{proof}

This identity immediately explains the sign of the income effect:

\begin{itemize}
    \item If $z_1(p)>0$ (the consumer is a net buyer of good 1), then
    \[
    p'\cdot \omega < m^s.
    \]
    So the consumer is poorer than under Slutsky compensation.
    \item If $z_1(p)<0$ (the consumer is a net seller of good 1), then
    \[
    p'\cdot \omega > m^s.
    \]
    So the consumer is richer than under Slutsky compensation.
\end{itemize}

\begin{theorem}{Welfare and direction of the budget shift}{ch3-buyer-seller}
Suppose the price of good 1 rises.

\begin{enumerate}
    \item If the consumer is a net buyer of good 1 at the original prices, then welfare falls.
    \item If the consumer is a net seller of good 1 at the original prices, then welfare rises.
\end{enumerate}
\end{theorem}

\begin{proof}
Let $\tilde x=x(p,p\cdot \omega)$.

If the consumer is a net buyer of good 1, then $\tilde x_1>\omega_1$. Hence
\[
p'\cdot \tilde x - p'\cdot \omega = (p_1'-p_1)(\tilde x_1-\omega_1)>0.
\]
So the old bundle $\tilde x$ is no longer affordable at the new prices and new wealth. Therefore the consumer cannot attain the old optimum, so welfare falls.

If the consumer is a net seller of good 1, then $\tilde x_1<\omega_1$. Hence
\[
p'\cdot \tilde x - p'\cdot \omega = (p_1'-p_1)(\tilde x_1-\omega_1)<0.
\]
So the old bundle is affordable at the new prices and new wealth, with money left over. Since preferences are increasing, the consumer can buy something strictly better. Therefore welfare rises.
\end{proof}

\begin{theorem}{A seller of the good whose price rises cannot become a buyer}{ch3-cannot-become-buyer}
Suppose the consumer is a net seller of good 1 at the original prices:
\[
x_1(p,p\cdot \omega)<\omega_1.
\]
If the price of good 1 rises, then at the new prices the consumer cannot choose a bundle with
\[
x_1(p',p'\cdot \omega)>\omega_1.
\]
\end{theorem}

\begin{proof}
Let $\tilde x=x(p,p\cdot \omega)$ and suppose for contradiction that the new choice $x'$ satisfies $x_1'>\omega_1$.

Since $\tilde x_1<\omega_1$, at the new prices
\[
p'\cdot \tilde x < p'\cdot \omega = p'\cdot x'.
\]
So $x'$ is chosen when $\tilde x$ is strictly cheaper, hence
\[
x' P \tilde x.
\]

On the other hand, because $x_1'>\omega_1$,
\[
p\cdot x' = p'\cdot x' - (p_1'-p_1)x_1'
< p'\cdot \omega - (p_1'-p_1)\omega_1
= p\cdot \omega
= p\cdot \tilde x.
\]
Thus under the original prices, $x'$ was strictly cheaper than the old chosen bundle, so
\[
\tilde x P x'.
\]
This violates WARP. Therefore the new choice cannot satisfy $x_1'>\omega_1$.
\end{proof}

\begin{theorem}{Law of demand with endowments for a normal good and a net buyer}{ch3-law-demand-endowment}
Suppose good 1 is normal, and the consumer is a net buyer of good 1 at the original prices:
\[
z_1(p)\ge 0.
\]
If the price of good 1 rises from $p_1$ to $p_1' > p_1$, then
\[
x_1(p',p'\cdot \omega)\le x_1(p,p\cdot \omega).
\]
\end{theorem}

\begin{proof}
By Definition \ref{def:ch3-endowment-slutsky},
\[
\Delta x_1=\Delta x_1^s+\Delta x_1^n.
\]
The substitution effect satisfies
\[
\Delta x_1^s\le 0.
\]
By Theorem \ref{thm:ch3-income-identity},
\[
p'\cdot \omega = m^s-(p_1'-p_1)z_1(p)\le m^s
\]
because $z_1(p)\ge 0$. Since good 1 is normal, a lower income at fixed prices $p'$ weakly lowers its demand. Therefore
\[
\Delta x_1^n = x_1(p',p'\cdot \omega)-x_1(p',m^s)\le 0.
\]
So
\[
\Delta x_1=\Delta x_1^s+\Delta x_1^n\le 0.
\]
This is exactly
\[
x_1(p',p'\cdot \omega)\le x_1(p,p\cdot \omega).
\]
\end{proof}

\begin{examplebox}[title={Worked Example: seller of good 1 and the sign of the income effect}]
Let
\[
U(x_1,x_2)=\sqrt{x_1}+\sqrt{x_2},
\qquad
x(p_1,p_2,m)=\frac{m}{p_1+p_2}\left(\frac{p_2}{p_1},\frac{p_1}{p_2}\right),
\]
and let the endowment be
\[
\omega=(3,1).
\]

\textbf{Step 1: original prices.}  
At
\[
p=(1,1),
\]
income is
\[
m_0=(1,1)\cdot(3,1)=4.
\]
Hence
\[
x^0=x(1,1,4)=\frac{4}{2}(1,1)=(2,2).
\]
So
\[
z(1,1)=x^0-\omega=(2,2)-(3,1)=(-1,1).
\]
The consumer is a net seller of good 1.

\textbf{Step 2: new prices.}  
Let
\[
p'=(2,1).
\]
New endowment income is
\[
m_1=(2,1)\cdot(3,1)=7.
\]
So
\[
x^1=x(2,1,7)=\frac{7}{3}\left(\frac12,2\right)=\left(\frac76,\frac{14}{3}\right).
\]

\textbf{Step 3: Slutsky compensation.}  
The compensated income needed to still afford the original bundle is
\[
m^s=p'\cdot x^0=(2,1)\cdot(2,2)=6.
\]
So the compensated demand is
\[
x^c=x(2,1,6)=\frac{6}{3}\left(\frac12,2\right)=(1,4).
\]

\textbf{Step 4: decomposition for good 1.}
\[
\Delta x_1^s=x_1^c-x_1^0=1-2=-1,
\]
\[
\Delta x_1^n=x_1^1-x_1^c=\frac76-1=\frac16.
\]
Thus
\[
\Delta x_1=\Delta x_1^s+\Delta x_1^n=-1+\frac16=-\frac56.
\]

So
\[
\boxed{
\Delta x_1^s=-1,
\qquad
\Delta x_1^n=\frac16.
}
\]

Interpretation:
\begin{itemize}
    \item the substitution effect is negative, as usual;
    \item because the consumer is a seller of good 1, the price increase raises wealth, so the income effect is positive;
    \item the total change is still negative here because substitution dominates.
\end{itemize}
\end{examplebox}

\begin{policynote}[title={Policy Application: wages and labour supply}]
Lecture 3 interprets an endowment model as a labour-supply problem. Let good 1 be \emph{time/leisure} and let the endowment be
\[
\omega=(24,0),
\]
meaning the consumer begins each day with 24 hours of time and no composite consumption good. If the price vector is
\[
p=(w,1),
\]
then $w$ is the wage. A worker who supplies labour is a net seller of leisure/time, so a wage increase raises the value of the endowment.

Two forces then conflict:
\begin{itemize}
    \item the substitution effect makes leisure more expensive, pushing the consumer toward less leisure and more labour;
    \item the income effect from higher wages makes the consumer richer, which can increase demand for leisure if leisure is normal.
\end{itemize}

Hence a higher wage can increase or decrease labour supplied. This is the standard intuition behind backward-bending labour supply.
\end{policynote}

\subsection{Chapter Summary}

\begin{keyformula}[title={Key Formula: Lecture 3 Summary}]
\begin{align*}
m^s &= p' \cdot x(p,m), \\
m^h &= \min\{\, p' \cdot y : U(y)\ge U(x(p,m)) \,\}, \\
m^h &\le m^s, \\
\Delta x_1 &= \Delta x_1^h + \Delta x_1^q, \\
MRS_{12} &= -\frac{\partial f(x_1,u)}{\partial x_1}
= \frac{MU_1}{MU_2}, \\
\text{interior optimum: } \quad \frac{MU_1}{MU_2} &= \frac{p_1}{p_2}, \\
z(p) &= x(p,p\cdot \omega)-\omega, \\
p\cdot z(p) &= 0, \\
p'\cdot \omega &= m^s-(p_1'-p_1)z_1(p).
\end{align*}
If good 1 is normal and the consumer is a net buyer of good 1, then a rise in $p_1$ implies
\[
x_1(p',p'\cdot \omega)\le x_1(p,p\cdot \omega).
\]
\end{keyformula}

\begin{examtip}
The highest-yield points from Lecture 3 are:
\begin{itemize}
    \item Slutsky and Hicksian compensation are different; Hicksian compensation is weakly smaller.
    \item Hicksian substitution always reduces compensated demand for the good whose price rises.
    \item Know how to compute $MU_1$, $MU_2$, $MRS_{12}$, and the tangency condition.
    \item In endowment models, income changes with prices because wealth equals $p\cdot \omega$.
    \item Use net demand
    \[
    z(p)=x(p,p\cdot \omega)-\omega
    \]
    to determine whether the consumer is a buyer or seller.
    \item For a net buyer, a price increase hurts welfare; for a net seller, it helps.
    \item Labour supply is ambiguous because the substitution and income effects work in opposite directions.
\end{itemize}
\end{examtip}
\newpage
% chapters/ch04_intertemporal_consumption_and_asset_pricing.tex

\section{Intertemporal Consumption and Asset Pricing}

This chapter studies choice across time. The first part rewrites intertemporal consumption as an ordinary consumer-choice problem with dated consumption goods and relative prices determined by the interest rate. The second part introduces discounted utility and the Euler-style optimality condition for present versus future consumption. The final part applies present-value logic to asset pricing, net present value, perpetuities, installment loans, and asset returns under taxes.

\subsection{Intertemporal Consumption: Two-Period Model}

Consider a consumer who lives for two periods, $t=0,1$. Let
\[
x=(x_0,x_1)\in \mathbb{R}_+^2
\]
denote the consumption stream, where $x_t$ is expenditure on consumption in period $t$. Let
\[
y=(y_0,y_1)\in \mathbb{R}_+^2
\]
denote the income stream. The consumer can save or borrow between the two periods at interest rate $r\ge 0$.

\begin{definition}{Two-period savings formulation}{ch4-two-period-savings}
Let $s_0 \in \mathbb{R}$ denote period-0 saving, where
\[
s_0>0 \text{ means saving}, \qquad s_0<0 \text{ means borrowing}.
\]
The two-period budget equations are
\[
x_0+s_0=y_0,
\qquad
x_1=(1+r)s_0+y_1.
\]
\end{definition}

The first equation says current income is split between current consumption and saving. The second says future consumption equals future income plus the gross return on current saving.

\begin{theorem}{Present-value and future-value budget constraints}{ch4-pv-fv}
The two-period savings system
\[
x_0+s_0=y_0,
\qquad
x_1=(1+r)s_0+y_1
\]
is equivalent to either of the following:
\[
(1+r)x_0+x_1=(1+r)y_0+y_1,
\]
or
\[
x_0+\frac{x_1}{1+r}=y_0+\frac{y_1}{1+r}.
\]
\end{theorem}

\begin{proof}
From
\[
x_0+s_0=y_0,
\]
multiply by $1+r$:
\[
(1+r)x_0+(1+r)s_0=(1+r)y_0.
\]
Add the second budget equation:
\[
x_1=(1+r)s_0+y_1.
\]
Cancelling $(1+r)s_0$ gives
\[
(1+r)x_0+x_1=(1+r)y_0+y_1.
\]
Dividing through by $1+r$ yields
\[
x_0+\frac{x_1}{1+r}=y_0+\frac{y_1}{1+r}.
\]
\end{proof}

\begin{definition}{Present value and future value}{ch4-pv-fv-def}
For a two-period stream $(a_0,a_1)$:
\begin{itemize}
    \item its \emph{future value} at date 1 is
    \[
    FV(a_0,a_1)=(1+r)a_0+a_1;
    \]
    \item its \emph{present value} at date 0 is
    \[
    PV(a_0,a_1)=a_0+\frac{a_1}{1+r}.
    \]
\end{itemize}
Thus an intertemporal consumption plan is affordable if and only if
\[
PV(x_0,x_1)=PV(y_0,y_1).
\]
\end{definition}

This turns intertemporal choice into an ordinary consumer problem with two goods:
\[
p=(p_0,p_1)=\left(1,\frac{1}{1+r}\right),
\]
so that
\[
p\cdot x = p\cdot y.
\]
The ``price'' of future consumption is therefore
\[
p_1=\frac{1}{1+r}.
\]
A rise in $r$ lowers the price of future consumption relative to present consumption.

\begin{keyformula}[title={Key Formula: Two-Period Intertemporal Budget}]
\[
x_0+\frac{x_1}{1+r}=y_0+\frac{y_1}{1+r}.
\]
Equivalently, with
\[
p=\left(1,\frac{1}{1+r}\right),
\]
the budget set is
\[
B(p,p\cdot y)=\{x\in \mathbb{R}_+^2 : p\cdot x \le p\cdot y\}.
\]
\end{keyformula}

\begin{examplebox}[title={Worked Example: present value, future value, and saving}]
Suppose
\[
y=(190,11), \qquad r=0.1,
\]
and utility is
\[
U(x_0,x_1)=x_0x_1.
\]
The present value of income is
\[
PV(y)=190+\frac{11}{1.1}=190+10=200.
\]
The future value of income is
\[
FV(y)=1.1(190)+11=209+11=220.
\]

For Cobb--Douglas utility $U(x_0,x_1)=x_0x_1$, the optimal demand with prices
\[
\left(1,\frac{1}{1+r}\right)
\]
and present-value wealth $m$ is
\[
x(r,m)=\left(\frac{m}{2},\frac{(1+r)m}{2}\right).
\]
Hence, with $m=200$ and $r=0.1$,
\[
x(0.1,200)=\left(100,\frac{1.1\cdot 200}{2}\right)=(100,110).
\]

So the consumer saves
\[
s_0=y_0-x_0=190-100=90
\]
in period 0, and this becomes
\[
(1+r)s_0=1.1\cdot 90=99
\]
in period 1. Since future income is $11$, future consumption is
\[
99+11=110.
\]

Therefore the optimal plan is
\[
\boxed{x=(100,110), \qquad s_0=90.}
\]
\end{examplebox}

\subsection{Comparative Statics with Respect to the Interest Rate}

Because
\[
p_1=\frac{1}{1+r},
\]
a rise in the interest rate lowers the relative price of future consumption. Thus intertemporal comparative statics are a direct application of the consumer theory from Chapters 2 and 3.

\begin{definition}{Saver, borrower, and net demand for future consumption}{ch4-saver-borrower}
Given income stream $y=(y_0,y_1)$ and chosen consumption stream $x=(x_0,x_1)$:
\begin{itemize}
    \item the consumer is a \emph{saver} if
    \[
    x_1>y_1
    \quad \Longleftrightarrow \quad
    s_0>0;
    \]
    \item the consumer is a \emph{borrower} if
    \[
    x_1<y_1
    \quad \Longleftrightarrow \quad
    s_0<0.
    \]
\end{itemize}
Thus a saver is a net buyer of future consumption, while a borrower is a net seller of future consumption.
\end{definition}

This matches the endowment logic from Chapter 3 with dated consumption as the two goods and income stream $y$ as the endowment.

\begin{theorem}{Effect of an increase in the interest rate on saver/borrower status}{ch4-saver-status}
Suppose the interest rate rises from $r$ to $r'>r$.
\begin{enumerate}
    \item A saver stays a saver.
    \item A borrower may remain a borrower or may become a saver.
\end{enumerate}
\end{theorem}

\begin{proof}
Future consumption has price
\[
q(r)=\frac{1}{1+r}.
\]
If $r'>r$, then
\[
q(r')<q(r),
\]
so the price of future consumption falls. By Chapter 3, when the price of a good falls, a net buyer of that good remains a net buyer. Thus a saver stays a saver. A net seller may remain a seller or may switch to being a buyer.
\end{proof}

To apply the Slutsky decomposition, write Marshallian intertemporal demand as
\[
x(r,m)=\bigl(x_0(r,m),x_1(r,m)\bigr),
\]
where $m$ is present-value wealth. If the interest rate changes from $r$ to $r'$, then the present value of the income stream changes from
\[
m=y_0+\frac{y_1}{1+r}
\]
to
\[
m'=y_0+\frac{y_1}{1+r'}.
\]

\begin{definition}{Slutsky decomposition for intertemporal choice}{ch4-slutsky-r}
Let
\[
\Delta x_1 := x_1(r',m')-x_1(r,m)
\]
be the total change in future consumption. Define the Slutsky compensated present-value wealth at the new interest rate by
\[
m^s := x_0(r,m)+\frac{x_1(r,m)}{1+r'}.
\]
Then
\[
\Delta x_1 = \Delta x_1^s + \Delta x_1^n,
\]
where
\[
\Delta x_1^s := x_1(r',m^s)-x_1(r,m)
\]
is the substitution effect, and
\[
\Delta x_1^n := x_1(r',m')-x_1(r',m^s)
\]
is the income effect.
\end{definition}

Since $r'>r$ implies the price of future consumption falls, the substitution effect for future consumption is nonnegative:
\[
\Delta x_1^s \ge 0.
\]

If future consumption is normal and the consumer is a saver, then the income effect is also nonnegative. The reason is that a saver benefits when the interest rate rises: the price of the good she is a net buyer of falls.

\begin{theorem}{Interest-rate comparative statics for savers}{ch4-saver-theorem}
Suppose future consumption is normal and the consumer is a saver at the initial interest rate. If $r$ increases, then
\[
x_1(r',m') \ge x_1(r,m).
\]
\end{theorem}

\begin{proof}
When $r$ rises, the price of future consumption falls, so the substitution effect satisfies
\[
\Delta x_1^s \ge 0.
\]
Because the consumer is a saver, she is a net buyer of future consumption. Thus the fall in its price raises her welfare and, if future consumption is normal, the associated income effect also satisfies
\[
\Delta x_1^n \ge 0.
\]
Hence
\[
\Delta x_1=\Delta x_1^s+\Delta x_1^n\ge 0.
\]
\end{proof}

\begin{examtip}
A rise in $r$ does \emph{not} mean current wealth necessarily rises. The correct logic is:
\begin{enumerate}
    \item the price of future consumption falls;
    \item substitution pushes toward more future consumption;
    \item the welfare effect depends on whether the consumer is a saver or borrower;
    \item for a saver, both substitution and income effects push toward more future consumption if future consumption is normal.
\end{enumerate}
For a borrower, the overall effect on future consumption is ambiguous.
\end{examtip}

\subsection{Multi-Period Intertemporal Choice}

Now suppose the consumer lives for periods
\[
t=0,1,\dots,T.
\]
Let
\[
x=(x_0,x_1,\dots,x_T)\in \mathbb{R}_+^{T+1},
\qquad
y=(y_0,y_1,\dots,y_T)\in \mathbb{R}_+^{T+1}.
\]
Let $s_t$ denote assets carried from period $t$ to period $t+1$, with
\[
s_{-1}=0, \qquad s_T=0.
\]

\begin{definition}{Multi-period intertemporal budget equations}{ch4-multi-budget}
For each period $t=0,1,\dots,T$,
\[
x_t+s_t=(1+r)s_{t-1}+y_t.
\]
\end{definition}

\begin{theorem}{Multi-period present-value budget constraint}{ch4-multi-pv}
The multi-period budget system is equivalent to
\[
\sum_{t=0}^{T}\frac{x_t}{(1+r)^t}
=
\sum_{t=0}^{T}\frac{y_t}{(1+r)^t}.
\]
Equivalently, with price vector
\[
p=\left(1,\frac{1}{1+r},\frac{1}{(1+r)^2},\dots,\frac{1}{(1+r)^T}\right),
\]
the budget constraint is
\[
p\cdot x = p\cdot y.
\]
\end{theorem}

\begin{proof}
Divide each period budget equation
\[
x_t+s_t=(1+r)s_{t-1}+y_t
\]
by $(1+r)^t$:
\[
\frac{x_t}{(1+r)^t}+\frac{s_t}{(1+r)^t}
=
\frac{s_{t-1}}{(1+r)^{t-1}}+\frac{y_t}{(1+r)^t}.
\]
Summing from $t=0$ to $T$ causes the asset terms to telescope:
\[
\sum_{t=0}^{T}\frac{s_t}{(1+r)^t}
=
\sum_{t=0}^{T}\frac{s_{t-1}}{(1+r)^{t-1}}
\]
because $s_{-1}=s_T=0$. Hence
\[
\sum_{t=0}^{T}\frac{x_t}{(1+r)^t}
=
\sum_{t=0}^{T}\frac{y_t}{(1+r)^t}.
\]
\end{proof}

\begin{theorem}{Only present value matters for demand}{ch4-only-pv-matters}
Suppose the consumer solves
\[
\max_{x\in \mathbb{R}_+^{T+1}} U(x)
\quad \text{subject to} \quad
p\cdot x \le p\cdot y.
\]
If two income streams $y$ and $y'$ satisfy
\[
p\cdot y = p\cdot y',
\]
then they generate the same budget set and therefore the same demand.
\end{theorem}

\begin{proof}
The feasible set depends only on the scalar wealth level $p\cdot y$. If
\[
p\cdot y = p\cdot y',
\]
then
\[
\{x : p\cdot x \le p\cdot y\}
=
\{x : p\cdot x \le p\cdot y'\}.
\]
Hence the maximization problem is identical, so the solution is the same.
\end{proof}

\begin{examplebox}[title={Worked Example: different income timings with the same present value}]
Let $T=2$ and $r=\tfrac12$. Then
\[
p=\left(1,\frac{2}{3},\left(\frac{2}{3}\right)^2\right).
\]
Consider the three income streams
\[
y=(2,3,9), \qquad y'=(8,0,0), \qquad y''=(0,0,18).
\]
Their present values are:
\[
p\cdot y = 2+\frac{2}{3}\cdot 3 + \left(\frac{2}{3}\right)^2\cdot 9 = 2+2+4=8,
\]
\[
p\cdot y' = 8,
\]
\[
p\cdot y'' = \left(\frac{2}{3}\right)^2\cdot 18 = \frac{4}{9}\cdot 18 = 8.
\]
Thus all three income streams have the same present value. Therefore they generate the same intertemporal budget set and the same optimal demand.

So the correct conclusion is:
\[
\boxed{
(2,3,9),\ (8,0,0),\ \text{and } (0,0,18)
\text{ are equivalent for choice purposes at } r=\tfrac12.
}
\]
\end{examplebox}

\begin{commonmistake}
Do not say that the \emph{timing} of income never matters. It matters through its effect on present value. What Lecture 4 shows is that if the present value is the same, then the timing alone does not change the budget set under perfect borrowing and lending at rate $r$.
\end{commonmistake}

\subsection{Inflation and the Real Interest Rate}

Suppose the monetary price of one unit of consumption rises by factor $1+\pi$ each period. Then one unit of consumption costs:
\[
1 \text{ in period 0}, \quad 1+\pi \text{ in period 1}, \quad (1+\pi)^2 \text{ in period 2}, \dots
\]

\begin{definition}{Intertemporal budget with inflation}{ch4-inflation-budget}
With inflation rate $\pi$, the period-$t$ budget equation becomes
\[
(1+\pi)^t x_t + s_t = (1+r)s_{t-1} + (1+\pi)^t y_t.
\]
\end{definition}

Dividing by $(1+r)^t$ and summing over time yields the \emph{real} present-value budget:

\begin{theorem}{Real present-value budget}{ch4-real-pv}
With inflation rate $\pi$ and nominal interest rate $r$,
\[
\sum_{t=0}^{T}\left(\frac{1+\pi}{1+r}\right)^t x_t
=
\sum_{t=0}^{T}\left(\frac{1+\pi}{1+r}\right)^t y_t.
\]
Hence the real price vector is
\[
p=\left(1,\frac{1+\pi}{1+r},\left(\frac{1+\pi}{1+r}\right)^2,\dots,\left(\frac{1+\pi}{1+r}\right)^T\right).
\]
\end{theorem}

\begin{definition}{Real interest rate}{ch4-real-interest}
The \emph{real interest rate} $\rho$ is defined by
\[
1+\rho = \frac{1+r}{1+\pi}.
\]
Equivalently,
\[
\rho = \frac{1+r}{1+\pi}-1 = \frac{r-\pi}{1+\pi} \approx r-\pi.
\]
\end{definition}

The approximation
\[
\rho \approx r-\pi
\]
is accurate when inflation is moderate. It is the standard ``nominal minus inflation'' rule.

\begin{examplebox}[title={Worked Example: computing the real interest rate}]
Suppose
\[
r=0.32, \qquad \pi=0.20.
\]
Then
\[
1+\rho = \frac{1.32}{1.20}=1.1,
\]
so
\[
\rho = 0.1.
\]
If the income stream is
\[
y=(100,132),
\]
then its real present value is
\[
100+\frac{132}{1.1}=100+120=220.
\]
Hence
\[
\boxed{\rho=10\%, \qquad PV_{\text{real}}(y)=220.}
\]
\end{examplebox}

\subsection{Discounted Utility and Optimal Consumption Smoothing}

Lecture 4 uses the standard discounted-utility specification
\[
U(x_0,x_1)=u(x_0)+\delta u(x_1),
\]
where $u$ is the felicity function and $\delta \in [0,1]$ is the discount factor.

\begin{definition}{Discounted utility}{ch4-discounted-utility}
A two-period discounted-utility representation is
\[
U(x_0,x_1)=u(x_0)+\delta u(x_1),
\]
where:
\begin{itemize}
    \item $u:\mathbb{R}_+\to \mathbb{R}$ is increasing;
    \item $u$ is concave, so marginal utility is diminishing;
    \item $\delta \in [0,1]$ measures the value placed on future felicity.
\end{itemize}
\end{definition}

\begin{theorem}{Concavity implies diminishing marginal utility}{ch4-dmu}
If $u$ is concave and differentiable, then $u'$ is decreasing.
\end{theorem}

\begin{proof}
Concavity means that for any $x>x'$ and any $\alpha>0$,
\[
u(x+\alpha)-u(x) \le u(x'+\alpha)-u(x').
\]
Divide by $\alpha$ and let $\alpha \to 0$ to obtain
\[
u'(x) \le u'(x').
\]
Hence $u'$ is decreasing.
\end{proof}

\begin{theorem}{Euler-type first-order condition for discounted utility}{ch4-euler}
Suppose the consumer maximizes
\[
u(x_0)+\delta u(x_1)
\]
subject to
\[
x_0+\frac{x_1}{1+r}=m,
\]
and the optimum is interior. Then
\[
\frac{u'(x_0)}{\delta u'(x_1)} = 1+r,
\]
equivalently,
\[
u'(x_0) = (1+r)\delta u'(x_1).
\]
\end{theorem}

\begin{proof}
The marginal utilities are
\[
MU_0=u'(x_0), \qquad MU_1=\delta u'(x_1).
\]
By the interior tangency condition from Chapter 3,
\[
\frac{MU_0}{MU_1} = \frac{p_0}{p_1}.
\]
Here
\[
p_0=1, \qquad p_1=\frac{1}{1+r},
\]
so
\[
\frac{p_0}{p_1}=1+r.
\]
Hence
\[
\frac{u'(x_0)}{\delta u'(x_1)}=1+r,
\]
which rearranges to
\[
u'(x_0)=(1+r)\delta u'(x_1).
\]
\end{proof}

This condition compares the marginal value of consuming one more unit today with the discounted marginal value of consuming one more unit tomorrow, adjusted by the gross return from saving.

\begin{theorem}{Consumption pattern under discounted utility}{ch4-pattern}
Assume $u$ is differentiable and strictly concave, so $u'$ is strictly decreasing.
\begin{enumerate}
    \item If $(1+r)\delta = 1$, then $x_0=x_1$.
    \item If $(1+r)\delta > 1$, then $x_1>x_0$.
    \item If $(1+r)\delta < 1$, then $x_0>x_1$.
\end{enumerate}
\end{theorem}

\begin{proof}
From Theorem \ref{thm:ch4-euler},
\[
u'(x_0)=(1+r)\delta u'(x_1).
\]
If $(1+r)\delta=1$, then
\[
u'(x_0)=u'(x_1),
\]
and strict monotonicity of $u'$ implies
\[
x_0=x_1.
\]
If $(1+r)\delta>1$, then
\[
u'(x_0)>u'(x_1).
\]
Because $u'$ is decreasing, this implies
\[
x_0<x_1.
\]
Similarly, if $(1+r)\delta<1$, then
\[
u'(x_0)<u'(x_1),
\]
so
\[
x_0>x_1.
\]
\end{proof}

\begin{examplebox}[title={Worked Example: logarithmic felicity}]
Suppose
\[
u(x)=\ln x.
\]
Then
\[
u'(x)=\frac{1}{x}.
\]
The first-order condition becomes
\[
\frac{1}{x_0} = (1+r)\delta \frac{1}{x_1},
\]
so
\[
x_1=(1+r)\delta x_0.
\]

If
\[
r=0.1, \qquad \delta=0.9,
\]
then
\[
(1+r)\delta = 1.1\times 0.9 = 0.99,
\]
hence
\[
x_1=0.99x_0.
\]
So optimal future consumption is slightly below present consumption:
\[
\boxed{x_1 < x_0.}
\]
\end{examplebox}

\begin{examtip}
The single most important condition in discounted-utility questions is
\[
u'(x_0)=(1+r)\delta u'(x_1).
\]
Always interpret the product $(1+r)\delta$:
\begin{itemize}
    \item high $\delta$ means patience;
    \item high $r$ makes future consumption relatively cheap;
    \item both push consumption toward the future.
\end{itemize}
\end{examtip}

\subsection{Asset Pricing by No-Arbitrage}

An asset is a stream of future monetary payoffs. Lecture 4 prices assets by the no-arbitrage principle.

\begin{definition}{Arbitrage}{ch4-arbitrage}
An \emph{arbitrage opportunity} is a trading strategy that yields a nonnegative payoff in every contingency and a strictly positive payoff in at least one contingency, with no net cost or riskless profit at the prevailing prices.
\end{definition}

\begin{theorem}{Present-value pricing of a certain asset}{ch4-pv-asset}
Suppose an asset pays deterministic cash flow
\[
y=(y_0,y_1,\dots,y_T).
\]
If borrowing and lending are possible at interest rate $r$ and there is no arbitrage, then the asset's time-0 price must be
\[
P_0 = \sum_{t=0}^{T}\frac{y_t}{(1+r)^t}.
\]
\end{theorem}

\begin{proof}
Let
\[
PV(y)=\sum_{t=0}^{T}\frac{y_t}{(1+r)^t}.
\]

If $P_0<PV(y)$, borrow $P_0$ at time 0 and buy the asset. The asset's cash flows are sufficient to service the debt and still leave a strictly positive surplus in present-value terms, which is an arbitrage.

If $P_0>PV(y)$, short the asset, receive $P_0$ at time 0, invest the proceeds, and use the investment returns to cover the promised asset payments. The excess $P_0-PV(y)$ yields arbitrage.

Hence no-arbitrage requires
\[
P_0=PV(y).
\]
\end{proof}

\begin{definition}{Net present value}{ch4-npv}
If acquiring an asset requires payment stream
\[
m=(m_0,m_1,\dots,m_T)
\]
and yields payoff stream
\[
y=(y_0,y_1,\dots,y_T),
\]
then its \emph{net present value} is
\[
NPV=\sum_{t=0}^{T}\frac{y_t-m_t}{(1+r)^t}.
\]
\end{definition}

A project is attractive at discount rate $r$ if its net present value is positive.

\begin{theorem}{Perpetuity formula}{ch4-perpetuity}
A perpetuity that pays a fixed amount $m$ each period forever, starting in period 1, has present value
\[
PV = \frac{m}{r}.
\]
\end{theorem}

\begin{proof}
Let
\[
PV=\frac{m}{1+r}+\frac{m}{(1+r)^2}+\frac{m}{(1+r)^3}+\cdots.
\]
Multiply both sides by $\frac{1}{1+r}$:
\[
\frac{PV}{1+r}
=
\frac{m}{(1+r)^2}+\frac{m}{(1+r)^3}+\cdots.
\]
Subtract this from the original expression:
\[
PV-\frac{PV}{1+r}=\frac{m}{1+r}.
\]
So
\[
\frac{r}{1+r}PV=\frac{m}{1+r},
\]
hence
\[
PV=\frac{m}{r}.
\]
\end{proof}

\begin{examplebox}[title={Worked Example: net present value of a delayed perpetuity project}]
Suppose a project requires payment of $1$ in each of the first three periods and then pays $2$ forever from period 3 onward. Its net present value is
\[
NPV
=
-1-\frac{1}{1+r}-\frac{1}{(1+r)^2}
+\frac{2}{(1+r)^3}+\frac{2}{(1+r)^4}+\cdots.
\]
The infinite tail is a perpetuity starting at period 3, so
\[
\frac{2}{(1+r)^3}+\frac{2}{(1+r)^4}+\cdots
=
\frac{1}{(1+r)^2}\left(\frac{2}{1+r}+\frac{2}{(1+r)^2}+\cdots\right)
=
\frac{1}{(1+r)^2}\cdot \frac{2}{r}.
\]
Therefore
\[
NPV
=
\frac{2}{r(1+r)^2}
-
\left(
1+\frac{1}{1+r}+\frac{1}{(1+r)^2}
\right).
\]
So the correct compact expression is
\[
\boxed{
NPV
=
\frac{2}{r(1+r)^2}
-
\left(
1+\frac{1}{1+r}+\frac{1}{(1+r)^2}
\right).
}
\]
\end{examplebox}

\begin{examplebox}[title={Worked Example: installment loan interest rate}]
An installment loan gives $m=1000$ upfront and requires monthly payments of $y=100$ for $T=12$ months. The correct monthly interest rate $r$ solves
\[
1000
=
\frac{100}{1+r}
+
\frac{100}{(1+r)^2}
+\cdots+
\frac{100}{(1+r)^{12}}.
\]
This is \emph{not} the same as taking
\[
\frac{12\cdot 100}{1000}-1=20\%.
\]
That simple ratio ignores the fact that payments are made over time rather than all at the end.

Numerically, the solution is approximately
\[
r \approx 0.03
\]
per month, i.e. about $3\%$ monthly. The implied annual rate is therefore much larger than $20\%$; Lecture 4 reports it as roughly $41\%$ per year.

The exam logic is:
\begin{enumerate}
    \item write the present-value equation for the loan;
    \item solve for the periodic interest rate;
    \item do not treat the total repayment as if it occurred only in the final period.
\end{enumerate}
\end{examplebox}

\subsection{Asset Price Changes and Rates of Return}

Suppose an asset does not pay dividends in period 0 and its price rises from $p_0$ to $p_1$ with certainty by next period.

\begin{theorem}{No-arbitrage return condition}{ch4-return}
If the asset pays no interim cash flow and can be traded freely, then no arbitrage implies
\[
r=\frac{p_1-p_0}{p_0}.
\]
Equivalently,
\[
p_0=\frac{p_1}{1+r}.
\]
\end{theorem}

\begin{proof}
If
\[
\frac{p_1-p_0}{p_0}>r,
\]
borrow $p_0$, buy the asset, sell it next period for $p_1$, and repay $(1+r)p_0$. The surplus
\[
p_1-(1+r)p_0>0
\]
is an arbitrage.

If
\[
\frac{p_1-p_0}{p_0}<r,
\]
short the asset, receive $p_0$, invest at rate $r$, and next period buy back the asset for $p_1$. The surplus
\[
(1+r)p_0-p_1>0
\]
is an arbitrage.

Therefore only
\[
\frac{p_1-p_0}{p_0}=r
\]
can hold.
\end{proof}

\begin{policynote}[title={Policy Application: housing as an asset}]
A house pays housing services rather than direct cash. If the monetary rental value of current housing services is $y_0$ and next period's house price is $p_1$, then no-arbitrage implies
\[
p_0 = y_0 + \frac{p_1}{1+r}.
\]
So the value of a house is the value of current housing services plus the discounted resale value.
\end{policynote}

\begin{policynote}[title={Policy Application: depletable resources}]
If a depletable resource such as oil yields no extraction cost and must eventually be replaced by a perfect substitute with cost $C$ at terminal date $T$, then no-arbitrage implies
\[
p_t = \frac{p_{t+1}}{1+r}.
\]
Hence
\[
p_0 = \frac{p_T}{(1+r)^T} = \frac{C}{(1+r)^T}.
\]
This is the basic present-value logic behind the price path of a scarce exhaustible resource.
\end{policynote}

\subsection{Taxes and Asset Returns}

Lecture 4 also studies taxes on interest income and asset returns.

\begin{definition}{After-tax interest rate}{ch4-after-tax}
If the tax rate on interest income is $\tau$, then one dollar saved becomes
\[
1+r(1-\tau)
\]
next period. Thus the effective after-tax interest rate is
\[
r(1-\tau).
\]
\end{definition}

This means taxes reduce the reward to saving and soften the incentive to transfer consumption into the future.

Now consider two assets:
\begin{itemize}
    \item asset $A$ is untaxed;
    \item asset $B$ faces tax rate $\tau$ on its capital gain.
\end{itemize}

\begin{theorem}{Required pre-tax return on a taxed asset}{ch4-taxed-asset}
Suppose asset $B$ has price $p_0^B$ today and price $p_1^B$ next period, and the tax rate on its return is $\tau$. No arbitrage implies
\[
p_0^B = \left(\frac{1-\tau}{1-\tau+r}\right)p_1^B,
\]
equivalently,
\[
\frac{p_1^B-p_0^B}{p_0^B}=\frac{r}{1-\tau}.
\]
\end{theorem}

\begin{proof}
If an investor buys asset $B$ at time 0 for $p_0^B$, then next period the gross proceeds net of tax are
\[
p_1^B - \tau(p_1^B-p_0^B).
\]
No arbitrage requires this to equal the gross return from investing $p_0^B$ at interest rate $r$:
\[
(1+r)p_0^B = p_1^B - \tau(p_1^B-p_0^B).
\]
Rearranging:
\[
(1+r)p_0^B = (1-\tau)p_1^B + \tau p_0^B,
\]
so
\[
(1-\tau+r)p_0^B = (1-\tau)p_1^B.
\]
Hence
\[
p_0^B = \left(\frac{1-\tau}{1-\tau+r}\right)p_1^B.
\]
Subtract $p_0^B$ from both sides of the equivalent expression
\[
(1-\tau+r)p_0^B=(1-\tau)p_1^B
\]
to get
\[
rp_0^B = (1-\tau)(p_1^B-p_0^B),
\]
which gives
\[
\frac{p_1^B-p_0^B}{p_0^B}=\frac{r}{1-\tau}.
\]
\end{proof}

Thus the taxed asset must offer a higher \emph{pre-tax} return than the untaxed asset in order to remain competitive after tax.

\subsection{Chapter Summary}

\begin{keyformula}[title={Key Formula: Lecture 4 Summary}]
\begin{align*}
x_0+s_0 &= y_0, \\
x_1 &= (1+r)s_0+y_1, \\
x_0+\frac{x_1}{1+r} &= y_0+\frac{y_1}{1+r}, \\
p &= \left(1,\frac{1}{1+r}\right), \\
\sum_{t=0}^{T}\frac{x_t}{(1+r)^t} &= \sum_{t=0}^{T}\frac{y_t}{(1+r)^t}, \\
1+\rho &= \frac{1+r}{1+\pi}, \\
U(x_0,x_1) &= u(x_0)+\delta u(x_1), \\
u'(x_0) &= (1+r)\delta u'(x_1), \\
P_0 &= \sum_{t=0}^{T}\frac{y_t}{(1+r)^t}, \\
NPV &= \sum_{t=0}^{T}\frac{y_t-m_t}{(1+r)^t}, \\
PV_{\text{perpetuity}} &= \frac{m}{r}, \\
r &= \frac{p_1-p_0}{p_0}, \\
p_0^B &= \left(\frac{1-\tau}{1-\tau+r}\right)p_1^B.
\end{align*}
\end{keyformula}

\begin{examtip}
The highest-yield skills from Lecture 4 are:
\begin{itemize}
    \item convert intertemporal budgets into present-value form;
    \item interpret $1/(1+r)$ as the price of future consumption;
    \item distinguish savers from borrowers and use price-change logic correctly;
    \item use
    \[
    u'(x_0)=(1+r)\delta u'(x_1)
    \]
    for discounted-utility problems;
    \item compute present values, real interest rates, and net present values cleanly;
    \item apply no-arbitrage to price assets and derive required returns;
    \item adjust return formulas correctly when taxes are imposed.
\end{itemize}
\end{examtip}
\newpage
% chapters/ch05_risk_state_prices_and_expected_utility.tex

\section{Risk, State Prices, and Expected Utility}

This chapter extends consumer theory to uncertain environments. Instead of choosing bundles across goods or across time, the consumer now chooses contingent consumption across states of the world. This allows a unified treatment of insurance, asset markets, diversification, and expected utility. The central ideas are:
\begin{enumerate}
    \item uncertain payoffs can be represented as contingent consumption bundles;
    \item state prices and asset prices are linked by no-arbitrage;
    \item diversification and risk sharing can smooth consumption across states;
    \item expected utility provides a tractable model of risk preferences.
\end{enumerate}

\subsection{Risk and Contingent Consumption}

Suppose there are $S$ mutually exclusive states of the world. State $s$ occurs with probability $\pi_s$, where
\[
\pi_s \ge 0
\qquad \text{and} \qquad
\sum_{s=1}^S \pi_s = 1.
\]

\begin{definition}{Contingent consumption bundle}{ch5-contingent-bundle}
A \emph{contingent consumption bundle} is a vector
\[
x=(x_1,x_2,\dots,x_S)\in \mathbb{R}_+^S,
\]
where $x_s$ is the amount of money or consumption received if state $s$ occurs.
\end{definition}

This is the correct representation for risky prospects. A lottery, job outcome, insurance payoff, or asset return can all be written as a contingent bundle.

\begin{definition}{State prices and the contingent-claims budget set}{ch5-state-prices}
Let
\[
p=(p_1,p_2,\dots,p_S)\in \mathbb{R}_{++}^S
\]
be the \emph{state price vector}, where $p_s$ is the price today of one unit of consumption delivered only in state $s$.

If wealth is $m>0$, the consumer chooses $x$ subject to
\[
p\cdot x = \sum_{s=1}^S p_s x_s \le m.
\]
If the consumer has endowment
\[
\omega=(\omega_1,\dots,\omega_S),
\]
then the budget constraint becomes
\[
p\cdot x \le p\cdot \omega.
\]
\end{definition}

State prices play exactly the same role here that ordinary prices played in Chapters 1--4. The difference is that the coordinates now refer to states rather than goods or dates.

\begin{examplebox}[title={Worked Example: contingent consumption from an uncertain career outcome}]
Suppose there are two states:
\begin{itemize}
    \item state 1: a high-paying job occurs with probability $0.2$;
    \item state 2: onion farming occurs with probability $0.8$.
\end{itemize}
If the corresponding incomes are
\[
10^6 \quad \text{and} \quad 10^3,
\]
then the risky prospect is represented by
\[
x=(10^6,10^3),
\qquad
\pi=(0.2,0.8).
\]
This is a contingent consumption bundle because the realized payoff depends on which state occurs.
\end{examplebox}

\subsection{Insurance and State Prices}

Insurance can be interpreted as trading contingent claims.

\begin{examplebox}[title={Worked Example: insurance as contingent consumption}]
Suppose the consumer is endowed with
\[
\omega=(8,2),
\qquad
\pi=(0.9,0.1),
\]
so state 1 is the good state and state 2 is the bad state. A unit of insurance:
\begin{itemize}
    \item costs $2$ today;
    \item pays $3$ in state 2.
\end{itemize}

Relative to the endowment, buying one unit of insurance reduces state-1 consumption by $2$ and raises state-2 consumption by $1$. If $\theta$ units are purchased, then the resulting contingent bundle is
\[
x(\theta) = (8-2\theta,\, 2+\theta).
\]

If shorting insurance is allowed, then $\theta$ may be negative. For example,
\[
\theta=-2
\]
gives
\[
x(-2)=(12,0).
\]

Now observe that
\[
x_1 + 2x_2 = (8-2\theta)+2(2+\theta)=12
\]
for every $\theta$. Since
\[
12 = 8 + 2\cdot 2 = (1,2)\cdot (8,2),
\]
this budget line is exactly the same as the contingent-claims budget with state price vector
\[
\boxed{p=(1,2).}
\]

So insurance trading is equivalent to buying state-contingent consumption at prices $(1,2)$.
\end{examplebox}

\begin{examtip}
When an insurance contract changes payoffs by a fixed amount across states, rewrite the set of attainable bundles as a linear equation. That linear equation usually reveals the implied state prices immediately.
\end{examtip}

\subsection{Assets, Portfolios, and Arbitrage}

Assets can also be viewed as contingent consumption bundles.

\begin{definition}{Assets, portfolios, and shorting}{ch5-assets}
An asset is a contingent payoff vector
\[
a^k=(a_1^k,\dots,a_S^k)\in \mathbb{R}^S.
\]
If its price is $q_k$, then holding one unit of asset $k$ costs $q_k$ today and delivers $a_s^k$ in state $s$.

A \emph{portfolio} is a vector
\[
\theta=(\theta_1,\dots,\theta_K),
\]
where $\theta_k$ is the quantity of asset $k$ held.

The portfolio costs
\[
q\cdot \theta = \sum_{k=1}^K q_k \theta_k
\]
and delivers contingent consumption
\[
\sum_{k=1}^K \theta_k a^k.
\]

If $\theta_k<0$ is allowed, then shorting is allowed.
\end{definition}

\begin{definition}{Arbitrage}{ch5-arbitrage}
A portfolio $\theta$ is an \emph{arbitrage opportunity} if:
\begin{enumerate}
    \item it costs nothing or yields money upfront:
    \[
    q\cdot \theta \le 0;
    \]
    \item its payoff is nonnegative in every state:
    \[
    \sum_{k=1}^K \theta_k a^k \ge 0;
    \]
    \item and it is strictly positive in at least one state.
\end{enumerate}
\end{definition}

Arbitrage matters because no-arbitrage is the key pricing principle of the chapter: if two portfolios have the same state-contingent payoffs, they must have the same price.

\begin{theorem}{Pricing by replication}{ch5-replication}
Suppose a new asset $\tilde a$ can be replicated by an existing portfolio $\theta$, i.e.
\[
\tilde a = \sum_{k=1}^K \theta_k a^k.
\]
If there is no arbitrage, then the price of the new asset must be
\[
\tilde q = q\cdot \theta.
\]
\end{theorem}

\begin{proof}
If $\tilde q < q\cdot \theta$, then buy the new asset and short the replicating portfolio. This gives a positive profit today and zero net payoff in every state, which is an arbitrage.

If $\tilde q > q\cdot \theta$, then short the new asset and buy the replicating portfolio. Again, the payoffs cancel state by state, leaving a sure profit today.

So no-arbitrage requires
\[
\tilde q = q\cdot \theta.
\]
\end{proof}

\begin{definition}{State-price representation}{ch5-state-representation}
If Arrow securities exist, where:
\[
e^1=(1,0,\dots,0),\ e^2=(0,1,\dots,0),\ \dots,\ e^S=(0,\dots,0,1),
\]
and their prices are
\[
p=(p_1,\dots,p_S),
\]
then any asset
\[
a=(a_1,\dots,a_S)
\]
has no-arbitrage price
\[
q(a)=\sum_{s=1}^S p_s a_s = p\cdot a.
\]
\end{definition}

So state prices are the building blocks of all asset prices.

\begin{examplebox}[title={Worked Example: extracting state prices from asset prices}]
Suppose there are two assets:
\[
a^1=(1,2), \qquad a^2=(1,1),
\]
with prices
\[
q_1=3, \qquad q_2=2.
\]

To find the price of one unit of state-1 consumption, construct the Arrow security $(1,0)$. Observe that
\[
2a^2-a^1 = 2(1,1)-(1,2) = (1,0).
\]
Hence its no-arbitrage price is
\[
2q_2-q_1 = 2\cdot 2 - 3 = 1.
\]
So
\[
p_1=1.
\]

To find the price of one unit of state-2 consumption, construct
\[
a^1-a^2 = (1,2)-(1,1) = (0,1),
\]
whose price is
\[
q_1-q_2 = 3-2 = 1.
\]
So
\[
p_2=1.
\]

Therefore the state price vector is
\[
\boxed{p=(1,1).}
\]
\end{examplebox}

\begin{examplebox}[title={Worked Example: contingent claims from a risky and a risk-free asset}]
Suppose there are two assets:
\begin{itemize}
    \item a risk-free asset with payoff
    \[
    a^f=(1,1)
    \]
    and price
    \[
    q_f=3;
    \]
    \item a risky asset with payoff
    \[
    a^r=(2,1)
    \]
    and price
    \[
    q_r=4.
    \]
\end{itemize}
If wealth is
\[
m=12,
\]
then spending all wealth on the risk-free asset buys
\[
\frac{12}{3}=4
\]
units, which yields
\[
(4,4).
\]
Spending all wealth on the risky asset buys
\[
\frac{12}{4}=3
\]
units, which yields
\[
(6,3).
\]

If a fraction $\alpha$ of wealth is spent on the risky asset and $1-\alpha$ on the risk-free asset, the payoff is
\[
(4,4)+\alpha\bigl((6,3)-(4,4)\bigr)
=
(4,4)+\alpha(2,-1).
\]
So each increase of $2$ units in state-1 payoff costs a reduction of $1$ unit in state-2 payoff. This is exactly the tradeoff implied by state prices
\[
\boxed{p=(1,2).}
\]
\end{examplebox}

\begin{commonmistake}
Do not confuse asset prices with state prices. Asset prices apply to entire payoff vectors; state prices apply to one unit of payoff in a single state. Asset prices are recovered from state prices by
\[
q(a)=p\cdot a.
\]
\end{commonmistake}

\subsection{Diversification and Risk Sharing}

Asset markets allow agents to reallocate consumption across states. This is valuable when people dislike large payoff fluctuations.

\begin{definition}{Diversification}{ch5-diversification}
Diversification means holding a portfolio of assets rather than concentrating wealth in one risky payoff. The aim is to reduce the variance of payoffs across states without sacrificing too much expected return.
\end{definition}

The most effective diversification combines assets with offsetting payoffs, ideally negatively correlated assets.

\begin{examplebox}[title={Worked Example: diversification can create a safe payoff}]
Suppose
\[
a=(1,0),
\qquad
a'=(0,1),
\qquad
\pi=\left(\frac12,\frac12\right).
\]
Each asset alone is risky:
\begin{itemize}
    \item $a$ pays only in state 1;
    \item $a'$ pays only in state 2.
\end{itemize}
But the portfolio
\[
a+a'=(1,1)
\]
is risk-free.

A second example: let
\[
\tilde a=(1,0),
\qquad
\tilde a'=(2,1).
\]
Both are risky, but if shorting is allowed,
\[
\tilde a' - \tilde a = (2,1)-(1,0)=(1,1),
\]
which is again risk-free.

So diversification works not only by buying several assets, but also by combining long and short positions when markets permit it.
\end{examplebox}

\begin{policynote}[title={Policy Application: village risk sharing}]
Suppose each villager faces idiosyncratic crop-failure risk. If one villager's crops fail while others' do not, the villagers can insure one another by trading state-contingent claims. This is risk sharing: each person gives up some consumption in good states in exchange for more consumption in bad states.

However, if \emph{everyone} is hit by the same shock, such as a flood or drought, then internal insurance becomes ineffective. This is \emph{correlated risk}. It explains why some risks, especially large-scale natural-disaster risks, are difficult or expensive to insure.
\end{policynote}

\subsection{Expected Value and Expected Utility}

To rank contingent bundles, the lecture introduces expected value and expected utility.

\begin{definition}{Expected value}{ch5-expected-value}
Given contingent bundle
\[
x=(x_1,\dots,x_S)
\]
with probabilities
\[
\pi=(\pi_1,\dots,\pi_S),
\]
its expected value is
\[
E[x] = \sum_{s=1}^S \pi_s x_s.
\]
\end{definition}

Expected value is a natural benchmark, but by itself it does not capture attitudes toward risk. Many people prefer a safer bundle with lower expected value to a riskier bundle with higher expected value.

\begin{definition}{Expected utility and Bernoulli utility}{ch5-expected-utility}
Let
\[
u:\mathbb{R}_+ \to \mathbb{R}
\]
be a Bernoulli utility function. The expected utility of contingent bundle
\[
x=(x_1,\dots,x_S)
\]
is
\[
U(x)=E[u(x)] = \sum_{s=1}^S \pi_s u(x_s).
\]
\end{definition}

The lecture assumes that $u$ is increasing, and often also concave.

\begin{definition}{Increasing and concave Bernoulli utility}{ch5-inc-concave}
A Bernoulli utility function $u$ is:
\begin{itemize}
    \item \emph{increasing} if more money is better:
    \[
    x>x' \implies u(x)>u(x');
    \]
    \item \emph{concave} if marginal utility is diminishing.
\end{itemize}
If $u$ is differentiable and concave, then $u'$ is decreasing.
\end{definition}

\begin{theorem}{Interior optimality under expected utility}{ch5-focexpected}
Suppose the consumer maximizes
\[
U(x)=\sum_{s=1}^S \pi_s u(x_s)
\]
subject to
\[
p\cdot x \le m,
\]
and the optimum is interior. Then for any two states $i,j$,
\[
\frac{\pi_i u'(x_i)}{\pi_j u'(x_j)} = \frac{p_i}{p_j}.
\]
\end{theorem}

\begin{proof}
The marginal utility of state-$s$ consumption is
\[
MU_s = \pi_s u'(x_s).
\]
At an interior optimum, the tangency condition from earlier chapters gives
\[
\frac{MU_i}{MU_j} = \frac{p_i}{p_j}.
\]
Substituting the marginal utilities yields
\[
\frac{\pi_i u'(x_i)}{\pi_j u'(x_j)} = \frac{p_i}{p_j}.
\]
\end{proof}

This can be rewritten as
\[
\frac{u'(x_i)}{u'(x_j)} = \frac{p_i/\pi_i}{p_j/\pi_j}.
\]
So what matters is the price of state consumption relative to its probability.

\begin{examplebox}[title={Worked Example: optimal contingent consumption with logarithmic utility}]
Suppose there are two states with probabilities
\[
\pi_1=\pi_2=\frac12,
\]
state prices
\[
p=(1,2),
\]
wealth
\[
m=12,
\]
and Bernoulli utility
\[
u(x)=\ln x.
\]
Then expected utility is
\[
U(x_1,x_2)=\frac12 \ln x_1 + \frac12 \ln x_2.
\]

At an interior optimum,
\[
\frac{\pi_1 u'(x_1)}{\pi_2 u'(x_2)}=\frac{p_1}{p_2}.
\]
Since
\[
u'(x)=\frac{1}{x}
\quad \text{and} \quad
\pi_1=\pi_2,
\]
this becomes
\[
\frac{1/x_1}{1/x_2} = \frac{1}{2}
\quad \Longrightarrow \quad
\frac{x_2}{x_1}=\frac12.
\]
So
\[
x_2=\frac{x_1}{2}.
\]
Using the budget constraint
\[
x_1 + 2x_2 = 12,
\]
we obtain
\[
x_1 + 2\left(\frac{x_1}{2}\right)=12
\quad \Longrightarrow \quad
2x_1=12
\quad \Longrightarrow \quad
x_1=6,
\]
and hence
\[
x_2=3.
\]

Therefore the optimal contingent bundle is
\[
\boxed{x=(6,3).}
\]
\end{examplebox}

\subsection{First-Order Stochastic Dominance}

Expected utility with increasing Bernoulli utility respects a natural ranking of risky prospects.

\begin{definition}{Survival probability and first-order stochastic dominance}{ch5-fosd}
Given contingent bundle
\[
x=(x_1,\dots,x_S)
\]
and probability vector $\pi$, define
\[
\pi(x,\alpha)
\]
to be the probability that $x$ pays at least $\alpha$.

We say that $x$ \emph{first-order stochastically dominates} $x'$, written
\[
x \succeq_{FOSD} x',
\]
if
\[
\pi(x,\alpha) \ge \pi(x',\alpha)
\qquad \text{for every } \alpha.
\]
\end{definition}

This means that for every threshold $\alpha$, the probability of receiving at least $\alpha$ is weakly higher under $x$ than under $x'$.

\begin{theorem}{FOSD theorem}{ch5-fosd-theorem}
Suppose preferences take the expected-utility form
\[
U(x)=\sum_{s=1}^S \pi_s u(x_s),
\]
where $u$ is increasing. If
\[
x \succeq_{FOSD} x',
\]
then
\[
U(x)\ge U(x').
\]
\end{theorem}

So FOSD is a robust welfare ranking: every consumer with increasing expected utility agrees that the dominating bundle is at least as good.

\begin{examplebox}[title={Worked Example: checking first-order stochastic dominance}]
Let
\[
\pi=\left(\frac14,\frac14,\frac12\right),
\qquad
x=(7,5,1),
\qquad
x'=(1,0,4).
\]
To test whether
\[
x \succeq_{FOSD} x',
\]
it is enough to check the payoff thresholds appearing in $x'$.

For $\alpha=1$:
\[
\pi(x,1)=1,
\qquad
\pi(x',1)=\frac34.
\]

For $\alpha=0$:
\[
\pi(x,0)=1,
\qquad
\pi(x',0)=1.
\]

For $\alpha=4$:
\[
\pi(x,4)=\frac12,
\qquad
\pi(x',4)=\frac12.
\]

Thus
\[
\pi(x,\alpha)\ge \pi(x',\alpha)
\]
for all relevant thresholds, so
\[
\boxed{x \succeq_{FOSD} x'.}
\]
Therefore every consumer with increasing expected utility weakly prefers $x$ to $x'$.
\end{examplebox}

\subsection{Risk Attitudes}

Concavity of Bernoulli utility captures dislike of risk.

\begin{theorem}{Jensen inequality for expected utility}{ch5-jensen}
If $u$ is concave, then for every contingent bundle $x$,
\[
u(E[x]) \ge E[u(x)].
\]
Equivalently,
\[
u(E[x]) \ge U(x).
\]
\end{theorem}

\begin{proof}
This is Jensen's inequality applied to the concave function $u$.
\end{proof}

The economic meaning is crucial: if utility is concave, then receiving the expected value for sure is at least as good as facing the risky prospect itself.

\begin{definition}{Risk attitudes}{ch5-risk-attitudes}
Let
\[
U(x)=E[u(x)].
\]
Then:
\begin{itemize}
    \item the consumer is \emph{risk averse} if
    \[
    u(E[x]) \ge U(x)
    \quad \text{for all } x;
    \]
    \item the consumer is \emph{risk neutral} if
    \[
    u(E[x]) = U(x)
    \quad \text{for all } x;
    \]
    \item the consumer is \emph{risk loving} if
    \[
    u(E[x]) \le U(x)
    \quad \text{for all } x.
    \]
\end{itemize}
\end{definition}

\begin{theorem}{Linear utility implies risk neutrality}{ch5-risk-neutral}
A consumer is risk neutral if the Bernoulli utility function is affine:
\[
u(x)=ax+b
\qquad \text{with } a>0.
\]
In particular, $u(x)=x$ is risk neutral.
\end{theorem}

\begin{proof}
If
\[
u(x)=ax+b,
\]
then
\[
U(x)=E[u(x)] = E[ax+b] = aE[x]+b = u(E[x]).
\]
So the consumer is risk neutral.
\end{proof}

\begin{examplebox}[title={Worked Example: same expected value, different risk}]
Suppose
\[
\pi=\left(\frac12,\frac12\right).
\]
Compare the safe bundle
\[
x=(10^7,10^7)
\]
with the risky bundle
\[
x'=(4\times 10^7,0).
\]
Their expected values are
\[
E[x]=10^7,
\qquad
E[x']=2\times 10^7.
\]
So $x'$ has the higher expected value.

However, if the consumer has increasing and concave Bernoulli utility, expected utility need not follow expected value. A sufficiently risk-averse consumer may prefer the certain bundle $x$ to the much riskier $x'$ despite its lower expected value.

This is exactly why expected utility rather than expected value is the central decision criterion in the lecture.
\end{examplebox}

\begin{examplebox}[title={Worked Example: certainty equivalent under concavity}]
Suppose
\[
\pi=\left(\frac12,\frac14,\frac14\right),
\qquad
x=(4,8,12).
\]
Then the expected value is
\[
E[x]=\frac12\cdot 4 + \frac14\cdot 8 + \frac14\cdot 12 = 2+2+3=7.
\]
If $u$ is increasing and concave, Jensen's inequality gives
\[
u(7) = u(E[x]) \ge U(x).
\]
So the consumer weakly prefers receiving $7$ for sure to receiving the contingent bundle $(4,8,12)$ itself.

This is the standard exam interpretation of concavity:
\[
\boxed{\text{concavity } \Longleftrightarrow \text{ risk aversion}.}
\]
\end{examplebox}

\begin{commonmistake}
Do not conclude that a bundle with higher expected value must be preferred. That is true only for risk-neutral consumers. Under concave expected utility, a lower-mean but less risky bundle can be preferred.
\end{commonmistake}

\subsection{Chapter Summary}

\begin{keyformula}[title={Key Formula: Lecture 5 Summary}]
\begin{align*}
x &= (x_1,\dots,x_S) && \text{contingent consumption bundle}, \\
p\cdot x &\le m && \text{state-contingent budget}, \\
q\cdot \theta &= \sum_{k=1}^K q_k \theta_k && \text{portfolio cost}, \\
\sum_{k=1}^K \theta_k a^k &&& \text{portfolio payoff}, \\
q(\tilde a) &= q\cdot \theta && \text{if } \tilde a=\sum_{k=1}^K \theta_k a^k, \\
q(a) &= p\cdot a && \text{state-price representation}, \\
E[x] &= \sum_{s=1}^S \pi_s x_s, \\
U(x) &= E[u(x)] = \sum_{s=1}^S \pi_s u(x_s), \\
\frac{\pi_i u'(x_i)}{\pi_j u'(x_j)} &= \frac{p_i}{p_j} && \text{interior optimum}, \\
x \succeq_{FOSD} x' &\Longrightarrow U(x)\ge U(x') && \text{if } u \text{ is increasing}, \\
u(E[x]) &\ge E[u(x)] && \text{if } u \text{ is concave}.
\end{align*}
\end{keyformula}

\begin{examtip}
The highest-yield skills from Lecture 5 are:
\begin{itemize}
    \item represent uncertain outcomes as contingent consumption bundles;
    \item translate insurance or asset trading into state-contingent budgets;
    \item identify arbitrage and use no-arbitrage to price replicated assets;
    \item recover state prices from asset prices in simple two-state settings;
    \item use expected utility rather than expected value when risk attitudes matter;
    \item test first-order stochastic dominance carefully using payoff thresholds;
    \item interpret concavity of Bernoulli utility as risk aversion.
\end{itemize}
\end{examtip}
\newpage
% chapters/ch06_part1_summary.tex

\section{Part I Summary}

This chapter consolidates the main ideas from Part I. The goal is not to re-teach Chapters 1--5 in full, but to compress the most examinable material into a revision-oriented framework. The unifying theme is that all five chapters study \emph{choice under constraints}: first with ordinary consumption bundles, then with price changes and endowments, then across time, and finally across risky states of the world.

\subsection{Core Definitions and Frameworks}

\begin{definition}{Part I master framework}{ch6-master-framework}
A Part I consumer problem can usually be written in the form
\[
\max_{x \in X} U(x)
\quad \text{subject to} \quad
p \cdot x \le m,
\]
or, with an endowment $\omega$,
\[
\max_{x \in X} U(x)
\quad \text{subject to} \quad
p \cdot x \le p \cdot \omega.
\]
The interpretation of $x$, $p$, and $\omega$ changes by topic:
\begin{itemize}
    \item \textbf{ordinary consumption:} $x$ is a bundle of goods;
    \item \textbf{endowments:} $\omega$ is the initial bundle owned;
    \item \textbf{intertemporal choice:} coordinates of $x$ are dated consumptions;
    \item \textbf{contingent consumption:} coordinates of $x$ are state-contingent payoffs.
\end{itemize}
\end{definition}

\begin{keyformula}[title={Key Formula: Part I notation map}]
\[
x=(x_1,\dots,x_L)\in \mathbb{R}_+^L,
\qquad
p=(p_1,\dots,p_L)\in \mathbb{R}_{++}^L,
\qquad
p\cdot x=\sum_{\ell=1}^L p_\ell x_\ell.
\]
Core interpretations:
\begin{itemize}
    \item $x$: chosen bundle / plan / payoff;
    \item $p$: price vector / discounted price vector / state price vector;
    \item $m$: expenditure or wealth;
    \item $\omega$: endowment;
    \item $z=x-\omega$: net demand.
\end{itemize}
\end{keyformula}

\begin{examtip}
A large fraction of Part I can be solved by first identifying exactly what the coordinates of $x$ mean. Once that is clear, the problem is usually just a standard consumer-choice problem with a suitably interpreted price vector.
\end{examtip}

\subsection{High-Yield Definitions and Results}

\subsubsection{Preferences and revealed preference}

\begin{keyformula}[title={Key Formula: preference language}]
\[
x \succeq y \quad \text{weak preference}, \qquad
x \succ y \quad \text{strict preference}, \qquad
x \sim y \quad \text{indifference}.
\]
A preference relation is usually assumed to be:
\begin{itemize}
    \item \textbf{complete};
    \item \textbf{transitive};
    \item \textbf{increasing:} if $x>y$, then $x\succ y$.
\end{itemize}
\end{keyformula}

\begin{definition}{Revealed preference notation}{ch6-rp-def}
If bundle $x^t$ is chosen at prices $p^t$, then:
\[
x^t R x
\iff
p^t \cdot x^t \ge p^t \cdot x
\]
means $x$ was affordable when $x^t$ was chosen, and
\[
x^t P x
\iff
p^t \cdot x^t > p^t \cdot x
\]
means $x$ was strictly cheaper than the chosen bundle.
\end{definition}

\begin{keyformula}[title={Key Formula: WARP and GARP}]
If choices maximize an increasing utility function, then:
\[
x^t R x \implies U(x^t)\ge U(x),
\qquad
x^t P x \implies U(x^t)>U(x).
\]
Hence:

\textbf{WARP:}
\[
x^t R x^s \implies \text{not }(x^s P x^t).
\]

\textbf{GARP:}
if
\[
x^t R^* x^s,
\]
then not
\[
x^s P x^t.
\]
Equivalently, utility maximization with increasing utility forbids a weak revealed-preference chain in one direction together with a strict reversal.
\end{keyformula}

\begin{examtip}
For revealed-preference questions:
\begin{enumerate}
    \item compute each observed expenditure $p^t\cdot x^t$;
    \item test cross-affordability $p^t\cdot x^s$;
    \item translate into $R$ or $P$;
    \item check for WARP or GARP violations.
\end{enumerate}
Never start with utility before checking the affordability relations.
\end{examtip}

\subsubsection{Demand response to price and expenditure changes}

\begin{keyformula}[title={Key Formula: demand classifications}]
For Marshallian demand $x(p,m)$:
\begin{itemize}
    \item \textbf{ordinary good 1:}
    \[
    p_1 \uparrow \implies x_1 \downarrow;
    \]
    \item \textbf{Giffen good 1:}
    \[
    p_1 \uparrow \implies x_1 \uparrow;
    \]
    \item \textbf{gross substitutes:}
    \[
    p_1 \uparrow \implies x_2 \uparrow;
    \]
    \item \textbf{gross complements:}
    \[
    p_1 \uparrow \implies x_2 \downarrow;
    \]
    \item \textbf{normal good 1:}
    \[
    m \uparrow \implies x_1 \uparrow;
    \]
    \item \textbf{inferior good 1:}
    \[
    m \uparrow \implies x_1 \downarrow.
    \]
\end{itemize}
\end{keyformula}

\begin{keyformula}[title={Key Formula: Slutsky decomposition}]
If good 1's price rises from $p_1$ to $p_1'>p_1$, with $p'=(p_1',p_2)$, then
\[
\Delta x_1 = x_1(p',m)-x_1(p,m),
\]
\[
m^s = p' \cdot x(p,m),
\]
\[
\Delta x_1^s = x_1(p',m^s)-x_1(p,m),
\qquad
\Delta x_1^n = x_1(p',m)-x_1(p',m^s),
\]
and
\[
\Delta x_1 = \Delta x_1^s + \Delta x_1^n.
\]
Always:
\[
\Delta x_1^s \le 0
\quad \text{when } p_1 \uparrow.
\]
If the good is normal, then also
\[
\Delta x_1^n \le 0.
\]
Therefore normal goods satisfy the law of demand:
\[
p_1 \uparrow \implies x_1 \downarrow.
\]
\end{keyformula}

\begin{keyformula}[title={Key Formula: Hicksian substitution}]
For the same price increase, let $u_0=U(x(p,m))$ and let $m^h$ be the minimum expenditure at prices $p'$ needed to attain utility $u_0$. Then
\[
\Delta x_1 = \Delta x_1^h + \Delta x_1^q,
\]
where
\[
\Delta x_1^h = h_1(p',u_0)-x_1(p,m) \le 0.
\]
Also
\[
m^h \le m^s.
\]
\end{keyformula}

\subsubsection{Marginal utility, MRS, and interior optimality}

\begin{keyformula}[title={Key Formula: MRS toolkit}]
For differentiable utility $U(x_1,x_2)$:
\[
MU_1=\frac{\partial U}{\partial x_1},
\qquad
MU_2=\frac{\partial U}{\partial x_2}.
\]
At an interior optimum with budget line
\[
p_1x_1+p_2x_2=m,
\]
the tangency condition is
\[
MRS_{12}=\frac{MU_1}{MU_2}=\frac{p_1}{p_2}.
\]
\end{keyformula}

\begin{examtip}
The tangency condition is only valid for an \emph{interior} optimum. If the tangency candidate is infeasible or gives negative demand, switch immediately to checking corners.
\end{examtip}

\subsubsection{Endowments and net demand}

\begin{definition}{Net demand}{ch6-net-demand-def}
Given endowment $\omega$ and chosen bundle $x$,
\[
z=x-\omega.
\]
Equivalently, if $x(p,m)$ is Marshallian demand and wealth comes from selling the endowment,
\[
z(p)=x(p,p\cdot \omega)-\omega.
\]
Interpretation:
\[
z_\ell>0 \Rightarrow \text{net buyer of good }\ell,
\qquad
z_\ell<0 \Rightarrow \text{net seller of good }\ell.
\]
\end{definition}

\begin{keyformula}[title={Key Formula: endowment budget and Walras' law}]
Affordability with endowments:
\[
p\cdot x \le p\cdot \omega
\quad \Longleftrightarrow \quad
p\cdot z \le 0.
\]
At an optimum with increasing utility,
\[
p\cdot x = p\cdot \omega
\quad \Longrightarrow \quad
p\cdot z = 0.
\]
\end{keyformula}

\begin{keyformula}[title={Key Formula: endowment price-change logic}]
When the price of good 1 rises:
\begin{itemize}
    \item substitution effect still satisfies
    \[
    \Delta x_1^s \le 0;
    \]
    \item the sign of the income effect depends on whether the consumer is a buyer or seller of good 1.
\end{itemize}
For normal demand:
\begin{center}
\renewcommand{\arraystretch}{1.2}
\begin{tabular}{|c|c|c|c|}
\hline
Change & Status in good 1 & Income effect & Definite total sign? \\
\hline
$p_1 \uparrow$ & buyer & $\le 0$ & yes: $\Delta x_1 \le 0$ \\
\hline
$p_1 \uparrow$ & seller & $\ge 0$ & ambiguous \\
\hline
$p_1 \downarrow$ & buyer & $\ge 0$ & yes: $\Delta x_1 \ge 0$ \\
\hline
$p_1 \downarrow$ & seller & $\le 0$ & ambiguous \\
\hline
\end{tabular}
\end{center}
\end{keyformula}

\begin{examtip}
Under endowments, the most important first step is to classify the consumer as a \emph{net buyer} or \emph{net seller} of the good whose price changes. Without that step, the sign of the income effect is usually wrong.
\end{examtip}

\subsubsection{Intertemporal consumption}

\begin{keyformula}[title={Key Formula: intertemporal budget}]
For periods $t=0,\dots,T$ and interest rate $r\ge 0$,
\[
\sum_{t=0}^{T} \frac{x_t}{(1+r)^t}
=
\sum_{t=0}^{T} \frac{y_t}{(1+r)^t}.
\]
With
\[
p_t=\left(\frac{1}{1+r}\right)^t,
\]
this is simply
\[
p\cdot x = p\cdot y.
\]
\end{keyformula}

\begin{definition}{Intertemporal net demand}{ch6-intertemporal-net}
For income stream $y=(y_0,\dots,y_T)$ and chosen stream $x=(x_0,\dots,x_T)$,
\[
z=x-y.
\]
If $T=1$, then:
\begin{itemize}
    \item $z_0<0$ means saving in period 0;
    \item $z_0>0$ means borrowing in period 0;
    \item $z_1$ records the opposite side of that saving/borrowing decision in period 1.
\end{itemize}
\end{definition}

\begin{keyformula}[title={Key Formula: interest-rate comparative statics}]
In the two-period case,
\[
p_1=\frac{1}{1+r}.
\]
So:
\begin{itemize}
    \item $r \uparrow$ means future consumption becomes cheaper;
    \item a saver (net buyer of future consumption) stays a saver;
    \item if $r \downarrow$, a borrower (net seller of future consumption) stays a borrower.
\end{itemize}
\end{keyformula}

\begin{keyformula}[title={Key Formula: discounted utility}]
For
\[
U(x_0,x_1)=u(x_0)+\delta u(x_1),
\]
the interior first-order condition is
\[
u'(x_0)=(1+r)\delta u'(x_1).
\]
Hence:
\begin{itemize}
    \item if $(1+r)\delta=1$, then $x_0=x_1$;
    \item if $(1+r)\delta>1$, then $x_1>x_0$;
    \item if $(1+r)\delta<1$, then $x_0>x_1$.
\end{itemize}
\end{keyformula}

\subsubsection{Contingent consumption, state prices, and risk}

\begin{definition}{Contingent consumption and expected utility}{ch6-risk-def}
With $S$ states, probabilities $\pi=(\pi_1,\dots,\pi_S)$, and contingent consumption
\[
x=(x_1,\dots,x_S),
\]
expected utility is
\[
U(x)=\sum_{s=1}^S \pi_s u(x_s).
\]
\end{definition}

\begin{keyformula}[title={Key Formula: state prices and asset pricing}]
If $p=(p_1,\dots,p_S)$ is the state-price vector, then a contingent bundle $x$ costs
\[
p\cdot x.
\]
If an asset has payoff vector
\[
a=(a_1,\dots,a_S),
\]
then its no-arbitrage price is
\[
q(a)=p\cdot a.
\]
More generally, if a new asset can be replicated by portfolio $\theta$, its price must equal the cost of the replicating portfolio.
\end{keyformula}

\begin{keyformula}[title={Key Formula: expected value, FOSD, and risk attitudes}]
Expected value:
\[
E[x]=\sum_{s=1}^S \pi_s x_s.
\]

Interior FOC under expected utility:
\[
\frac{\pi_i u'(x_i)}{\pi_j u'(x_j)}=\frac{p_i}{p_j}.
\]

First-order stochastic dominance:
\[
x \succeq_{FOSD} x'
\quad \Longrightarrow \quad
U(x)\ge U(x')
\]
for all increasing $u$.

If $u$ is concave, then
\[
u(E[x]) \ge E[u(x)] = U(x).
\]
Interpretation: the consumer is risk averse.

Risk classifications:
\[
u(E[x]) \ge U(x) \Rightarrow \text{risk averse},
\]
\[
u(E[x]) = U(x) \Rightarrow \text{risk neutral},
\]
\[
u(E[x]) \le U(x) \Rightarrow \text{risk loving}.
\]
\end{keyformula}

\subsection{Comparison of Central Models}

\begin{center}
\renewcommand{\arraystretch}{1.25}
\begin{tabular}{|p{3.0cm}|p{3.4cm}|p{4.2cm}|p{4.0cm}|}
\hline
\textbf{Topic} & \textbf{Choice vector} & \textbf{Price interpretation} & \textbf{Core technique} \\
\hline
Revealed preference & ordinary bundle $x$ & market prices $p$ & affordability tests, $R/P$, WARP, GARP \\
\hline
Price changes & Marshallian demand $x(p,m)$ & relative goods prices & Slutsky decomposition \\
\hline
Endowments & consumption $x$ with endowment $\omega$ & goods prices applied to both $x$ and $\omega$ & net demand $z=x-\omega$, buyer/seller classification \\
\hline
Intertemporal choice & dated consumption $(x_0,\dots,x_T)$ & discounted prices $p_t=(1+r)^{-t}$ & present value, saving/borrowing, Euler condition \\
\hline
Risk / contingent claims & state payoffs $(x_1,\dots,x_S)$ & state prices $p_s$ & expected utility, no-arbitrage pricing, FOSD \\
\hline
\end{tabular}
\end{center}

\begin{examtip}
A strong way to remember Part I is:
\[
\text{ordinary goods} \rightarrow \text{endowments} \rightarrow \text{time} \rightarrow \text{states}.
\]
The mathematics stays almost the same. Only the economic interpretation of each coordinate changes.
\end{examtip}

\subsection{Standard Problem Types and Solution Patterns}

\subsubsection{Type A: revealed-preference consistency}

\begin{examtip}[title={Method Pattern: WARP / GARP problems}]
Use the following template:
\begin{enumerate}
    \item list each observation $(p^t,x^t)$;
    \item compute
    \[
    p^t \cdot x^t
    \quad \text{and} \quad
    p^t \cdot x^s
    \]
    for other bundles $x^s$;
    \item infer $R$ or $P$ relations;
    \item look for:
    \[
    x^t R x^s \text{ and } x^s P x^t
    \]
    for WARP, or
    \[
    x^t R^* x^s \text{ and } x^s P x^t
    \]
    for GARP.
\end{enumerate}
\end{examtip}

\subsubsection{Type B: price change and demand decomposition}

\begin{examtip}[title={Method Pattern: Slutsky questions}]
When a price changes:
\begin{enumerate}
    \item compute the original demand $x(p,m)$;
    \item compute the compensated expenditure
    \[
    m^s = p' \cdot x(p,m);
    \]
    \item compute the compensated demand $x(p',m^s)$;
    \item split the total effect into substitution and income effects;
    \item use normal/inferior status to sign the income effect.
\end{enumerate}
\end{examtip}

\subsubsection{Type C: endowment comparative statics}

\begin{examtip}[title={Method Pattern: endowment questions}]
Always do this first:
\[
z = x-\omega.
\]
Then identify whether the consumer is a buyer or seller of the good whose price changes. Only after that should you analyze the substitution and income effects.
\end{examtip}

\subsubsection{Type D: intertemporal optimization}

\begin{examtip}[title={Method Pattern: intertemporal questions}]
Standard order:
\begin{enumerate}
    \item convert the stream budget into present-value form;
    \item rewrite as a standard consumer problem with prices
    \[
    p_t=(1+r)^{-t};
    \]
    \item if discounted utility is given, apply
    \[
    u'(x_0)=(1+r)\delta u'(x_1);
    \]
    \item use the budget equation to solve.
\end{enumerate}
\end{examtip}

\subsubsection{Type E: risk and asset pricing}

\begin{examtip}[title={Method Pattern: contingent consumption / assets}]
When a risk problem involves insurance or asset trade:
\begin{enumerate}
    \item write the final contingent consumption bundle;
    \item identify the linear tradeoff across states;
    \item infer the implied state prices or use no-arbitrage replication;
    \item if expected utility is given, solve
    \[
    \frac{\pi_i u'(x_i)}{\pi_j u'(x_j)}=\frac{p_i}{p_j}.
    \]
\end{enumerate}
\end{examtip}

\subsection{Common Mistakes and Exam Traps}

\begin{commonmistake}
In revealed-preference questions, do not mix up the meanings of $R$ and $P$. The lecture convention is:
\[
x^t R x \iff x \text{ was affordable when } x^t \text{ was chosen},
\]
\[
x^t P x \iff x \text{ was strictly cheaper than the chosen bundle}.
\]
\end{commonmistake}

\begin{commonmistake}
Do not conclude that a good is normal from a negative own-price effect. Normality is about the sign of the \emph{expenditure} effect, not the total price effect.
\end{commonmistake}

\begin{commonmistake}
Do not apply the tangency condition automatically. It requires an interior optimum. If the candidate is infeasible or violates nonnegativity, check corners.
\end{commonmistake}

\begin{commonmistake}
In endowment problems, forgetting to classify buyer versus seller is usually fatal. The substitution effect has its usual sign, but the income effect reverses between buyers and sellers.
\end{commonmistake}

\begin{commonmistake}
In intertemporal questions, do not keep the raw interest-rate algebra longer than necessary. Rewrite immediately as
\[
p_t=(1+r)^{-t}
\]
and reduce to a standard budget problem.
\end{commonmistake}

\begin{commonmistake}
In risky-choice problems, do not use expected value when the question specifies expected utility. Higher expected value does \emph{not} imply higher welfare for a risk-averse consumer.
\end{commonmistake}

\subsection{Integrated Worked Review Examples}

\begin{examplebox}[title={Worked Review Example 1: GARP logic from a chain of choices}]
Suppose three bundles are observed:
\[
x,\ y,\ z.
\]
You find that
\[
x R y,
\qquad
y R z,
\qquad
z P x.
\]
Then
\[
x R^* z
\]
because there is a revealed-preference chain from $x$ to $z$. Since also
\[
z P x,
\]
the data violate GARP.

Why? Under increasing utility maximization,
\[
x R^* z \implies U(x)\ge U(z),
\]
but
\[
z P x \implies U(z)>U(x),
\]
which is impossible.

So the correct conclusion is:
\[
\boxed{\text{The dataset cannot be rationalized by an increasing utility function.}}
\]
\end{examplebox}

\begin{examplebox}[title={Worked Review Example 2: Slutsky decomposition and the law of demand}]
Suppose a two-good consumer has demand
\[
x(p_1,p_2,m)=\left(\frac{m}{2p_1},\frac{m}{2p_2}\right).
\]
This is the Cobb--Douglas demand from $U(x_1,x_2)=x_1x_2$.

Let
\[
p=(1,1), \qquad m=8,
\qquad
p'=(2,1).
\]
Original demand:
\[
x(p,m)=\left(\frac{8}{2},\frac{8}{2}\right)=(4,4).
\]
New uncompensated demand:
\[
x(p',m)=\left(\frac{8}{4},\frac{8}{2}\right)=(2,4).
\]
So the total change in demand for good 1 is
\[
\Delta x_1 = 2-4=-2.
\]

Now compute Slutsky compensation:
\[
m^s = p' \cdot x(p,m) = 2\cdot 4 + 1\cdot 4 = 12.
\]
Compensated demand:
\[
x(p',m^s)=\left(\frac{12}{4},\frac{12}{2}\right)=(3,6).
\]
Hence
\[
\Delta x_1^s = 3-4=-1,
\qquad
\Delta x_1^n = 2-3=-1.
\]
Therefore
\[
\Delta x_1=\Delta x_1^s+\Delta x_1^n=-1+(-1)=-2.
\]

Interpretation:
\begin{itemize}
    \item substitution effect is negative;
    \item income effect is negative because the good is normal;
    \item law of demand holds.
\end{itemize}
\end{examplebox}

\begin{examplebox}[title={Worked Review Example 3: endowment buyer versus seller}]
Suppose the same demand
\[
x(p_1,p_2,m)=\left(\frac{m}{2p_1},\frac{m}{2p_2}\right)
\]
applies, and the endowment is
\[
\omega=(1,6).
\]
At prices
\[
p=(1,1),
\]
wealth is
\[
m=p\cdot \omega=7.
\]
So demand is
\[
x(1,1,7)=\left(\frac72,\frac72\right).
\]
Hence net demand is
\[
z = x-\omega = \left(\frac72-1,\frac72-6\right)=\left(\frac52,-\frac52\right).
\]
So the consumer is a \emph{buyer of good 1}.

If $p_1$ rises, then:
\begin{itemize}
    \item substitution effect on good 1 is negative;
    \item because the consumer is a buyer of good 1, the income effect is also negative if demand is normal.
\end{itemize}
Therefore
\[
\boxed{\Delta x_1 \le 0.}
\]

The exam lesson is that the buyer/seller classification determines the sign of the income effect in endowment problems.
\end{examplebox}

\begin{examplebox}[title={Worked Review Example 4: intertemporal demand and discounting}]
Suppose the income stream is
\[
y=(190,11),
\qquad
r=0.1,
\]
and utility is
\[
U(x_0,x_1)=x_0x_1.
\]
Present-value wealth is
\[
m = 190 + \frac{11}{1.1}=200.
\]
For this utility, intertemporal demand is
\[
x(r,m)=\left(\frac{m}{2},\frac{(1+r)m}{2}\right).
\]
So
\[
x(0.1,200)=\left(100,\frac{1.1\cdot 200}{2}\right)=(100,110).
\]
Hence the consumer saves
\[
190-100=90
\]
in period 0.

The revision point is that once the present value is computed, the problem is just a standard two-good demand problem with prices
\[
\left(1,\frac{1}{1+r}\right).
\]
\end{examplebox}

\begin{examplebox}[title={Worked Review Example 5: state prices and expected-utility demand}]
Suppose there are two states with probabilities
\[
\pi=\left(\frac13,\frac23\right),
\]
and Bernoulli utility
\[
u(x)=\ln x.
\]
For log utility under expected utility,
\[
x(p_1,p_2,m)=\left(\frac{\pi_1 m}{p_1},\frac{\pi_2 m}{p_2}\right).
\]

First suppose the implied state prices are
\[
p=(3,1)
\]
and wealth is
\[
m=30.
\]
Then
\[
x(3,1,30)=\left(\frac{(1/3)\cdot 30}{3},\frac{(2/3)\cdot 30}{1}\right)=\left(\frac{10}{3},20\right).
\]

Now suppose instead two assets generate state prices
\[
p=(6,1)
\]
and the consumer has wealth
\[
m=18.
\]
Then demand is
\[
x(6,1,18)=\left(\frac{(1/3)\cdot 18}{6},\frac{(2/3)\cdot 18}{1}\right)=(1,12).
\]

These examples show the standard review workflow:
\begin{enumerate}
    \item recover the state prices from insurance or assets;
    \item plug them into the expected-utility demand formula.
\end{enumerate}
\end{examplebox}

\subsection{Part I Revision Checklist}

\begin{examtip}[title={Part I final checklist}]
Before the test, make sure you can do all of the following quickly:
\begin{enumerate}
    \item define $R$, $P$, $R^*$, WARP, and GARP precisely;
    \item classify goods as ordinary/Giffen, normal/inferior, substitutes/complements;
    \item derive and use the Slutsky identity;
    \item explain the difference between Slutsky and Hicksian compensation;
    \item compute $MU_1$, $MU_2$, and apply
    \[
    \frac{MU_1}{MU_2}=\frac{p_1}{p_2}
    \]
    when appropriate;
    \item form net demand
    \[
    z=x-\omega
    \]
    and classify buyer versus seller;
    \item write any intertemporal budget in present-value form;
    \item apply
    \[
    u'(x_0)=(1+r)\delta u'(x_1)
    \]
    in discounted-utility problems;
    \item price assets by replication or state prices;
    \item compute expected value, expected utility, and identify risk aversion from concavity;
    \item test first-order stochastic dominance.
\end{enumerate}
\end{examtip}

\begin{keyformula}[title={Key Formula: Part I memory sheet}]
\begin{align*}
x^t R x &\iff p^t\cdot x^t \ge p^t\cdot x, \\
x^t P x &\iff p^t\cdot x^t > p^t\cdot x, \\
\Delta x_1 &= \Delta x_1^s + \Delta x_1^n, \\
z &= x-\omega, \\
p\cdot z &= 0 \quad \text{at an optimum with increasing utility}, \\
\sum_{t=0}^{T}\frac{x_t}{(1+r)^t} &= \sum_{t=0}^{T}\frac{y_t}{(1+r)^t}, \\
u'(x_0) &= (1+r)\delta u'(x_1), \\
U(x) &= \sum_{s=1}^S \pi_s u(x_s), \\
\frac{\pi_i u'(x_i)}{\pi_j u'(x_j)} &= \frac{p_i}{p_j}, \\
u(E[x]) &\ge E[u(x)] \quad \text{if } u \text{ is concave}.
\end{align*}
\end{keyformula}
\newpage
% chapters/ch07_measuring_welfare.tex

\section{Measuring Welfare}

This chapter begins Part II by introducing monetary measures of welfare. Utility is the conceptual object that captures welfare, but utility itself is not directly observable. For policy analysis, market intervention, and social-welfare comparisons, it is therefore useful to translate changes in utility into dollar terms. The main measures in this chapter are reservation prices, gains to trade, consumer surplus, compensating variation, equivalent variation, and producer surplus.

\subsection{Why Measure Welfare in Monetary Terms?}

A central goal of welfare economics is to compare alternative allocations, prices, and policies. Since utility is not directly observed, economists often use monetary measures that track how much income a consumer would be willing to give up or require in compensation.

\begin{definition}{Monetary welfare measure}{ch7-monetary-welfare}
A \emph{monetary welfare measure} is a dollar amount that is equivalent to a change in utility. It answers questions such as:
\begin{itemize}
    \item How much is a consumer willing to pay for a good or opportunity?
    \item How much compensation is needed after a harmful price change?
    \item How much welfare is gained from access to a market?
\end{itemize}
\end{definition}

\begin{examtip}
The point of monetary welfare measures is not to replace utility, but to express changes in utility in a common unit that can be compared with costs, benefits, taxes, subsidies, or producer gains.
\end{examtip}

\subsection{Reservation Prices and Gains to Trade}

The most direct monetary measure of value is a reservation price.

\begin{definition}{Reservation price}{ch7-reservation-price}
A \emph{reservation price} is the maximum amount of money a consumer is willing to pay for a good, action, or marginal increase in consumption.
\end{definition}

For an indivisible object, such as buying one iPhone, the reservation price is defined by indifference between buying and not buying.

\begin{keyformula}[title={Key Formula: Reservation Price for a Discrete Object}]
Suppose utility depends on cash holdings and ownership of an object:
\[
U(\text{cash}, \text{ownership}).
\]
If income is $I$, then the reservation price $r$ satisfies
\[
U(I-r,\text{own object}) = U(I,\text{do not own object}).
\]
Hence:
\begin{itemize}
    \item if price $p<r$, buy;
    \item if price $p>r$, do not buy.
\end{itemize}
\end{keyformula}

For repeated units of a good, reservation prices can be defined unit by unit.

\begin{definition}{Reservation price sequence and gains to trade}{ch7-rp-sequence}
Let $r_n$ denote the reservation price for the $n$th unit of a good, given that the consumer already has $n-1$ units. If the market price is $p$, then the consumer buys the $n$th unit only if
\[
r_n - p > 0.
\]
If the consumer buys $n$ units, the total dollar value of net utility gains from trade is
\[
\sum_{i=1}^{n} (r_i-p)
=
\left(\sum_{i=1}^{n} r_i\right) - np.
\]
\end{definition}

When the good is infinitely divisible, this becomes the area between the reservation price curve and the price line.

\begin{examplebox}[title={Worked Example: reservation price for nuggets}]
Suppose Tom has utility
\[
U(\text{cash held}, \#\text{nuggets}) = \text{cash held} + 10 \times \#\text{nuggets}.
\]

Then the gain from one extra nugget is always $10$ units of utility, and one dollar of cash is worth $1$ unit of utility. Therefore the reservation price for the first nugget is
\[
\boxed{10}.
\]

Because utility is linear in the number of nuggets, the marginal utility of nuggets does not change with the number already consumed. Hence the reservation price does \emph{not} change across units:
\[
r_1=r_2=r_3=\cdots=10.
\]
\end{examplebox}

\subsection{Consumer Surplus and Its Interpretation}

Consumer surplus is often used to approximate the gains to trade from market participation.

\begin{definition}{Consumer surplus}{ch7-cs}
Let $x(p)$ denote Marshallian demand for a good and let $P(x)$ denote inverse demand. Consumer surplus at price $p$ is
\[
CS = \int_0^{x(p)} P(q)\,dq - p\,x(p).
\]
Graphically, it is the area under the inverse demand curve and above the price line, up to the chosen quantity.
\end{definition}

For a price increase from $p'$ to $p''>p'$, the loss in consumer surplus is the area under Marshallian demand between the two prices:
\[
\Delta CS_{\text{loss}} = \int_{p'}^{p''} x(p)\,dp.
\]

\begin{commonmistake}
Consumer surplus is \emph{not} generally the same as the exact utility gain from trade. It is computed from the inverse demand curve, whereas exact gains to trade are based on the reservation price curve.
\end{commonmistake}

\subsection{Reservation Price Curve versus Inverse Demand Curve}

These two curves look similar, but they are conceptually different.

\begin{definition}{Reservation price curve}{ch7-rp-curve}
A \emph{reservation price curve} records the willingness to pay for successive marginal units of a good, holding the reference level of money available for other goods fixed.
\end{definition}

\begin{definition}{Inverse demand curve}{ch7-inverse-demand}
An \emph{inverse demand curve} gives the price at which a total quantity $x$ is the optimal quantity purchased when all $x$ units are bought at that price.
\end{definition}

The difference comes from income effects:
\begin{itemize}
    \item along the reservation price curve, the reference cash level is held fixed when valuing marginal units;
    \item along inverse demand, the amount of money left over changes with the quantity chosen.
\end{itemize}
Therefore the two curves need not coincide.

\begin{examtip}
If the question asks whether consumer surplus is an \emph{exact} welfare measure, the correct answer is usually ``not in general.'' It becomes exact only under special assumptions, especially quasi-linear utility.
\end{examtip}

\subsection{Quasi-Linear Utility and Exact Consumer Surplus}

The main special case where consumer surplus becomes exact is quasi-linearity.

\begin{definition}{Quasi-linear utility}{ch7-quasilinear}
A two-good utility function is \emph{quasi-linear in $x_2$} if it takes the form
\[
U(x_1,x_2)=v(x_1)+\alpha x_2,
\]
where $v$ is increasing and typically concave, and $\alpha>0$ is constant.
\end{definition}

If utility is quasi-linear in $x_2$, then the marginal utility of $x_2$ is constant. This removes the income effect on $x_1$.

\begin{theorem}{No income effect under quasi-linear utility}{ch7-no-income-effect}
If
\[
U(x,m)=v(x)+m,
\]
where $m$ is money left over for other goods, then demand for $x$ is independent of income. In particular, the inverse demand curve and reservation price curve coincide.
\end{theorem}

\begin{proof}
Consider
\[
\max_{x,m} \ v(x)+m
\quad \text{subject to} \quad
px+m=I.
\]
Substituting $m=I-px$ gives
\[
\max_x \ v(x)+I-px.
\]
The first-order condition is
\[
v'(x)=p.
\]
This condition depends on $x$ and $p$, but not on income $I$. Hence there is no income effect on $x$. Since willingness to pay for the marginal unit is also given by
\[
v'(x),
\]
the inverse demand and reservation price curves coincide.
\end{proof}

\begin{theorem}{Consumer surplus equals gains to trade under quasi-linear utility}{ch7-cs-exact}
If
\[
U(x,m)=v(x)+m,
\]
then consumer surplus is an exact measure of gains to trade:
\[
CS = v(x)-v(0)-px.
\]
\end{theorem}

\begin{proof}
Under quasi-linear utility, inverse demand is
\[
p=v'(x).
\]
Hence
\[
CS = \int_0^x v'(q)\,dq - px
= v(x)-v(0)-px.
\]
This is exactly the net utility gain from obtaining $x$ units rather than zero units.
\end{proof}

\begin{examplebox}[title={Worked Example: identifying quasi-linear utility}]
Consider the following utility functions:
\[
U(x,y)=x+y, \qquad U(x,y)=x+4y, \qquad U(x,y)=x+3\sqrt{y}.
\]

\begin{itemize}
    \item $U(x,y)=x+y$ is quasi-linear in both goods, since it is linear in both.
    \item $U(x,y)=x+4y$ is also quasi-linear in both goods.
    \item $U(x,y)=x+3\sqrt{y}$ is quasi-linear in $x$, since it can be written as
    \[
    U(x,y)=x+v(y)
    \quad \text{with} \quad
    v(y)=3\sqrt{y}.
    \]
\end{itemize}

However, the zero-income-effect conclusion must be attached to the good that enters through the \emph{nonlinear} part in the standard quasi-linear setup. For example, if
\[
U(x,y)=y+3\sqrt{x},
\]
then utility is quasi-linear in $y$, and demand for $x$ has zero income effect.
\end{examplebox}

\begin{commonmistake}
Do not say ``quasi-linear means no income effects for both goods.'' The standard result is that the good appearing inside the nonlinear part has zero income effect when utility is quasi-linear in the other good.
\end{commonmistake}

\subsection{Price Changes and Welfare: Indirect Utility, CV, and EV}

Consumer surplus is convenient, but it is not always exact. Two exact monetary measures of welfare change are compensating variation and equivalent variation.

\begin{definition}{Indirect utility}{ch7-indirect-utility}
The \emph{indirect utility function} is
\[
V(p,I)=\max_{x} U(x)
\quad \text{subject to} \quad
p\cdot x \le I.
\]
It gives the maximum utility attainable at prices $p$ and income $I$.
\end{definition}

Suppose the price of good 1 rises from $p_1'$ to $p_1''$, with $p_2$ fixed and income $I$ fixed.

\begin{definition}{Compensating variation and equivalent variation}{ch7-cv-ev}
For a price increase:
\begin{itemize}
    \item \emph{Compensating variation (CV)} is the least extra income needed at the \emph{new} prices to restore the consumer to the \emph{original} utility level:
    \[
    V(p_1'',p_2, I+CV)=V(p_1',p_2,I).
    \]
    \item \emph{Equivalent variation (EV)} is the greatest amount of income that can be taken away at the \emph{old} prices while leaving the consumer at the \emph{final} utility level:
    \[
    V(p_1',p_2, I-EV)=V(p_1'',p_2,I).
    \]
\end{itemize}
\end{definition}

For a price increase, both $CV$ and $EV$ are positive if they are interpreted as magnitudes of welfare loss.

\begin{examtip}
The difference between CV and EV is the reference point:
\begin{itemize}
    \item \textbf{CV:} new prices, old utility;
    \item \textbf{EV:} old prices, new utility.
\end{itemize}
A good mnemonic is:
\begin{itemize}
    \item \emph{compensating} = money added after the bad change;
    \item \emph{equivalent} = money taken away before the bad change.
\end{itemize}
\end{examtip}

\begin{examplebox}[title={Worked Example: indirect utility, EV, and CV for perfect complements}]
Suppose
\[
U(x_1,x_2)=\min\{x_1,x_2\},
\qquad
I=100,
\qquad
p_2=10,
\]
and the price of good 1 rises from
\[
p_1'=5
\quad \text{to} \quad
p_1''=10.
\]

Since utility is Leontief, the optimal bundle satisfies
\[
x_1=x_2.
\]
Using the budget constraint
\[
p_1x_1+p_2x_2=I,
\]
we get
\[
x_1=x_2=\frac{I}{p_1+p_2}.
\]
Therefore the indirect utility is
\[
V(p_1,p_2,I)=\frac{I}{p_1+p_2}.
\]

\textbf{Initial utility:}
\[
V(5,10,100)=\frac{100}{15}=\frac{20}{3}.
\]

\textbf{Final utility:}
\[
V(10,10,100)=\frac{100}{20}=5.
\]

\textbf{Equivalent variation:} solve
\[
V(5,10,100-EV)=5.
\]
So
\[
\frac{100-EV}{15}=5
\quad \Longrightarrow \quad
100-EV=75
\quad \Longrightarrow \quad
EV=25.
\]

\textbf{Compensating variation:} solve
\[
V(10,10,100+CV)=\frac{20}{3}.
\]
So
\[
\frac{100+CV}{20}=\frac{20}{3}
\quad \Longrightarrow \quad
100+CV=\frac{400}{3}
\quad \Longrightarrow \quad
CV=\frac{100}{3}.
\]

Hence
\[
\boxed{EV=25, \qquad CV=\frac{100}{3}.}
\]
Both are positive, as expected for a welfare-reducing price increase.
\end{examplebox}

\subsection{Comparing \texorpdfstring{$\Delta CS$}{Delta CS}, CV, and EV}

These three measures are closely related.

\begin{theorem}{Relationship between consumer surplus, CV, and EV}{ch7-cs-cv-ev}
For a price change in the two-good case, the consumer-surplus measure lies between compensating variation and equivalent variation:
\[
\min\{|CV|,|EV|\} \le |\Delta CS| \le \max\{|CV|,|EV|\}.
\]
When utility is quasi-linear in the other good, all three measures coincide.
\end{theorem}

\begin{examplebox}[title={Worked Example: comparing the three welfare measures}]
Continue the previous example with
\[
U(x_1,x_2)=\min\{x_1,x_2\},
\qquad
I=100,
\qquad
p_2=10,
\qquad
p_1: 5 \to 10.
\]
We already found
\[
EV=25,
\qquad
CV=\frac{100}{3}\approx 33.33.
\]

Marshallian demand for good 1 is
\[
x_1(p_1)=\frac{100}{p_1+10}.
\]
So the loss in consumer surplus from the price increase is
\[
\Delta CS_{\text{loss}}
=
\int_{5}^{10}\frac{100}{p_1+10}\,dp_1
=
100\ln\left(\frac{20}{15}\right)
=
100\ln\left(\frac{4}{3}\right)
\approx 28.77.
\]

Hence
\[
25 = EV < \Delta CS_{\text{loss}} \approx 28.77 < CV \approx 33.33.
\]
So $\Delta CS$ lies between EV and CV, exactly as theory predicts.
\end{examplebox}

\begin{commonmistake}
Do not treat
\[
x_1 \cdot (p_1''-p_1')
\]
as the welfare loss from a price increase. That ignores substitution effects and therefore generally overstates the true utility loss.
\end{commonmistake}

\subsection{Producer Surplus}

Welfare analysis also applies to producers. The lecture treats producer surplus as the producer-side analogue of consumer surplus.

\begin{definition}{Producer surplus}{ch7-ps}
If a firm supplies quantity $q$ at price $p$ and has marginal cost curve $MC(q)$, then producer surplus is
\[
PS = pq - VC(q)
= \int_0^q \bigl(p-MC(\tilde q)\bigr)\,d\tilde q,
\]
where $VC(q)$ is variable cost.
\end{definition}

Under standard profit maximization, the competitive supply curve is the marginal-cost curve above shutdown. Hence producer surplus is the area above supply and below price, up to the quantity sold.

\begin{examplebox}[title={Worked Example: producer surplus from marginal cost}]
Suppose a competitive firm's marginal cost is
\[
MC(q)=2q,
\]
and market price is
\[
p=10.
\]
The profit-maximizing quantity satisfies
\[
p=MC(q)
\quad \Longrightarrow \quad
10=2q
\quad \Longrightarrow \quad
q=5.
\]

Producer surplus is
\[
PS=\int_0^5 (10-2q)\,dq
= \left[10q-q^2\right]_0^5
=50-25
=25.
\]
Therefore
\[
\boxed{PS=25.}
\]
\end{examplebox}

\begin{policynote}
Consumer surplus and producer surplus are central because later welfare analysis often aggregates them to study taxes, subsidies, quotas, and other market interventions. This chapter provides the monetary foundations for that aggregation.
\end{policynote}

\subsection{Standard Problem Types and Solution Patterns}

\begin{examtip}[title={Method Pattern: reservation price questions}]
For discrete or indivisible objects:
\begin{enumerate}
    \item write the consumer's utility before and after purchase;
    \item set them equal to identify the reservation price;
    \item compare the actual price with the reservation price.
\end{enumerate}
\end{examtip}

\begin{examtip}[title={Method Pattern: quasi-linear utility}]
If utility is
\[
U(x,m)=v(x)+m,
\]
then:
\begin{enumerate}
    \item use
    \[
    v'(x)=p
    \]
    to find demand;
    \item conclude there is no income effect on $x$;
    \item use consumer surplus as an exact welfare measure.
\end{enumerate}
\end{examtip}

\begin{examtip}[title={Method Pattern: EV/CV questions}]
For price-change questions:
\begin{enumerate}
    \item derive or identify the indirect utility function $V(p,I)$;
    \item compute initial and final utility levels;
    \item solve
    \[
    V(\text{new prices}, I+CV)=V(\text{old prices}, I)
    \]
    for $CV$;
    \item solve
    \[
    V(\text{old prices}, I-EV)=V(\text{new prices}, I)
    \]
    for $EV$.
\end{enumerate}
\end{examtip}

\subsection{Chapter Summary}

\begin{keyformula}[title={Key Formula: Lecture 7 Summary}]
\begin{align*}
U(I-r,\text{own}) &= U(I,\text{not own}) && \text{reservation price}, \\
\text{Gains to trade} &= \sum_{i=1}^{n}(r_i-p), \\
CS &= \int_0^{x(p)} P(q)\,dq - p\,x(p), \\
V(p,I) &= \max_x U(x) \ \text{s.t.}\ p\cdot x \le I, \\
V(p_1'',p_2,I+CV) &= V(p_1',p_2,I), \\
V(p_1',p_2,I-EV) &= V(p_1'',p_2,I), \\
\min\{|CV|,|EV|\} &\le |\Delta CS| \le \max\{|CV|,|EV|\}, \\
U(x,m)=v(x)+m &\Rightarrow v'(x)=p \text{ and } CS \text{ is exact}, \\
PS &= \int_0^q \bigl(p-MC(\tilde q)\bigr)\,d\tilde q.
\end{align*}
\end{keyformula}

\begin{examtip}
The highest-yield takeaways from this chapter are:
\begin{itemize}
    \item utility changes are often measured in money because utility is not directly observable;
    \item reservation prices measure willingness to pay for marginal utility gains;
    \item consumer surplus usually approximates gains to trade, but becomes exact under quasi-linear utility;
    \item EV and CV are exact welfare measures for price changes and account for substitution effects;
    \item for a price increase, EV and CV are positive magnitudes of welfare loss;
    \item producer surplus measures firm welfare when firms maximize profit.
\end{itemize}
\end{examtip}
\newpage
% chapters/ch08_simple_welfare_analysis.tex

\section{Simple Welfare Analysis}

This chapter moves from individual welfare to aggregate welfare. Chapter 7 introduced monetary measures of consumer and producer welfare. Here, those measures are summed across agents to study welfare at the market level. The main applications are welfare analysis in demand--supply diagrams, tax regimes, deadweight loss, tax incidence, and comparative statics.

\subsection{Aggregating Welfare Across Consumers and Producers}

Consider a single-good market. Let individual Marshallian demands be
\[
q_i^D(p),
\]
and individual supplies be
\[
q_j^S(p).
\]
The market curves are obtained by horizontal summation.

\begin{definition}{Aggregate demand and aggregate supply}{ch8-agg-curves}
The \emph{aggregate demand} and \emph{aggregate supply} functions are
\[
Q^D(p)=\sum_i q_i^D(p),
\qquad
Q^S(p)=\sum_j q_j^S(p).
\]
Their inverse forms are denoted by
\[
P_D(Q), \qquad P_S(Q),
\]
whenever these are well-defined.
\end{definition}

Aggregate welfare is measured by adding up individual surpluses.

\begin{definition}{Aggregate consumer surplus, producer surplus, and social welfare}{ch8-sw-def}
If market quantity is $Q$ and consumers pay price $p$, aggregate consumer surplus is
\[
CS(Q)=\int_0^Q P_D(q)\,dq - pQ.
\]
If producers receive price $p$, aggregate producer surplus is
\[
PS(Q)=pQ-\int_0^Q P_S(q)\,dq.
\]
Hence \emph{social welfare} in the single-good market is
\[
SW(Q)=CS(Q)+PS(Q)
=
\int_0^Q \bigl(P_D(q)-P_S(q)\bigr)\,dq.
\]
\end{definition}

The last expression is the key one: social welfare is the area between inverse demand and inverse supply up to the traded quantity.

\begin{theorem}{Competitive equilibrium maximizes $CS+PS$}{ch8-sw-max}
Suppose inverse demand is downward sloping and inverse supply is upward sloping. Then social welfare
\[
SW(Q)=\int_0^Q \bigl(P_D(q)-P_S(q)\bigr)\,dq
\]
is maximized at the competitive quantity $Q^*$ satisfying
\[
P_D(Q^*)=P_S(Q^*).
\]
\end{theorem}

\begin{proof}
Differentiate:
\[
\frac{dSW(Q)}{dQ}=P_D(Q)-P_S(Q).
\]
Thus welfare rises whenever
\[
P_D(Q)>P_S(Q),
\]
falls whenever
\[
P_D(Q)<P_S(Q),
\]
and is maximized when
\[
P_D(Q)=P_S(Q).
\]
Under the usual slopes, this point is unique and coincides with the competitive equilibrium quantity.
\end{proof}

\begin{examtip}
The fastest welfare argument in a competitive single-good market is:
\[
SW(Q)=\int_0^Q \bigl(P_D(q)-P_S(q)\bigr)\,dq.
\]
Hence any deviation from the competitive quantity loses the positive gains from trade on units for which demand exceeds supply, or forces inefficient units for which supply exceeds demand.
\end{examtip}

\begin{examplebox}[title={Worked Example: aggregate demand and aggregate consumer surplus}]
Suppose there are 30 identical buyers, each with inverse demand
\[
p=100-2q.
\]
Equivalently, each individual's demand is
\[
q_i(p)=\frac{100-p}{2}
\qquad \text{for } p\le 100,
\]
and $q_i(p)=0$ for $p>100$.

Therefore aggregate demand is
\[
Q^D(p)=30\cdot \frac{100-p}{2}=1500-15p
\qquad \text{for } p\le 100.
\]
In inverse form,
\[
p=100-\frac{Q}{15}.
\]

Now suppose the market price is
\[
p=10.
\]
Then aggregate quantity demanded is
\[
Q=1500-15(10)=1350.
\]
Aggregate consumer surplus is
\[
CS=\frac12 \times (100-10)\times 1350
=
\frac12 \times 90 \times 1350
=
60750.
\]

So the correct answers are
\[
\boxed{Q^D(p)=1500-15p}
\quad \text{and} \quad
\boxed{CS=60750.}
\]
\end{examplebox}

\subsection{Deadweight Welfare Loss}

Because the competitive equilibrium maximizes $CS+PS$, any policy that changes the traded quantity away from $Q^*$ creates a welfare loss.

\begin{definition}{Deadweight welfare loss}{ch8-dwl-def}
If a policy reduces market quantity from the efficient level $Q^*$ to some lower quantity $Q'$, the \emph{deadweight welfare loss} is
\[
DWL
=
\int_{Q'}^{Q^*} \bigl(P_D(q)-P_S(q)\bigr)\,dq.
\]
\end{definition}

This is the value of mutually beneficial trades that no longer occur.

\begin{keyformula}[title={Key Formula: Deadweight Loss under a Per-Unit Tax}]
If a per-unit tax $t$ creates a constant wedge between buyer and seller prices and the market quantity falls from $Q^*$ to $Q'$, then graphically
\[
DWL=\frac12\, t \,(Q^*-Q')
\]
whenever the relevant region is triangular, in particular under linear demand and supply.
\end{keyformula}

\begin{theorem}{No quantity distortion implies no deadweight loss}{ch8-no-dwl}
If a tax does not change the equilibrium quantity, then it does not create deadweight loss.
\end{theorem}

\begin{proof}
By definition,
\[
DWL=\int_{Q'}^{Q^*} \bigl(P_D(q)-P_S(q)\bigr)\,dq.
\]
If $Q'=Q^*$, the interval of integration is empty, so
\[
DWL=0.
\]
\end{proof}

\begin{examtip}
A standard implication is: if either demand or supply is perfectly inelastic, then a tax can be pure transfer with no deadweight loss because the traded quantity does not change.
\end{examtip}

\subsection{Quantity Taxes}

A quantity tax, or per-unit tax, imposes a fixed dollar tax $t$ on each unit traded.

\begin{definition}{Quantity taxes, excise taxes, and sales taxes}{ch8-quantity-tax}
A \emph{quantity tax} is a fixed amount $t>0$ per unit traded.
\begin{itemize}
    \item If sellers pay it, it is an \emph{excise tax}.
    \item If buyers pay it, it is a \emph{sales tax}.
\end{itemize}
Let
\[
p_b = \text{price paid by buyers},
\qquad
p_s = \text{price received by sellers}.
\]
Then in either case,
\[
p_b=p_s+t.
\]
\end{definition}

The statutory side of the tax does not matter for the equilibrium outcome when the tax is a fixed amount per unit. What matters is the wedge condition
\[
p_b-p_s=t.
\]

\begin{theorem}{Quantity excise and quantity sales taxes are equivalent}{ch8-quantity-equivalence}
A per-unit excise tax of size $t$ and a per-unit sales tax of size $t$ generate the same equilibrium allocation and the same buyer and seller prices.
\end{theorem}

\begin{proof}
In both cases equilibrium requires
\[
Q^D(p_b)=Q^S(p_s)
\]
together with
\[
p_b=p_s+t.
\]
These two equations are identical across the two statutory tax regimes, so the equilibrium values of $p_b$, $p_s$, and $Q$ are the same.
\end{proof}

\subsubsection{Graphical effect of a quantity excise tax}

Let the original supply curve be
\[
Q=S(p).
\]
If the tax is levied on sellers, then when buyers pay price $p$, sellers receive only $p-t$. Therefore supply becomes
\[
Q=S(p-t).
\]
In inverse form, if the original inverse supply is
\[
p=P_S(Q),
\]
the post-tax inverse supply is
\[
p=P_S(Q)+t.
\]
So a quantity excise tax shifts inverse supply vertically upward by $t$.

\subsubsection{Graphical effect of a quantity sales tax}

Let the original demand curve be
\[
Q=D(p).
\]
If the tax is levied on buyers, then when sellers receive price $p$, buyers pay $p+t$. Therefore demand becomes
\[
Q=D(p+t).
\]
In inverse form, if the original inverse demand is
\[
p=P_D(Q),
\]
the post-tax inverse demand is
\[
p=P_D(Q)-t.
\]
So a quantity sales tax shifts inverse demand vertically downward by $t$.

\begin{keyformula}[title={Key Formula: Welfare Accounting under a Quantity Tax}]
With equilibrium quantity $Q'$, buyer price $p_b$, seller price $p_s$, and original demand and supply curves:
\[
CS^{tax}=\int_0^{Q'} \bigl(P_D(q)-p_b\bigr)\,dq,
\]
\[
PS^{tax}=\int_0^{Q'} \bigl(p_s-P_S(q)\bigr)\,dq,
\]
\[
TR=tQ',
\]
\[
DWL=\int_{Q'}^{Q^*} \bigl(P_D(q)-P_S(q)\bigr)\,dq.
\]
\end{keyformula}

\subsection{Ad-Valorem Taxes}

An ad-valorem tax is proportional to value rather than fixed per unit.

\begin{definition}{Ad-valorem tax}{ch8-ad-valorem}
An \emph{ad-valorem tax} of rate $\tau$ satisfies $0<\tau<1$ and is proportional to price.
\begin{itemize}
    \item Under an ad-valorem excise tax on sellers,
    \[
    p_s=(1-\tau)p_b.
    \]
    \item Under an ad-valorem sales tax on buyers,
    \[
    p_b=(1+\tau)p_s.
    \]
\end{itemize}
\end{definition}

Unlike fixed quantity taxes, equal percentage excise and sales taxes are not exactly the same tax wedge, though they are approximately similar for small $\tau$.

\subsubsection{Ad-valorem excise tax}

If original supply is
\[
Q=S(p_s),
\]
and buyers pay price $p_b$, sellers receive
\[
p_s=(1-\tau)p_b.
\]
Hence the taxed supply curve is
\[
Q=S((1-\tau)p_b).
\]
In inverse form,
\[
p_b=\frac{P_S(Q)}{1-\tau}.
\]
So the inverse supply curve pivots upward.

\subsubsection{Ad-valorem sales tax}

If original demand is
\[
Q=D(p_b),
\]
and sellers receive price $p_s$, buyers pay
\[
p_b=(1+\tau)p_s.
\]
Thus demand becomes
\[
Q=D((1+\tau)p_s),
\]
or in inverse form,
\[
p_s=\frac{P_D(Q)}{1+\tau}.
\]
So the inverse demand curve pivots downward.

\begin{commonmistake}
For ad-valorem taxes, the pivoted curve is useful for locating the new equilibrium, but it is \emph{not} the correct object for measuring producer surplus or deadweight loss. Once equilibrium is found, welfare must be measured using the original supply and demand curves together with the actual buyer and seller prices.
\end{commonmistake}

\begin{keyformula}[title={Key Formula: Welfare Accounting under an Ad-Valorem Tax}]
If the post-tax equilibrium is $(Q',p_b,p_s)$, then:
\[
CS^{tax}=\int_0^{Q'} \bigl(P_D(q)-p_b\bigr)\,dq,
\]
\[
PS^{tax}=\int_0^{Q'} \bigl(p_s-P_S(q)\bigr)\,dq,
\]
\[
TR=(p_b-p_s)Q',
\]
\[
DWL=\int_{Q'}^{Q^*} \bigl(P_D(q)-P_S(q)\bigr)\,dq.
\]
\end{keyformula}

\subsection{Tax Incidence}

A tax lowers welfare for both consumers and producers. The division of the burden is called tax incidence.

\begin{definition}{Tax incidence}{ch8-incidence}
For a unit tax with pre-tax equilibrium price $p^*$ and post-tax prices $(p_b,p_s)$ at traded quantity $Q'$, the tax burden borne by:
\begin{itemize}
    \item consumers is measured by
    \[
    (p_b-p^*)Q';
    \]
    \item producers is measured by
    \[
    (p^*-p_s)Q'.
    \]
\end{itemize}
Their sum is the total tax transfer on traded units:
\[
\bigl((p_b-p^*)+(p^*-p_s)\bigr)Q'
=
(p_b-p_s)Q'
=
tQ'
\]
under a quantity tax.
\end{definition}

\begin{theorem}{More elastic side bears less tax incidence}{ch8-elasticity-incidence}
All else equal:
\begin{itemize}
    \item more elastic demand implies lower consumer tax incidence;
    \item more elastic supply implies lower producer tax incidence.
\end{itemize}
In the limiting cases:
\begin{itemize}
    \item perfectly inelastic demand means consumers bear all of the tax;
    \item perfectly inelastic supply means producers bear all of the tax.
\end{itemize}
\end{theorem}

\begin{proof}
The less elastic side has fewer quantity adjustments available, so its price changes more under a tax wedge. Geometrically, that side experiences the larger movement between pre-tax and post-tax price. The formal linear-demand proof appears in the next subsection.
\end{proof}

\subsection{Comparative Statics with Linear Demand and Supply}

The lecture gives a simple algebraic comparative-statics framework.

\begin{definition}{Linear market environment}{ch8-linear-env}
Suppose market demand and supply are
\[
Q^D(p_b)=a-bp_b,
\qquad
Q^S(p_s)=c+dp_s,
\]
with
\[
b>0, \qquad d>0.
\]
Here $b$ and $d$ capture the price responsiveness of demand and supply in this linear specification.
\end{definition}

\subsubsection{No tax}

Equilibrium satisfies
\[
a-bp^*=c+dp^*.
\]
Hence
\[
p^*=\frac{a-c}{b+d},
\qquad
Q^*=\frac{ad+bc}{b+d}.
\]

\subsubsection{Quantity tax}

Let a per-unit tax $t$ create the wedge
\[
p_b=p_s+t.
\]
Then equilibrium satisfies
\[
a-bp_b=c+dp_s,
\qquad
p_b=p_s+t.
\]

\begin{theorem}{Linear equilibrium under a quantity tax}{ch8-linear-tax}
Under
\[
Q^D(p_b)=a-bp_b,
\qquad
Q^S(p_s)=c+dp_s,
\qquad
p_b=p_s+t,
\]
the post-tax equilibrium is
\[
p_b=\frac{a-c+dt}{b+d},
\qquad
p_s=\frac{a-c-bt}{b+d},
\qquad
Q'=\frac{ad+bc-bdt}{b+d}.
\]
\end{theorem}

\begin{proof}
Substitute
\[
p_b=p_s+t
\]
into demand:
\[
Q^D=a-b(p_s+t)=a-bp_s-bt.
\]
Equilibrium requires
\[
a-bp_s-bt=c+dp_s.
\]
So
\[
a-c-bt=(b+d)p_s,
\]
hence
\[
p_s=\frac{a-c-bt}{b+d}.
\]
Then
\[
p_b=p_s+t
=
\frac{a-c-bt}{b+d}+t
=
\frac{a-c+dt}{b+d}.
\]
Finally,
\[
Q'=c+dp_s
=
c+d\left(\frac{a-c-bt}{b+d}\right)
=
\frac{ad+bc-bdt}{b+d}.
\]
\end{proof}

\begin{theorem}{Tax incidence shares in the linear model}{ch8-linear-incidence}
Under the same linear quantity-tax model,
\[
\frac{p_b-p^*}{t}=\frac{d}{b+d},
\qquad
\frac{p^*-p_s}{t}=\frac{b}{b+d}.
\]
So:
\begin{itemize}
    \item consumers bear share
    \[
    \frac{d}{b+d};
    \]
    \item producers bear share
    \[
    \frac{b}{b+d}.
    \]
\end{itemize}
\end{theorem}

\begin{proof}
From the formulas above,
\[
p_b-p^*
=
\frac{a-c+dt}{b+d}-\frac{a-c}{b+d}
=
\frac{dt}{b+d},
\]
so
\[
\frac{p_b-p^*}{t}=\frac{d}{b+d}.
\]
Similarly,
\[
p^*-p_s
=
\frac{a-c}{b+d}-\frac{a-c-bt}{b+d}
=
\frac{bt}{b+d},
\]
so
\[
\frac{p^*-p_s}{t}=\frac{b}{b+d}.
\]
\end{proof}

\begin{theorem}{Deadweight loss in the linear quantity-tax model}{ch8-linear-dwl}
In the same linear model,
\[
Q^*-Q'
=
\frac{bdt}{b+d}.
\]
Hence
\[
DWL
=
\frac12 t(Q^*-Q')
=
\frac{bd\,t^2}{2(b+d)}.
\]
\end{theorem}

\begin{proof}
Subtract the quantity formulas:
\[
Q^*-Q'
=
\frac{ad+bc}{b+d} - \frac{ad+bc-bdt}{b+d}
=
\frac{bdt}{b+d}.
\]
Under linear demand and supply, the deadweight-loss region is a triangle with base
\[
Q^*-Q'
\]
and height
\[
t.
\]
Therefore
\[
DWL=\frac12 t(Q^*-Q')=\frac{bd\,t^2}{2(b+d)}.
\]
\end{proof}

\begin{examplebox}[title={Worked Example: comparative statics of tax incidence}]
Suppose
\[
Q^D(p_b)=a-bp_b,
\qquad
Q^S(p_s)=c+dp_s,
\]
and there is a unit tax $t$.

The consumer share of incidence is
\[
\frac{d}{b+d}.
\]
Compare the benchmark case
\[
(b,d)=(5,5),
\]
for which the consumer share is
\[
\frac{5}{10}=\frac12.
\]

Now consider the following cases:
\[
(b,d)=(3,3), \quad (3,2), \quad (7,3), \quad (3,7).
\]

Compute the consumer share in each case:
\[
(3,3):\ \frac{3}{6}=\frac12,
\]
\[
(3,2):\ \frac{2}{5}=0.4,
\]
\[
(7,3):\ \frac{3}{10}=0.3,
\]
\[
(3,7):\ \frac{7}{10}=0.7.
\]

Therefore the only case in which consumers bear a strictly higher tax incidence than in the benchmark is
\[
\boxed{(b,d)=(3,7).}
\]
\end{examplebox}

\subsubsection{Behavioral bias and tax collected}

The lecture also studies a simple behavioral-taxation example.

\begin{definition}{Perceived tax parameter}{ch8-k-bias}
Suppose consumers behave as if the buyer price were
\[
p_b=p_s+kt,
\]
where $k>0$.
\begin{itemize}
    \item If $0<k<1$, consumers underweight the tax.
    \item If $k>1$, consumers overweight the tax.
\end{itemize}
\end{definition}

With linear demand and supply,
\[
Q^D=a-bp_b,
\qquad
Q^S=c+dp_s,
\]
equilibrium becomes
\[
a-b(p_s+kt)=c+dp_s.
\]

\begin{theorem}{Equilibrium and tax revenue with perceived-tax bias}{ch8-bias}
Under the perceived-price rule
\[
p_b=p_s+kt,
\]
the equilibrium is
\[
p_s=\frac{a-c-bkt}{b+d},
\qquad
p_b=\frac{a-c+dkt}{b+d},
\qquad
Q=\frac{bc+ad-bdkt}{b+d}.
\]
Tax revenue is
\[
T=tQ
=
\frac{bc+ad-bdkt}{b+d}\,t.
\]
Therefore
\[
\frac{\partial T}{\partial k}<0.
\]
\end{theorem}

\begin{proof}
The equilibrium formulas are obtained exactly as in the ordinary quantity-tax case, with $t$ replaced by $kt$ in the demand wedge.

Since
\[
T(k)=\frac{bc+ad-bdkt}{b+d}\,t,
\]
differentiate with respect to $k$:
\[
\frac{\partial T}{\partial k}
=
-\frac{bdt^2}{b+d}<0.
\]
Thus greater perceived tax salience lowers quantity traded and therefore reduces tax revenue.
\end{proof}

\begin{examplebox}[title={Worked Example: effect of underweighting and overweighting taxes}]
Using the previous formula,
\[
T(k)=\frac{bc+ad-bdkt}{b+d}\,t.
\]
Compare tax revenue under perceived-bias parameter $k$ with the benchmark $k=1$:
\[
\Delta T
=
T(k)-T(1)
=
\frac{bc+ad-bdkt}{b+d}t - \frac{bc+ad-bdt}{b+d}t.
\]
So
\[
\Delta T
=
\frac{bd(1-k)t^2}{b+d}.
\]

Hence:
\begin{itemize}
    \item if $k<1$ (underweighting), then
    \[
    \Delta T>0,
    \]
    so more tax is collected than in the unbiased case;
    \item if $k>1$ (overweighting), then
    \[
    \Delta T<0,
    \]
    so less tax is collected.
\end{itemize}

Therefore the lecture's conclusion is
\[
\boxed{\text{Underweighting taxes raises tax revenue; overweighting taxes lowers it.}}
\]
\end{examplebox}

\subsection{Common Mistakes and Exam Traps}

\begin{commonmistake}
Do not confuse \emph{statutory incidence} with \emph{economic incidence}. The side legally paying the tax need not be the side that bears most of the burden.
\end{commonmistake}

\begin{commonmistake}
For quantity taxes, excise and sales taxes of the same dollar amount are equivalent. For ad-valorem taxes, equal percentage excise and sales taxes are not exactly identical.
\end{commonmistake}

\begin{commonmistake}
Do not use the pivoted supply curve under an ad-valorem tax to calculate producer surplus. It is a device for locating the new equilibrium, not the true marginal cost curve.
\end{commonmistake}

\begin{commonmistake}
When computing deadweight loss, do not forget tax revenue. The welfare decomposition is:
\[
\text{loss in CS}+\text{loss in PS}
=
TR + DWL.
\]
Only the part not transferred as revenue is deadweight loss.
\end{commonmistake}

\subsection{Chapter Summary}

\begin{keyformula}[title={Key Formula: Lecture 8 Summary}]
\begin{align*}
Q^D(p) &= \sum_i q_i^D(p), \\
Q^S(p) &= \sum_j q_j^S(p), \\
SW(Q) &= \int_0^Q \bigl(P_D(q)-P_S(q)\bigr)\,dq, \\
p_b &= p_s+t \qquad \text{(quantity tax)}, \\
p_s &= (1-\tau)p_b \qquad \text{(ad-valorem excise tax)}, \\
p_b &= (1+\tau)p_s \qquad \text{(ad-valorem sales tax)}, \\
CS^{tax} &= \int_0^{Q'} \bigl(P_D(q)-p_b\bigr)\,dq, \\
PS^{tax} &= \int_0^{Q'} \bigl(p_s-P_S(q)\bigr)\,dq, \\
TR &= (p_b-p_s)Q', \\
DWL &= \int_{Q'}^{Q^*} \bigl(P_D(q)-P_S(q)\bigr)\,dq.
\end{align*}
Under
\[
Q^D=a-bp_b, \qquad Q^S=c+dp_s,
\]
a quantity tax gives
\[
p_b=\frac{a-c+dt}{b+d},
\qquad
p_s=\frac{a-c-bt}{b+d},
\qquad
Q'=\frac{ad+bc-bdt}{b+d},
\]
consumer incidence share
\[
\frac{d}{b+d},
\]
producer incidence share
\[
\frac{b}{b+d},
\]
and
\[
DWL=\frac{bdt^2}{2(b+d)}.
\]
\end{keyformula}

\begin{examtip}
The highest-yield takeaways from this chapter are:
\begin{itemize}
    \item aggregate welfare in a single-good market is $CS+PS$;
    \item competitive equilibrium maximizes this welfare under the standard assumptions;
    \item taxes create a wedge, reduce traded quantity, and generate deadweight loss;
    \item the more inelastic side of the market bears more of the tax burden;
    \item quantity excise and sales taxes are equivalent, but ad-valorem taxes require more care;
    \item in tax questions, keep buyer price, seller price, tax revenue, and welfare triangles clearly separated.
\end{itemize}
\end{examtip}
\newpage
% chapters/ch09_perfect_competition_and_welfare.tex

\section{Perfect Competition and Welfare}

This chapter studies welfare in a more general market environment than the single-market partial-equilibrium framework of Chapter 8. Instead of measuring welfare by summing consumer and producer surplus in one market, we consider an exchange economy with multiple goods and multiple consumers, and use Pareto efficiency as the central efficiency concept. The chapter develops the simple exchange economy, the Edgeworth box, competitive equilibrium, the contract curve, envy-freeness, and the two fundamental welfare theorems.

\subsection{The Simple Exchange Economy}

We begin with a pure exchange economy: goods are redistributed across consumers, but there is no production.

\begin{definition}{Simple exchange economy}{ch9-exchange-economy}
A simple exchange economy consists of:
\begin{itemize}
    \item $N$ consumers, indexed by $i=1,\dots,N$;
    \item $M$ goods, indexed by $j=1,\dots,M$;
    \item for each consumer $i$, an endowment
    \[
    \omega^i = (\omega_1^i,\dots,\omega_M^i) \in \mathbb{R}_+^M;
    \]
    \item for each consumer $i$, preferences represented by a utility function
    \[
    U_i(x^i),
    \qquad
    x^i=(x_1^i,\dots,x_M^i)\in \mathbb{R}_+^M.
    \]
\end{itemize}
There is no production, so goods cannot be transformed into one another.
\end{definition}

The total resource of each good is the sum of individual endowments.

\begin{definition}{Aggregate resources and feasible allocations}{ch9-feasible-allocation}
For each good $j$, total resources are
\[
e_j = \sum_{i=1}^N \omega_j^i.
\]
An allocation
\[
X=(x^1,\dots,x^N)
\]
is \emph{feasible} if for every good $j$,
\[
\sum_{i=1}^N x_j^i = e_j.
\]
That is, the allocation must exactly distribute the economy's total endowment.
\end{definition}

A feasible allocation need not respect individual endowments. A consumer may consume more than her own initial endowment of some good, provided other consumers consume less and total resources are preserved.

\begin{examtip}
In exchange-economy questions, always distinguish:
\begin{itemize}
    \item \emph{individual endowment} $\omega^i$,
    \item \emph{individual allocation} $x^i$,
    \item \emph{aggregate resources} $e$.
\end{itemize}
Feasibility is about aggregate resources, not about each person staying within their original bundle.
\end{examtip}

\subsection{Pareto Improvements and Pareto Efficiency}

The chapter replaces the more subjective ``maximize $CS+PS$'' criterion with a more general efficiency concept.

\begin{definition}{Pareto improvement}{ch9-pareto-improvement}
Let
\[
X=(x^1,\dots,x^N)
\quad \text{and} \quad
Y=(y^1,\dots,y^N)
\]
be feasible allocations. We say that $Y$ \emph{Pareto improves on} $X$ if
\[
y^i \succeq_i x^i
\quad \text{for all } i,
\]
and
\[
y^j \succ_j x^j
\quad \text{for some } j.
\]
So $Y$ makes everyone weakly better off and at least one person strictly better off.
\end{definition}

\begin{definition}{Pareto efficient allocation}{ch9-pareto-efficient}
A feasible allocation $X$ is \emph{Pareto efficient} if there is no other feasible allocation that Pareto improves on it.
\end{definition}

The logic is simple: if a feasible reallocation can make someone better off without hurting anyone else, then the original allocation wastes gains from trade.

\begin{commonmistake}
If an allocation $Y$ Pareto improves on $X$, it does \emph{not} follow that $Y$ is Pareto efficient. Another feasible allocation $Z$ may still Pareto improve on $Y$.
\end{commonmistake}

\begin{examplebox}[title={Worked Example: checking Pareto efficiency}]
Suppose 10 apples and 10 oranges are split between consumers A and B.
Consumer A has utility
\[
U_A = \#\text{apples},
\]
while consumer B has utility
\[
U_B = \#\text{apples} + \#\text{oranges}.
\]

Check the following allocations.

\textbf{1. A gets 10 oranges, B gets 10 apples.}

A's utility is $0$ and B's utility is $10$. This is not Pareto efficient because one can transfer some oranges from A to B:
\begin{itemize}
    \item A does not care about oranges;
    \item B values oranges positively.
\end{itemize}
So B can be made strictly better off without hurting A.

\textbf{2. A gets 10 apples, B gets 10 oranges.}

A's utility is $10$ and B's utility is $10$. Any apple transferred from A to B hurts A, and A has no oranges to give B. So no Pareto improvement exists. This allocation is Pareto efficient.

\textbf{3. A gets 5 apples, B gets 5 apples and 10 oranges.}

A's utility is $5$ and B's utility is $15$. Since B already has all oranges, any move that gives A more apples must take apples away from B. That hurts B unless compensated, but A has no desirable oranges to offer. So this allocation is Pareto efficient.

\textbf{4. A gets 0 apples, B gets 10 apples and 10 oranges.}

A's utility is $0$ and B's utility is $20$. Although A is badly treated, any apple transferred from B to A lowers B's utility, and A has nothing valuable to compensate B with. Hence there is no Pareto improvement. This allocation is also Pareto efficient.

Therefore:
\[
\boxed{\text{Allocation 1 is not Pareto efficient; allocations 2, 3, and 4 are Pareto efficient.}}
\]
This example shows that Pareto efficiency and fairness are different concepts.
\end{examplebox}

\subsection{The Edgeworth Box}

For the two-good, two-consumer case, the key graphical tool is the Edgeworth box.

\begin{definition}{Two-good, two-consumer exchange economy}{ch9-2x2-economy}
Suppose there are two goods and two consumers, A and B. Let
\[
\omega^A=(\omega_1^A,\omega_2^A),
\qquad
\omega^B=(\omega_1^B,\omega_2^B).
\]
A feasible allocation consists of bundles
\[
x^A=(x_1^A,x_2^A),
\qquad
x^B=(x_1^B,x_2^B),
\]
satisfying
\[
x_1^A+x_1^B = \omega_1^A+\omega_1^B,
\qquad
x_2^A+x_2^B = \omega_2^A+\omega_2^B.
\]
\end{definition}

\begin{definition}{Edgeworth box}{ch9-edgeworth-box}
The \emph{Edgeworth box} is the rectangle with:
\[
\text{length } = \omega_1^A+\omega_1^B,
\qquad
\text{height } = \omega_2^A+\omega_2^B.
\]
Each point in the box represents a feasible allocation. The lower-left origin is consumer A's origin, while the upper-right origin is consumer B's origin.
\end{definition}

Thus every point in the box automatically satisfies feasibility. The endowment point
\[
(\omega_1^A,\omega_2^A)
\]
from A's perspective simultaneously determines B's endowment from the opposite origin.

\begin{definition}{Contract curve}{ch9-contract-curve}
The \emph{contract curve}, or \emph{Pareto set}, is the set of feasible allocations in the Edgeworth box that are Pareto efficient.
\end{definition}

At a smooth interior Pareto-efficient allocation, the consumers' indifference curves are tangent. Hence, for differentiable utilities,
\[
MRS_A = MRS_B
\]
at interior points of the contract curve.

\begin{theorem}{Tangency characterization of interior Pareto efficiency}{ch9-tangency}
Suppose preferences are locally nonsatiated and differentiable. If a feasible interior allocation in the Edgeworth box is Pareto efficient, then the indifference curves of A and B are tangent at that point. Equivalently,
\[
MRS_A = MRS_B.
\]
\end{theorem}

\begin{proof}
If the indifference curves cross rather than being tangent, then there is a local feasible direction that moves both consumers to higher indifference curves. This creates a Pareto improvement, contradicting Pareto efficiency. Therefore, at an interior Pareto-efficient point, the indifference curves must be tangent, so their slopes and hence their MRS values are equal.
\end{proof}

\begin{examtip}
To identify Pareto improvements in an Edgeworth box:
\begin{enumerate}
    \item draw both consumers' indifference curves through the allocation;
    \item find the region that each consumer weakly prefers;
    \item take the overlap.
\end{enumerate}
If the overlap contains other feasible allocations, the original point is not Pareto efficient.
\end{examtip}

\subsection{Fairness and Envy-Freeness}

Efficiency alone does not address distributional concerns.

\begin{definition}{Envy and envy-freeness}{ch9-envy-free}
In a two-consumer allocation $(x^A,x^B)$:
\begin{itemize}
    \item A envies B if
    \[
    x^B \succ_A x^A
    \quad \text{or equivalently} \quad
    U_A(x^B)>U_A(x^A);
    \]
    \item B envies A if
    \[
    x^A \succ_B x^B.
    \]
\end{itemize}
The allocation is \emph{envy-free} if neither consumer envies the other:
\[
U_A(x^A)\ge U_A(x^B),
\qquad
U_B(x^B)\ge U_B(x^A).
\]
\end{definition}

\begin{definition}{Fair allocation}{ch9-fair-allocation}
A feasible allocation is called \emph{fair} if it is both:
\begin{enumerate}
    \item Pareto efficient, and
    \item envy-free.
\end{enumerate}
\end{definition}

So fairness is stricter than efficiency. A Pareto-efficient allocation can still be highly unequal or envied.

\begin{commonmistake}
Do not equate Pareto efficiency with fairness. An allocation can be Pareto efficient simply because no Pareto improvement is possible, even if one consumer receives almost everything.
\end{commonmistake}

\subsection{Competitive Equilibrium in the Exchange Economy}

Markets now determine allocations through decentralized utility maximization.

\begin{definition}{Competitive equilibrium in a pure exchange economy}{ch9-competitive-equilibrium}
A \emph{competitive equilibrium} consists of a price vector
\[
p=(p_1,\dots,p_M)
\]
and a feasible allocation
\[
(x^1,\dots,x^N)
\]
such that:
\begin{enumerate}
    \item for each consumer $i$, $x^i$ maximizes utility subject to the budget constraint
    \[
    p\cdot x^i = p\cdot \omega^i;
    \]
    \item the allocation is feasible, i.e. markets clear:
    \[
    \sum_{i=1}^N x_j^i = \sum_{i=1}^N \omega_j^i
    \qquad \text{for each good } j.
    \]
\end{enumerate}
\end{definition}

In the two-good, two-consumer Edgeworth box, the budget line passes through the endowment point and has slope
\[
-\frac{p_1}{p_2}
\]
from A's perspective. The same physical line is also B's budget line when drawn from B's origin.

\begin{keyformula}[title={Key Formula: Budget Line in the Edgeworth Box}]
Consumer A's budget line is
\[
p_1 x_1^A + p_2 x_2^A = p_1 \omega_1^A + p_2 \omega_2^A.
\]
Hence
\[
x_2^A
=
\frac{p_1}{p_2}\omega_1^A + \omega_2^A - \frac{p_1}{p_2}x_1^A,
\]
so the slope is
\[
-\frac{p_1}{p_2}.
\]
This same line also represents B's budget constraint after converting coordinates using feasibility.
\end{keyformula}

\begin{examtip}
In exchange-equilibrium problems, only the \emph{price ratio} matters. If
\[
\frac{p_1^*}{p_2^*}=m,
\]
then any positive scalar multiple of $(p_1^*,p_2^*)$ gives the same budget-line slope and the same demand choices. One price is usually normalized to 1 as the numeraire.
\end{examtip}

\subsection{How to Solve for Competitive Equilibrium}

The lecture's solution method is entirely standard.

\begin{examtip}[title={Method Pattern: Solving a 2-good, 2-consumer equilibrium}]
To solve a competitive equilibrium algebraically:
\begin{enumerate}
    \item Derive each consumer's demand as a function of prices and wealth:
    \[
    x^A(p,p\cdot \omega^A), \qquad x^B(p,p\cdot \omega^B).
    \]
    \item Write the market-clearing conditions:
    \[
    x_1^A + x_1^B = \omega_1^A+\omega_1^B,
    \qquad
    x_2^A + x_2^B = \omega_2^A+\omega_2^B.
    \]
    \item Substitute the demand functions and solve for the equilibrium price ratio.
    \item Plug the equilibrium prices back into the demand functions to get the equilibrium allocation.
\end{enumerate}
\end{examtip}

\begin{examplebox}[title={Worked Example: Solving for competitive equilibrium}]
Consider a two-good, two-consumer exchange economy with endowments
\[
\omega^A=(4,8),
\qquad
\omega^B=(7,3).
\]
Consumer A has utility
\[
U_A(x_1,x_2)=\min\{x_1,x_2\},
\]
and consumer B has utility
\[
U_B(x_1,x_2)=x_1^{0.5}x_2^{0.5}.
\]

\textbf{Step 1: derive demands.}

Consumer A has perfect-complements preferences, so she chooses
\[
x_1^A=x_2^A.
\]
Her wealth is
\[
m_A = 4p_1+8p_2.
\]
Hence
\[
x_1^A=x_2^A=\frac{4p_1+8p_2}{p_1+p_2}.
\]

Consumer B has Cobb--Douglas utility with equal exponents, so his demand is
\[
x_1^B=\frac{m_B}{2p_1},
\qquad
x_2^B=\frac{m_B}{2p_2},
\]
where
\[
m_B=7p_1+3p_2.
\]
Therefore
\[
x_1^B=\frac{7p_1+3p_2}{2p_1},
\qquad
x_2^B=\frac{7p_1+3p_2}{2p_2}.
\]

\textbf{Step 2: impose market clearing.}

Total endowment is
\[
(11,11).
\]
Market 1 clearing requires
\[
\frac{4p_1+8p_2}{p_1+p_2}
+
\frac{7p_1+3p_2}{2p_1}
=
11.
\]
Solving gives
\[
p_1=p_2.
\]
By symmetry of total resources, this also clears market 2.

\textbf{Step 3: find the equilibrium allocation.}

Normalize prices so that
\[
p_1=p_2=1.
\]
Then
\[
x_1^A=x_2^A=\frac{4+8}{2}=6,
\]
and
\[
x_1^B=x_2^B=\frac{7+3}{2}=5.
\]

Therefore a competitive equilibrium is
\[
\boxed{
p_1^*=p_2^*,
\qquad
x^{A*}=(6,6),
\qquad
x^{B*}=(5,5).
}
\]
\end{examplebox}

\subsection{Competitive Equilibrium in the Edgeworth Box}

The Edgeworth-box geometry gives a graphical proof of the welfare property of competitive equilibrium.

At a competitive equilibrium:
\begin{itemize}
    \item each consumer maximizes utility subject to the common budget line through the endowment point, so each chosen bundle lies at a tangency between an indifference curve and that budget line;
    \item market clearing means the two consumers choose the same feasible point in the box;
    \item therefore both consumers' indifference curves are tangent to the same line at the same point.
\end{itemize}

Hence the equilibrium allocation lies on the contract curve.

\begin{theorem}{First Fundamental Welfare Theorem in the exchange economy}{ch9-fwt}
In the simple competitive exchange economy, a competitive equilibrium allocation is Pareto efficient.
\end{theorem}

\begin{proof}
At a competitive equilibrium in the two-good, two-consumer Edgeworth box, both consumers maximize utility on the same budget line. Therefore each consumer's indifference curve is tangent to the budget line at the equilibrium allocation. Since the allocation is feasible and both tangencies occur at the same point, the two indifference curves are tangent to each other there. A feasible interior allocation where the indifference curves are tangent admits no local Pareto improvement. Hence the equilibrium allocation is Pareto efficient.
\end{proof}

The lecture emphasizes that this theorem is much less subjective than maximizing $CS+PS$: it does not require attaching interpersonal weights to utility.

\subsection{Existence, Uniqueness, and the Role of Prices}

The lecture notes that under standard convex preferences, the exchange economy has well-behaved equilibrium properties.

\begin{theorem}{Standard existence and uniqueness statement used in this course}{ch9-existence}
Under standard convex preferences in the exchange economy considered in the lecture, a competitive equilibrium exists, and the equilibrium price ratio is unique.
\end{theorem}

\begin{examtip}
Be careful with the word \emph{unique}. In these exchange problems, the equilibrium \emph{price ratio} is unique, not the absolute price vector. If $(p_1^*,p_2^*)$ is an equilibrium price vector, then so is
\[
(\lambda p_1^*, \lambda p_2^*)
\]
for any $\lambda>0$.
\end{examtip}

\subsection{The Second Fundamental Welfare Theorem}

The contract curve usually contains many Pareto-efficient allocations. Competitive equilibrium selects one of them depending on the initial endowment point. The second welfare theorem explains how any desired Pareto-efficient allocation can be supported as a market equilibrium after suitable redistribution.

\begin{theorem}{Second Fundamental Welfare Theorem}{ch9-swt}
Under the standard convexity assumptions, for any Pareto-efficient allocation, there exists a redistribution of endowments such that the resulting competitive equilibrium yields that allocation.
\end{theorem}

\begin{proof}
Take a Pareto-efficient allocation in the Edgeworth box. At such a point, the consumers' indifference curves admit a common supporting tangent line. Choose a redistributed endowment point lying on this tangent budget line. At the prices corresponding to the slope of that line, each consumer can afford the Pareto-efficient allocation and optimally chooses it, since the line is tangent to the relevant indifference curve. Because the allocation is feasible, markets clear. Hence the Pareto-efficient allocation is decentralized as a competitive equilibrium from the redistributed endowment.
\end{proof}

So the theorem separates:
\begin{itemize}
    \item \textbf{efficiency}: achieved by competitive trade;
    \item \textbf{distribution}: achieved by the initial redistribution of endowments.
\end{itemize}

This is one of the lecture's main messages: if society has a preferred Pareto-efficient allocation, it may be possible to reach it not by distorting trade itself, but by redistributing endowments first and then letting markets operate.

\begin{examtip}
A standard Edgeworth-box interpretation of the SWT is:
\begin{itemize}
    \item pick a Pareto-efficient point on the contract curve;
    \item draw the common tangent budget line there;
    \item any redistributed endowment lying on that line can support that Pareto-efficient point as competitive equilibrium.
\end{itemize}
\end{examtip}

\subsection{Implicit Assumptions Behind the Welfare Theorems}

The welfare theorems are powerful, but they do not hold without structure.

\begin{definition}{Implicit assumptions emphasized in the lecture}{ch9-assumptions}
The exchange model in this chapter relies on the following assumptions:
\begin{itemize}
    \item \textbf{Complete markets:} all relevant goods can be traded.
    \item \textbf{No externalities:} each consumer's utility depends only on own consumption.
    \item \textbf{Standard convex preferences:} needed for the clean welfare-theorem results.
\end{itemize}
\end{definition}

If external effects exist but cannot be traded, or if some markets are missing, then the competitive equilibrium need not be Pareto efficient.

\begin{commonmistake}
Do not quote the welfare theorems without mentioning their assumptions. In particular, missing markets and externalities are standard reasons why competitive equilibrium can fail to be Pareto efficient.
\end{commonmistake}

\subsection{Standard Problem Types and Solution Patterns}

\begin{examtip}[title={Method Pattern: Checking Pareto efficiency}]
In allocation questions:
\begin{enumerate}
    \item identify total resources and feasibility;
    \item ask whether one consumer can be made strictly better off without making the other worse off;
    \item in Edgeworth-box settings, look for overlap of upper contour sets;
    \item at smooth interior points, check tangency:
    \[
    MRS_A=MRS_B.
    \]
\end{enumerate}
\end{examtip}

\begin{examtip}[title={Method Pattern: Exchange equilibrium algebra}]
For two-good, two-consumer economies:
\begin{enumerate}
    \item compute each consumer's wealth from endowment:
    \[
    m_i = p_1 \omega_1^i + p_2 \omega_2^i;
    \]
    \item derive Marshallian demand from each utility function;
    \item impose market clearing on one or both goods;
    \item solve for the price ratio and allocations.
\end{enumerate}
\end{examtip}

\begin{examtip}[title={Method Pattern: Fairness questions}]
If asked whether an allocation is fair:
\begin{enumerate}
    \item test Pareto efficiency;
    \item test envy-freeness:
    \[
    U_A(x^A)\ge U_A(x^B),
    \qquad
    U_B(x^B)\ge U_B(x^A).
    \]
\end{enumerate}
Only if both hold is the allocation fair.
\end{examtip}

\subsection{Chapter Summary}

\begin{keyformula}[title={Key Formula: Lecture 9 Summary}]
\begin{align*}
e_j &= \sum_{i=1}^N \omega_j^i, \\
\sum_{i=1}^N x_j^i &= e_j
\qquad \text{for each good } j, \\
Y \text{ Pareto improves on } X
&\iff
y^i \succeq_i x^i \ \forall i
\text{ and } y^j \succ_j x^j \text{ for some } j, \\
X \text{ Pareto efficient}
&\iff
\text{no feasible allocation Pareto improves on } X, \\
p\cdot x^i &= p\cdot \omega^i
\qquad \text{for each consumer in equilibrium}, \\
x_1^A+x_1^B &= \omega_1^A+\omega_1^B, \\
x_2^A+x_2^B &= \omega_2^A+\omega_2^B, \\
\text{interior Pareto efficiency} &\Rightarrow MRS_A=MRS_B, \\
\text{competitive equilibrium} &\Rightarrow \text{Pareto efficiency (FWT)}, \\
\text{any Pareto-efficient allocation}
&\Rightarrow
\text{can be decentralized after redistribution (SWT)}.
\end{align*}
\end{keyformula}

\begin{examtip}
The highest-yield takeaways from this chapter are:
\begin{itemize}
    \item feasibility in an exchange economy means total consumption equals total endowment good by good;
    \item Pareto efficiency is about the impossibility of Pareto improvements, not about fairness;
    \item the Edgeworth box turns a two-person, two-good exchange economy into a geometric object;
    \item the contract curve is the set of Pareto-efficient allocations;
    \item a competitive equilibrium is determined by utility maximization plus market clearing;
    \item in equilibrium, only the price ratio matters;
    \item the First Welfare Theorem says competitive equilibrium is Pareto efficient;
    \item the Second Welfare Theorem says any Pareto-efficient allocation can be supported after suitable redistribution, under the standard assumptions.
\end{itemize}
\end{examtip}
\newpage
% chapters/ch10_social_choice.tex

\section{Social Choice}

This chapter studies how a society should choose among alternatives when different individuals have different preferences. In earlier chapters, we developed efficiency concepts such as Pareto efficiency and showed that competitive equilibrium can support efficient allocations under standard assumptions. However, efficiency alone often does not determine a unique social choice. This chapter introduces social decision mechanisms, their desirable properties, Arrow's impossibility theorem, and social welfare functions. The main conceptual distinction is between:
\begin{itemize}
    \item \textbf{ordinal aggregation:} using only rankings of alternatives;
    \item \textbf{cardinal aggregation:} using utility levels or utility intensities.
\end{itemize}

\subsection{Alternatives and Individual Preferences}

A social choice problem begins with a set of alternatives and a collection of individual preferences over those alternatives.

\begin{definition}{Alternatives}{ch10-alternatives}
An \emph{alternative} is one feasible option from which society may choose. Examples include:
\begin{itemize}
    \item possible lunch venues;
    \item political candidates or parties;
    \item feasible allocations in an Edgeworth box;
    \item policy regimes such as different tax schedules.
\end{itemize}
\end{definition}

\begin{definition}{Individual preferences}{ch10-individual-preferences}
Let the set of alternatives be $A$. For individual $i$:
\begin{itemize}
    \item $x \succeq_i y$ means $i$ weakly prefers $x$ to $y$;
    \item $x \succ_i y$ means $i$ strictly prefers $x$ to $y$;
    \item $x \sim_i y$ means $i$ is indifferent between $x$ and $y$.
\end{itemize}
When the lecture restricts attention to strict preferences for simplicity, each individual can rank all alternatives from most preferred to least preferred.
\end{definition}

\begin{definition}{Rational preferences}{ch10-rational-preferences}
Preferences are \emph{well-defined} or \emph{rational} if they satisfy:
\begin{itemize}
    \item \textbf{Completeness:} for any two alternatives $x,y$, either $x \succeq_i y$ or $y \succeq_i x$ (or both);
    \item \textbf{Transitivity:} if $x \succeq_i y$ and $y \succeq_i z$, then $x \succeq_i z$.
\end{itemize}
\end{definition}

\begin{examtip}
In this chapter, many examples are written using \emph{strict} rankings only. When that happens, completeness means that the individual can rank every pair of alternatives, and transitivity means the ranking contains no cycles.
\end{examtip}

\begin{examplebox}[title={Worked Example: Checking completeness and transitivity}]
Suppose individual $j$ has preferences over $\{a,b,c,d\}$:
\[
a \succ_j b,\qquad a \succ_j c,\qquad a \succ_j d,\qquad b \succ_j c,\qquad c \succ_j d.
\]

These comparisons imply:
\[
a \succ_j b \succ_j c \succ_j d,
\]
and therefore also
\[
b \succ_j d,\qquad a \succ_j d,\qquad a \succ_j c.
\]
So the ranking is complete and transitive.

Now suppose individual $k$ has:
\[
a \succ_k b,\qquad a \succ_k c,\qquad a \succ_k d,\qquad b \succ_k c,\qquad d \succ_k b,\qquad c \succ_k d.
\]
Then we have
\[
b \succ_k c,\qquad c \succ_k d,\qquad d \succ_k b,
\]
which forms a cycle. Hence transitivity fails.

So:
\[
\boxed{\text{$j$'s preferences are complete and transitive, while $k$'s are not transitive.}}
\]
\end{examplebox}

\subsection{The Social Choice Problem}

Once individual preferences are specified, the next problem is to aggregate them into a social ranking.

\begin{definition}{Social choice problem}{ch10-social-choice-problem}
A social choice problem consists of:
\begin{itemize}
    \item a set of individuals $\{1,2,\dots,n\}$;
    \item a set of alternatives $A$;
    \item for each individual $i$, a preference ordering over $A$.
\end{itemize}
The objective is to use these individual preferences to generate a social preference or social choice.
\end{definition}

\begin{definition}{Social decision mechanism}{ch10-sdm}
A \emph{social decision mechanism} (SDM) is a rule that takes individual preference orderings
\[
\succ_1,\dots,\succ_n
\]
as inputs and outputs a social preference ordering
\[
\succ^*.
\]
\end{definition}

The key point is that an SDM uses only individuals' rankings of alternatives. It does \emph{not} use utility intensities.

\begin{examtip}
A social decision mechanism is an \emph{ordinal} aggregator. It uses rankings such as
\[
x \succ_i y \succ_i z,
\]
not numerical statements such as ``$i$ cares twice as much about $x$ as about $y$.'' This distinction becomes central when the chapter moves to social welfare functions.
\end{examtip}

\subsection{Examples of Social Decision Mechanisms}

\subsubsection{Pairwise majority voting}

\begin{definition}{Pairwise majority voting}{ch10-pairwise-majority}
For each pair of alternatives $(x,y)$, compare how many individuals prefer $x$ to $y$ and how many prefer $y$ to $x$. The socially preferred alternative is the one that wins the majority vote in that pairwise comparison.
\end{definition}

\begin{examplebox}[title={Worked Example: Pairwise majority voting}]
Suppose three individuals have rankings:
\[
\text{Harry: } X \succ Y \succ Z,
\qquad
\text{Ron: } Z \succ Y \succ X,
\qquad
\text{Hermione: } Z \succ X \succ Y.
\]

Compare pairwise:
\begin{itemize}
    \item $X$ vs $Y$: Harry and Hermione prefer $X$, so
    \[
    X \succ^* Y.
    \]
    \item $X$ vs $Z$: Ron and Hermione prefer $Z$, so
    \[
    Z \succ^* X.
    \]
    \item $Y$ vs $Z$: Ron and Hermione prefer $Z$, so
    \[
    Z \succ^* Y.
    \]
\end{itemize}

Thus the social ordering is
\[
\boxed{Z \succ^* X \succ^* Y.}
\]
\end{examplebox}

\subsubsection{Borda count}

\begin{definition}{Borda count}{ch10-borda}
Under the Borda count, each individual assigns rank-points to alternatives. With $m$ alternatives, the top-ranked alternative gets 1 point, the next gets 2 points, and so on until the last gets $m$ points. The social ranking orders alternatives from lowest total point score to highest.
\end{definition}

\begin{examplebox}[title={Worked Example: Borda count}]
Suppose 100 individuals fall into three preference types:
\[
40\%: X \succ Y \succ Z,
\qquad
25\%: Z \succ Y \succ X,
\qquad
35\%: Z \succ X \succ Y.
\]
With three alternatives, ranks 1, 2, and 3 correspond to 1, 2, and 3 points.

Compute total points:
\[
X: 40(1)+25(3)+35(2)=40+75+70=185,
\]
\[
Y: 40(2)+25(2)+35(3)=80+50+105=235,
\]
\[
Z: 40(3)+25(1)+35(1)=120+25+35=180.
\]
Lower points are better, so the social ranking is
\[
\boxed{Z \succ^* X \succ^* Y.}
\]
\end{examplebox}

\subsubsection{Plurality voting}

\begin{definition}{Plurality voting}{ch10-plurality}
Plurality voting ranks first the alternative that is placed top by the largest number of individuals. After assigning that rank, remove the chosen alternative and repeat the procedure among the remaining alternatives.
\end{definition}

\begin{examplebox}[title={Worked Example: Plurality voting}]
Using the same preference profile:
\[
40\%: X \succ Y \succ Z,
\qquad
25\%: Z \succ Y \succ X,
\qquad
35\%: Z \succ X \succ Y,
\]
the top-choice shares are:
\[
X:40\%,\qquad Y:0\%,\qquad Z:60\%.
\]
So $Z$ is ranked first. Remove $Z$ and compare $X$ and $Y$:
\begin{itemize}
    \item the first group prefers $X$ to $Y$;
    \item the second group prefers $Y$ to $X$;
    \item the third group prefers $X$ to $Y$.
\end{itemize}
Thus $X$ beats $Y$ by $75\%$ to $25\%$.

So the social ranking is
\[
\boxed{Z \succ^* X \succ^* Y.}
\]
\end{examplebox}

\subsection{Why Social Decision Mechanisms Can Fail}

An arbitrary rule that outputs a social ranking need not respect individual preferences or even produce a well-defined ordering.

\begin{commonmistake}
Do not assume that every rule taking rankings as inputs is a sensible social decision mechanism. The chapter's central message is precisely that many apparently reasonable rules fail some basic desiderata.
\end{commonmistake}

\subsubsection{Condorcet cycles}

\begin{definition}{Condorcet cycle}{ch10-condorcet}
A \emph{Condorcet cycle} occurs when pairwise majority comparisons produce a cycle such as
\[
X \succ^* Y,\qquad Y \succ^* Z,\qquad Z \succ^* X.
\]
This means the social ranking is not transitive.
\end{definition}

\begin{examplebox}[title={Worked Example: Condorcet cycle under pairwise majority voting}]
Suppose three individuals have preferences:
\[
\text{Type 1: } X \succ Z \succ Y,
\qquad
\text{Type 2: } Z \succ Y \succ X,
\qquad
\text{Type 3: } Y \succ X \succ Z.
\]

Then:
\begin{itemize}
    \item $Y$ beats $X$ because types 2 and 3 prefer $Y$ to $X$:
    \[
    Y \succ^* X.
    \]
    \item $X$ beats $Z$ because types 1 and 3 prefer $X$ to $Z$:
    \[
    X \succ^* Z.
    \]
    \item $Z$ beats $Y$ because types 1 and 2 prefer $Z$ to $Y$:
    \[
    Z \succ^* Y.
    \]
\end{itemize}

Hence
\[
Y \succ^* X,\qquad X \succ^* Z,\qquad Z \succ^* Y,
\]
which is a cycle. Therefore pairwise majority voting does not produce a well-defined transitive social ordering in this case.
\end{examplebox}

\subsubsection{Vote splitting and spoiling}

Some mechanisms are sensitive to the presence of other alternatives.

\begin{examplebox}[title={Worked Example: Plurality voting and vote splitting}]
Suppose preference types are:
\[
40\%: X \succ Y \succ Z,
\qquad
25\%: Z \succ Y \succ X,
\qquad
35\%: Y \succ X \succ Z.
\]
Under plurality voting, the top-choice shares are:
\[
X:40\%,\qquad Y:35\%,\qquad Z:25\%.
\]
So plurality ranks
\[
X \succ^* Y \succ^* Z.
\]

However, a direct comparison between $X$ and $Y$ gives:
\begin{itemize}
    \item the first type prefers $X$ to $Y$;
    \item the second and third types prefer $Y$ to $X$.
\end{itemize}
So $60\%$ prefer $Y$ to $X$.

Thus plurality makes $X$ the winner even though a majority would prefer $Y$ to $X$ in a head-to-head contest. The reason is that $Z$ splits support away from $Y$.
\end{examplebox}

\subsection{Desirable Properties of Social Decision Mechanisms}

The lecture focuses on three benchmark properties.

\begin{definition}{Pareto property for an SDM}{ch10-pareto-sdm}
An SDM satisfies \emph{Pareto} if whenever every individual prefers $x$ to $y$, the social ranking also prefers $x$ to $y$:
\[
x \succ_i y \ \text{for all } i
\quad \Longrightarrow \quad
x \succ^* y.
\]
\end{definition}

\begin{definition}{Universal Domain}{ch10-ud}
An SDM satisfies \emph{Universal Domain} if for every possible profile of well-defined individual preferences, the resulting social preference ordering is itself well-defined.
\end{definition}

\begin{definition}{Independence of Irrelevant Alternatives}{ch10-iia}
An SDM satisfies \emph{Independence of Irrelevant Alternatives} (IIA) if for any two alternatives $A$ and $B$, the social ranking of $A$ versus $B$ depends only on individuals' rankings of $A$ versus $B$, and not on how they rank other alternatives.
\end{definition}

\begin{examtip}
To check IIA:
\begin{itemize}
    \item keep every individual's ranking of $A$ versus $B$ fixed;
    \item change rankings involving other alternatives;
    \item see whether the social ranking of $A$ versus $B$ changes.
\end{itemize}
If it changes, IIA fails.
\end{examtip}

\subsection{How to Check Whether a Property Holds}

A property must hold for the \emph{entire multiverse} of possible preference profiles.

\begin{examtip}[title={Method Pattern: verifying SDM properties}]
To show a property \emph{does not hold}, it is enough to give a single counterexample.

To show a property \emph{does hold}, one must give a general logical argument explaining why no counterexample can exist.
\end{examtip}

\begin{examplebox}[title={Worked Example: Pairwise majority voting and the three properties}]
Consider pairwise majority voting.

\textbf{Pareto:} If every individual prefers $x$ to $y$, then in the pairwise vote between $x$ and $y$, every individual votes for $x$. Hence
\[
x \succ^* y.
\]
So pairwise majority voting satisfies Pareto.

\textbf{Universal Domain:} The Condorcet-cycle example gives well-defined individual preferences but produces a non-transitive social ordering. Hence pairwise majority voting does \emph{not} satisfy Universal Domain.

\textbf{IIA:} In pairwise majority voting, the social ranking of $x$ versus $y$ depends only on how individuals rank $x$ against $y$ in the pairwise contest. Changing rankings involving other alternatives does not affect that contest. Hence pairwise majority voting satisfies IIA.

Therefore:
\[
\boxed{\text{Pairwise majority voting satisfies Pareto and IIA, but not Universal Domain.}}
\]
\end{examplebox}

\begin{examplebox}[title={Worked Example: Pairwise minority voting}]
Define the pairwise minority voting rule as follows: for each pair of alternatives, the socially preferred alternative is the one that \emph{loses} the majority vote.

Check the three properties.

\textbf{Pareto:} If everyone prefers $x$ to $y$, then majority voting selects $x$, so minority voting selects $y$. Hence the social preference contradicts unanimous preference. So Pareto fails.

\textbf{Universal Domain:} Since pairwise minority voting simply reverses pairwise-majority outcomes, a Condorcet cycle under majority voting becomes another cycle under minority voting. Hence Universal Domain fails.

\textbf{IIA:} For any pair $(x,y)$, the minority-voting outcome depends only on individuals' rankings of $x$ versus $y$. So IIA holds.

Therefore:
\[
\boxed{\text{Pairwise minority voting satisfies IIA only; Pareto and Universal Domain both fail.}}
\]
\end{examplebox}

\subsection{Arrow's Impossibility Theorem}

The three conditions above are all plausible. Arrow's theorem shows that they are jointly incompatible with non-dictatorial social choice once there are at least three alternatives.

\begin{definition}{Dictatorship}{ch10-dictatorship}
An SDM is a \emph{dictatorship} if there exists some individual $d$ such that for every pair of alternatives $x,y$,
\[
x \succ_d y
\quad \Longrightarrow \quad
x \succ^* y.
\]
That is, the social ordering always follows the dictator's preferences.
\end{definition}

\begin{theorem}{Arrow's Impossibility Theorem}{ch10-arrow}
Suppose there are at least three alternatives. If a social decision mechanism satisfies:
\begin{enumerate}
    \item Pareto,
    \item Universal Domain,
    \item Independence of Irrelevant Alternatives,
\end{enumerate}
then it must be a dictatorship.
\end{theorem}

The lecture uses this theorem as a negative result: there is no perfect democratic ordinal rule satisfying all three desiderata for all preference profiles.

\begin{commonmistake}
Arrow's theorem does \emph{not} say that every voting rule is a dictatorship. It says that any SDM satisfying all three listed properties on the full domain of preferences must be dictatorial.
\end{commonmistake}

\begin{examtip}
A common interpretation is:
\begin{quote}
With at least three alternatives, one cannot have a social decision mechanism that is simultaneously fully general, respects unanimity, ignores irrelevant alternatives, and is non-dictatorial.
\end{quote}
So at least one condition must be weakened if one wants a non-dictatorial rule.
\end{examtip}

\subsection{Social Welfare Functions}

One way to escape Arrow's impossibility is to use more information than ordinal rankings alone.

\begin{definition}{Social welfare function}{ch10-swf}
A \emph{social welfare function} (SWF) is a function
\[
W(u_1,u_2,\dots,u_n)
\]
that takes individuals' utility levels as inputs and outputs a real number used to rank social outcomes. The function is increasing in each argument.
\end{definition}

This differs from an SDM:
\begin{itemize}
    \item an SDM aggregates \emph{ordinal} preference rankings;
    \item an SWF aggregates \emph{cardinal} utilities, and so can reflect preference intensities.
\end{itemize}

\begin{definition}{Common examples of SWFs}{ch10-swf-examples}
Common examples include:
\begin{itemize}
    \item \textbf{Utilitarian:}
    \[
    W(u_1,\dots,u_n)=\sum_{i=1}^n u_i;
    \]
    \item \textbf{Weighted utilitarian:}
    \[
    W(u_1,\dots,u_n)=\sum_{i=1}^n \alpha_i u_i,
    \qquad \alpha_i>0;
    \]
    \item \textbf{Rawlsian:}
    \[
    W(u_1,\dots,u_n)=\min\{u_1,\dots,u_n\}.
    \]
\end{itemize}
\end{definition}

\begin{examtip}
An SWF can be sensitive to \emph{preference intensities}. This is exactly what an IIA-style ordinal SDM rules out. That is one reason the SWF approach can bypass Arrow's impossibility.
\end{examtip}

\subsection{Utility Possibility Set and Social Optima}

To connect social welfare to Pareto efficiency, the lecture introduces the utility possibility set.

\begin{definition}{Utility possibility set}{ch10-ups}
The \emph{utility possibility set} is the set of utility vectors
\[
(u_A,u_B)
\]
that can be achieved by feasible allocations in the economy.
\end{definition}

\begin{definition}{Utility possibility frontier}{ch10-upf}
The \emph{utility possibility frontier} (UPF) is the boundary of the utility possibility set consisting of efficient utility pairs.
\end{definition}

A social optimum is obtained by maximizing a social welfare function over the utility possibility set.

\begin{definition}{Social optimum}{ch10-social-optimum}
A feasible allocation is \emph{socially optimal} for a given SWF if it maximizes
\[
W(u_1,\dots,u_n)
\]
over the feasible set of utility outcomes.
\end{definition}

Graphically, a social optimum occurs where the highest attainable iso-welfare curve touches the utility possibility frontier.

\begin{theorem}{Social optimum implies Pareto efficiency}{ch10-so-implies-pe}
If an allocation is socially optimal for some increasing social welfare function, then it is Pareto efficient.
\end{theorem}

\begin{proof}
Suppose a socially optimal allocation were not Pareto efficient. Then there would exist another feasible allocation that makes everyone weakly better off and at least one person strictly better off. Since the social welfare function is increasing in each individual's utility, the second allocation would yield strictly higher social welfare. This contradicts social optimality. Hence every social optimum must be Pareto efficient.
\end{proof}

\begin{theorem}{Pareto efficiency implies social optimality for some SWF}{ch10-pe-implies-so}
Under the regularity conditions assumed in the lecture, every Pareto-efficient allocation is socially optimal for some social welfare function.
\end{theorem}

The lecture emphasizes this as a positive converse: one need not fear that a Pareto-efficient allocation is unsupported by all welfare criteria. Some SWF can always rationalize it.

\subsection{Weighted Utilitarianism and Supporting Pareto-Efficient Points}

For two consumers, weighted utilitarian welfare takes the form
\[
W(u_A,u_B)=\gamma_A u_A + \gamma_B u_B,
\qquad
\gamma_A,\gamma_B>0.
\]
Its iso-welfare curves are straight lines.

\begin{keyformula}[title={Key Formula: Iso-Welfare Slope under Weighted Utilitarianism}]
For
\[
W(u_A,u_B)=\gamma_A u_A+\gamma_B u_B,
\]
an iso-welfare curve satisfies
\[
\gamma_A u_A+\gamma_B u_B=\bar W.
\]
Hence
\[
u_B=\frac{\bar W}{\gamma_B}-\frac{\gamma_A}{\gamma_B}u_A,
\]
so the slope is
\[
-\frac{\gamma_A}{\gamma_B}.
\]
At a smooth social optimum, this slope equals the slope of the utility possibility frontier.
\end{keyformula}

\begin{examplebox}[title={Worked Example: Which weights support a given social optimum?}]
Suppose point $X$ on the utility possibility frontier has tangent slope
\[
-\frac{1}{2}.
\]
A weighted utilitarian social welfare function supports $X$ if
\[
-\frac{\gamma_A}{\gamma_B}=-\frac12
\quad \Longrightarrow \quad
\frac{\gamma_A}{\gamma_B}=\frac12.
\]
So any positive weights in the ratio
\[
\gamma_A:\gamma_B=1:2
\]
work. Hence both
\[
(\gamma_A,\gamma_B)=(1,2)
\quad \text{and} \quad
(\gamma_A,\gamma_B)=(2,4)
\]
support $X$.

Now suppose point $Y$ has tangent slope
\[
-2.
\]
Then
\[
-\frac{\gamma_A}{\gamma_B}=-2
\quad \Longrightarrow \quad
\frac{\gamma_A}{\gamma_B}=2.
\]
So any positive weights in the ratio
\[
\gamma_A:\gamma_B=2:1
\]
work. Hence both
\[
(\gamma_A,\gamma_B)=(2,1)
\quad \text{and} \quad
(\gamma_A,\gamma_B)=(4,2)
\]
support $Y$.

Therefore:
\[
\boxed{X \text{ is supported by weights proportional to } (1,2),}
\]
\[
\boxed{Y \text{ is supported by weights proportional to } (2,1).}
\]
\end{examplebox}

\subsection{Relation to the Welfare Theorems}

Chapter 9 established that, under the standard assumptions:
\begin{itemize}
    \item competitive equilibrium is Pareto efficient;
    \item every Pareto-efficient allocation can be supported by a suitable redistribution of endowments followed by competitive trade.
\end{itemize}

This chapter adds:
\begin{itemize}
    \item every social optimum is Pareto efficient;
    \item every Pareto-efficient allocation is socially optimal for some SWF.
\end{itemize}

Combining the two chapters gives a bridge:
\begin{quote}
Under the standard assumptions and in the absence of externalities, one can connect competitive equilibria, Pareto-efficient allocations, and social optima.
\end{quote}

\begin{policynote}
The important limitation is that society must still choose \emph{which} social welfare function to use. The mathematical existence of an SWF supporting a Pareto-efficient point does not by itself solve the normative problem of deciding which efficient point is morally or politically preferable.
\end{policynote}

\subsection{Common Mistakes and Exam Traps}

\begin{commonmistake}
Do not confuse a social decision mechanism with a social welfare function. An SDM takes rankings as input; an SWF takes utility levels as input.
\end{commonmistake}

\begin{commonmistake}
Arrow's theorem is about SDMs defined on ordinal preference profiles. It does not directly rule out social welfare functions based on utility levels.
\end{commonmistake}

\begin{commonmistake}
Do not say that Pareto-efficient allocations are automatically fair or socially optimal for every SWF. The theorem says they are socially optimal for \emph{some} SWF.
\end{commonmistake}

\begin{commonmistake}
When checking IIA, do not change individuals' rankings of the pair being tested. The point is to keep their ranking of $A$ versus $B$ fixed and change only preferences involving other alternatives.
\end{commonmistake}

\subsection{Chapter Summary}

\begin{keyformula}[title={Key Formula: Lecture 10 Summary}]
\begin{align*}
x \succeq_i y &\quad \text{weak preference}, \\
x \succ_i y &\quad \text{strict preference}, \\
x \sim_i y &\quad \text{indifference}, \\
\text{Pareto for SDM: } 
x \succ_i y \ \forall i &\implies x \succ^* y, \\
\text{Universal Domain} &:\ \text{all well-defined individual preferences yield a well-defined social ranking}, \\
\text{IIA} &:\ \text{social ranking of } A \text{ vs } B \text{ depends only on rankings of } A \text{ vs } B, \\
W(u_1,\dots,u_n) &\quad \text{social welfare function}, \\
W(u_1,\dots,u_n)=\sum_i \alpha_i u_i &\quad \text{weighted utilitarian SWF}, \\
\text{iso-welfare slope} &= -\frac{\gamma_A}{\gamma_B}, \\
\text{social optimum} &\implies \text{Pareto efficiency}, \\
\text{Pareto efficiency} &\implies \text{social optimum for some SWF}.
\end{align*}
\end{keyformula}

\begin{examtip}
The highest-yield takeaways from this chapter are:
\begin{itemize}
    \item know the definitions of weak preference, strict preference, indifference, completeness, and transitivity;
    \item know the difference between SDMs and SWFs;
    \item be able to test Pareto, Universal Domain, and IIA;
    \item know why pairwise majority voting fails Universal Domain;
    \item understand Arrow's Impossibility Theorem and what it does \emph{not} say;
    \item understand why social optima lie on the Pareto frontier;
    \item be able to connect weighted utilitarian slopes
    \[
    -\frac{\gamma_A}{\gamma_B}
    \]
    to tangencies on the utility possibility frontier.
\end{itemize}
\end{examtip}
\newpage
% chapters/ch011_utilitarianism.tex

\section{Utilitarianism}

This chapter gives a mathematical argument for why utilitarianism can be a reasonable social welfare criterion. The key idea is that if utility is quasi-linear in money and society can implement balanced money transfers, then the policy problem separates into:
\begin{enumerate}
    \item \textbf{efficiency:} choose the allocation that maximizes the total utility generated by goods;
    \item \textbf{distribution:} use money transfers to redistribute that utility pie.
\end{enumerate}
Under these conditions, Pareto-efficient policies are exactly those built from utilitarian allocations.

\subsection{Mathematical Logic for Economic Statements}

Lecture 11 begins with a short review of mathematical logic because many economic results in this course take the form of conditional statements.

\begin{definition}{Conditional statement, contrapositive, and converse}{ch11-logic}
A \emph{conditional statement} is a proposition of the form
\[
A \to B,
\]
read as ``If $A$, then $B$.''

Its \emph{contrapositive} is
\[
\neg B \to \neg A,
\]
and its \emph{converse} is
\[
B \to A.
\]
\end{definition}

\begin{theorem}{Conditional statement and contrapositive are equivalent}{ch11-contrapositive}
The statements
\[
A \to B
\qquad \text{and} \qquad
\neg B \to \neg A
\]
are logically equivalent.
\end{theorem}

\begin{proof}
A conditional statement fails only in the case where $A$ is true and $B$ is false. But this is exactly the case where $\neg B$ is true and $\neg A$ is false, so the contrapositive also fails in precisely the same states of the world. Therefore the two statements always have the same truth value.
\end{proof}

\begin{commonmistake}
Do not confuse the \emph{contrapositive} with the \emph{converse}. In general,
\[
A \to B
\]
does \emph{not} imply
\[
B \to A.
\]
\end{commonmistake}

\begin{examplebox}[title={Worked Example: Using a contrapositive}]
Consider the statement:
\[
\text{If } x+y=1000,\ \text{then either } x\ge 500 \text{ or } y\ge 500.
\]
A direct proof must consider all cases with $x+y=1000$. It is easier to prove the contrapositive:
\[
\text{If } x<500 \text{ and } y<500,\ \text{then } x+y\neq 1000.
\]
But if both $x$ and $y$ are less than 500, then
\[
x+y<1000.
\]
So the contrapositive is true, and hence the original statement is also true.
\end{examplebox}

\begin{examtip}
In this lecture, several results are proved most cleanly by contraposition:
\[
A \to B
\quad \Longleftrightarrow \quad
\neg B \to \neg A.
\]
This is often easier than proving the original implication directly.
\end{examtip}

\subsection{Rawlsian and Utilitarian Social Welfare}

Lecture 10 introduced social welfare functions. Lecture 11 contrasts two central examples.

\begin{definition}{Rawlsian and utilitarian social welfare functions}{ch11-rawls-util}
For utility vector
\[
u=(u_1,\dots,u_n),
\]
two common social welfare functions are:
\begin{itemize}
    \item \textbf{Rawlsian:}
    \[
    W_R(u)=\min\{u_1,\dots,u_n\};
    \]
    \item \textbf{Utilitarian:}
    \[
    W_U(u)=\sum_{i=1}^n u_i.
    \]
\end{itemize}
\end{definition}

The Rawlsian criterion puts maximal emphasis on the worst-off individual, while the utilitarian criterion puts equal weight on increasing the sum of utilities.

\begin{policynote}
The lecture's goal is not to prove that utilitarianism is always morally correct. Rather, it gives a precise economic argument showing that utilitarian allocations have a special efficiency property once money transfers and quasi-linear utility are introduced.
\end{policynote}

\subsection{Social Choice with Money Transfers}

The lecture now enriches the social choice problem by allowing monetary transfers alongside allocations of goods.

\begin{definition}{Allocation, transfer scheme, and policy}{ch11-policy}
A \emph{transfer scheme} is a vector
\[
t=(t_1,\dots,t_n),
\]
where:
\begin{itemize}
    \item $t_i>0$ means individual $i$ receives money;
    \item $t_i<0$ means individual $i$ pays money.
\end{itemize}
A \emph{policy} is a pair
\[
(A,t),
\]
consisting of:
\begin{itemize}
    \item an allocation $A$ of goods;
    \item a monetary transfer scheme $t$.
\end{itemize}
\end{definition}

\begin{definition}{Balanced transfers}{ch11-balanced}
A transfer scheme is \emph{budget balanced} if
\[
\sum_{i=1}^n t_i = 0.
\]
Thus money is only redistributed across individuals; no outside surplus or deficit is created.
\end{definition}

\begin{theorem}{Why balanced transfers are needed}{ch11-balanced-needed}
If unrestricted transfer schemes are allowed, Pareto efficiency of policies becomes trivial: any policy can be Pareto improved upon by giving every individual additional money.
\end{theorem}

\begin{proof}
Suppose $(A,t)$ is any policy. If transfers need not balance, choose $\varepsilon>0$ and consider
\[
t'_i=t_i+\varepsilon
\qquad \text{for all } i.
\]
If utility is increasing in money, then every individual is strictly better off under $(A,t')$ than under $(A,t)$. Hence $(A,t)$ cannot be Pareto efficient. Therefore meaningful analysis requires balanced transfers.
\end{proof}

\begin{commonmistake}
Do not confuse an \emph{allocation} with a \emph{policy}. In this chapter:
\begin{itemize}
    \item an allocation concerns only goods;
    \item a policy concerns goods \emph{and} transfers.
\end{itemize}
Pareto efficiency is applied to policies, not just to allocations.
\end{commonmistake}

\subsection{Quasi-Linear Utility and Willingness to Pay}

The central assumption is that individuals value money directly and linearly.

\begin{definition}{Quasi-linear utility with transfers}{ch11-quasilinear}
Suppose that if allocation $A$ is chosen and individual $i$ receives transfer $t_i$, utility is
\[
U_i(A,t_i)=v_i(A)+t_i,
\]
where:
\begin{itemize}
    \item $v_i(A)$ is the utility derived from the allocation itself;
    \item $t_i$ is money, valued one-for-one in utility.
\end{itemize}
\end{definition}

This means money can be used to transfer utility across individuals in a direct way.

\begin{theorem}{Willingness to pay under quasi-linear utility}{ch11-wtp}
Suppose the status quo is allocation $A$, and the new allocation is $B$. If moving from $A$ to $B$ requires individual $i$ to pay $c_i$, then $i$ prefers the change if and only if
\[
c_i \le v_i(B)-v_i(A).
\]
Hence
\[
v_i(B)-v_i(A)
\]
is individual $i$'s maximum willingness to pay to move from $A$ to $B$.
\end{theorem}

\begin{proof}
Individual $i$ weakly prefers the new policy if
\[
U_i(B,-c_i)\ge U_i(A,0).
\]
Using quasi-linearity,
\[
v_i(B)-c_i \ge v_i(A),
\]
which is equivalent to
\[
c_i \le v_i(B)-v_i(A).
\]
\end{proof}

\begin{examplebox}[title={Worked Example: Movie ticket allocation with transfers}]
Suppose Joseph and Ivy both value a movie ticket, with willingness to pay
\[
WTP_J < WTP_I.
\]
If Joseph initially gets the ticket and no money changes hands, that allocation is Pareto efficient \emph{without} transfers: giving the ticket to Ivy would hurt Joseph.

Now allow money transfers. Give the ticket to Ivy and let Ivy pay Joseph an amount $x$ such that
\[
WTP_J \le x \le WTP_I.
\]
Then:
\begin{itemize}
    \item Joseph weakly prefers losing the ticket and receiving $x$;
    \item Ivy weakly prefers getting the ticket and paying $x$;
    \item at least one of them is strictly better off unless the inequalities are equalities at both ends.
\end{itemize}
So the initial allocation is Pareto improved upon.

Hence, once transfers are allowed, Pareto efficiency requires the ticket to go to the person with the highest willingness to pay.
\end{examplebox}

\subsection{Utilitarian Allocations}

Since transfers only change the distribution of utility, the relevant measure of the size of the ``utility pie'' generated by an allocation is the sum of $v_i(A)$ across individuals.

\begin{definition}{Utilitarian allocation}{ch11-utilitarian-allocation}
An allocation $A$ is \emph{utilitarian} if it maximizes
\[
\sum_{i=1}^n v_i(A)
\]
over all feasible allocations.
\end{definition}

Equivalently, society weakly prefers allocation $A$ to allocation $B$ under utilitarianism if and only if
\[
\sum_{i=1}^n v_i(A)\ge \sum_{i=1}^n v_i(B).
\]

\begin{examtip}
In this chapter, utilitarianism is applied to \emph{allocations}, not policies. Money transfers are dealt with separately through the vector $t$.
\end{examtip}

\subsection{Pareto Efficiency of Policies with Money Transfers}

We now define Pareto efficiency for policies, subject to balanced transfers.

\begin{definition}{Pareto improvement and Pareto efficiency for policies}{ch11-policy-pe}
A policy $(B,t')$ Pareto improves on policy $(A,t)$ if
\[
U_i(B,t_i') \ge U_i(A,t_i)
\quad \text{for all } i,
\]
and
\[
U_j(B,t_j') > U_j(A,t_j)
\quad \text{for some } j.
\]
A policy is \emph{Pareto efficient} if there is no other feasible policy with balanced transfers that Pareto improves on it.
\end{definition}

The key result of the lecture is that under quasi-linear utility and balanced transfers, Pareto-efficient policies and utilitarian allocations coincide.

\begin{theorem}{Utilitarianism and Pareto efficiency with money transfers}{ch11-main}
Suppose:
\begin{itemize}
    \item each individual has quasi-linear utility
    \[
    U_i(A,t_i)=v_i(A)+t_i;
    \]
    \item only balanced transfer schemes are allowed:
    \[
    \sum_{i=1}^n t_i=0.
    \]
\end{itemize}
Then a policy $(A,t)$ is Pareto efficient for some balanced transfer scheme $t$ if and only if $A$ is a utilitarian allocation.
\end{theorem}

\begin{proof}
We prove both directions by contraposition.

\paragraph{Direction 1.} If $A$ is utilitarian, then $(A,t)$ is Pareto efficient for some $t$.

It is enough to prove the contrapositive:
\[
\text{If $(A,t)$ is not Pareto efficient for any balanced $t$, then $A$ is not utilitarian.}
\]
Take any balanced $t$. If $(A,t)$ is not Pareto efficient, then there exists another balanced policy $(B,t')$ such that
\[
v_i(B)+t_i' \ge v_i(A)+t_i
\quad \text{for all } i,
\]
with strict inequality for at least one individual. Summing across all individuals gives
\[
\sum_{i=1}^n \bigl(v_i(B)+t_i'\bigr)
>
\sum_{i=1}^n \bigl(v_i(A)+t_i\bigr).
\]
Because both transfer schemes are balanced,
\[
\sum_{i=1}^n t_i'=\sum_{i=1}^n t_i=0,
\]
so
\[
\sum_{i=1}^n v_i(B) > \sum_{i=1}^n v_i(A).
\]
Therefore $A$ does not maximize total utility and is not utilitarian.

\paragraph{Direction 2.} If $(A,t)$ is Pareto efficient for some balanced $t$, then $A$ is utilitarian.

Again prove the contrapositive:
\[
\text{If $A$ is not utilitarian, then for any balanced $t$, $(A,t)$ is not Pareto efficient.}
\]
If $A$ is not utilitarian, then there exists some feasible allocation $B$ such that
\[
\sum_{i=1}^n v_i(B) > \sum_{i=1}^n v_i(A).
\]
Take any balanced transfer scheme $t$. Define a transfer vector
\[
\hat t_i = v_i(A)+t_i-v_i(B).
\]
Then for each individual $i$,
\[
U_i(B,\hat t_i)=v_i(B)+\hat t_i = v_i(A)+t_i = U_i(A,t_i).
\]
So each individual is indifferent between $(A,t)$ and $(B,\hat t)$.

Now sum these transfers:
\[
\sum_{i=1}^n \hat t_i
=
\sum_{i=1}^n v_i(A)+\sum_{i=1}^n t_i-\sum_{i=1}^n v_i(B)
=
\sum_{i=1}^n v_i(A)-\sum_{i=1}^n v_i(B)
<0,
\]
since $t$ is balanced and $B$ yields strictly higher total utility. Thus $\hat t$ runs a budget surplus. Let
\[
S = -\sum_{i=1}^n \hat t_i > 0.
\]
Redistribute this surplus equally by defining
\[
t_i'=\hat t_i + \frac{S}{n}.
\]
Then
\[
\sum_{i=1}^n t_i' = 0,
\]
so $t'$ is balanced. Also,
\[
U_i(B,t_i') = U_i(B,\hat t_i)+\frac{S}{n}
= U_i(A,t_i)+\frac{S}{n}
> U_i(A,t_i)
\]
for every individual $i$. Hence $(B,t')$ Pareto improves on $(A,t)$, so $(A,t)$ is not Pareto efficient.

This proves both directions.
\end{proof}

\begin{theorem}{Any balanced transfer supports a utilitarian allocation as Pareto efficient}{ch11-any-transfer}
Under the assumptions of Theorem \ref{thm:ch11-main}, if $A$ is utilitarian, then
\[
(A,t)
\]
is Pareto efficient for \emph{every} balanced transfer scheme $t$.
\end{theorem}

\begin{proof}
If some balanced policy $(B,t')$ Pareto improved on $(A,t)$, then by the summation argument in the proof of Theorem \ref{thm:ch11-main},
\[
\sum_{i=1}^n v_i(B) > \sum_{i=1}^n v_i(A),
\]
contradicting that $A$ is utilitarian.
\end{proof}

\subsection{Interpretation}

The theorem shows that, under the lecture's assumptions, the policy problem splits cleanly into two parts.

\begin{policynote}[title={Policy Interpretation: Efficiency then distribution}]
Under quasi-linear utility and balanced transfers:
\begin{enumerate}
    \item choose the allocation that maximizes
    \[
    \sum_i v_i(A),
    \]
    i.e. the largest utility pie;
    \item then use money transfers to determine how that pie is shared.
\end{enumerate}
So any non-utilitarian allocation is Pareto dominated by some policy built from a utilitarian allocation.
\end{policynote}

Graphically, the lecture interprets this as follows:
\begin{itemize}
    \item without transfers, different allocations generate a utility possibility frontier;
    \item with balanced transfers and quasi-linear utility, the frontier expands along lines of slope $-1$, because utility can be shifted across individuals through money;
    \item the outer frontier is attained by utilitarian allocations.
\end{itemize}

\begin{examtip}
The strongest intuition is:
\begin{quote}
Different allocations determine the \emph{size} of the utility pie. Transfers determine only the \emph{division} of that pie.
\end{quote}
Under quasi-linearity, any division of a smaller pie is Pareto dominated by some division of a larger pie.
\end{examtip}

\subsection{Worked Review Examples}

\begin{examplebox}[title={Worked Example: Which utility forms satisfy the utilitarianism argument?}]
Suppose policies involve allocations $(x_1^i,x_2^i)$ and transfers $t_i$. In which of the following cases does the lecture's utilitarianism argument apply?

\[
\text{(a)}\quad U_i = x_1^i x_2^i + t_i
\]
\[
\text{(b)}\quad U_i = x_1^i x_2^i
\]
\[
\text{(c)}\quad U_i = x_1^i x_2^i t_i
\]
\[
\text{(d)}\quad U_i = 0.5x_1^i + 0.5x_2^i + 0.5 t_i
\]

The argument requires utility to be quasi-linear in money, so that money enters additively with a constant marginal utility.

\begin{itemize}
    \item \textbf{(a)} Yes. This is of the form
    \[
    U_i = v_i(x)+t_i
    \]
    with
    \[
    v_i(x)=x_1^i x_2^i.
    \]
    \item \textbf{(b)} No. There is no monetary-transfer term.
    \item \textbf{(c)} No. Money enters multiplicatively, so utility is not quasi-linear in money.
    \item \textbf{(d)} Yes. This is still quasi-linear in money:
    \[
    U_i = v_i(x) + \alpha t_i
    \]
    with constant $\alpha=0.5>0$. Since the coefficient on money is constant and common, the same logic applies after a positive rescaling of utility.
\end{itemize}

Therefore the correct choices are
\[
\boxed{\text{(a) and (d).}}
\]
\end{examplebox}

\begin{examplebox}[title={Worked Example: Which allocations can be part of a Pareto-efficient policy?}]
Consider a two-person, one-good economy with balanced money transfers $(t_A,t_B)$ and utilities
\[
U_A=x_A+t_A,
\qquad
U_B=x_B+t_B.
\]
Suppose the feasible allocations are:
\[
\text{(a)}\ (x_A,x_B)=(4,4),
\qquad
\text{(b)}\ (x_A,x_B)=(5,3),
\qquad
\text{(c)}\ (x_A,x_B)=(3,5).
\]

Because
\[
v_A(x)=x_A,
\qquad
v_B(x)=x_B,
\]
the utilitarian sum for each allocation is
\[
x_A+x_B.
\]
For all three allocations,
\[
x_A+x_B=8.
\]
Hence all three allocations maximize the utilitarian objective over the listed feasible set.

By Theorem \ref{thm:ch11-main}, each utilitarian allocation can be part of a Pareto-efficient policy with some balanced transfers. In fact, by Theorem \ref{thm:ch11-any-transfer}, any balanced transfers work.

So the correct answer is
\[
\boxed{\text{(a), (b), and (c).}}
\]
\end{examplebox}

\subsection{Critical Assumptions and Limitations}

The lecture emphasizes two critical assumptions.

\begin{definition}{Critical assumptions behind the utilitarian result}{ch11-assumptions}
The utilitarian-efficiency argument requires:
\begin{enumerate}
    \item \textbf{Quasi-linear preferences in money:}
    \[
    U_i(A,t_i)=v_i(A)+t_i.
    \]
    \item \textbf{Correct implementation of balanced transfers:} society must be able to choose and enforce the needed transfer scheme after selecting the utilitarian allocation.
\end{enumerate}
\end{definition}

If these fail, then the theorem may no longer hold.

\begin{commonmistake}
Do not say that utilitarianism is always the best policy rule. The lecture's result is explicitly conditional:
\[
\text{quasi-linear utility} + \text{balanced transfers} + \text{correct implementation}.
\]
Without these, the equivalence between Pareto-efficient policies and utilitarian allocations can fail.
\end{commonmistake}

\begin{policynote}
Even when the utilitarian allocation is clear, society may still disagree over transfers. The theorem resolves the efficient \emph{allocation} question, but not the distributive question of how the gains from that allocation should be shared.
\end{policynote}

\subsection{Chapter Summary}

\begin{keyformula}[title={Key Formula: Lecture 11 Summary}]
\begin{align*}
A \to B &\Longleftrightarrow \neg B \to \neg A, \\
U_i(A,t_i) &= v_i(A)+t_i, \\
\sum_{i=1}^n t_i &= 0 \qquad \text{(balanced transfers)}, \\
\text{WTP}_i(A \to B) &= v_i(B)-v_i(A), \\
A \text{ utilitarian}
&\iff
\sum_{i=1}^n v_i(A) \ge \sum_{i=1}^n v_i(B)
\quad \text{for all feasible } B.
\end{align*}
Under quasi-linear utility and balanced transfers:
\[
(A,t) \text{ is Pareto efficient for some balanced } t
\iff
A \text{ is utilitarian.}
\]
Indeed, if $A$ is utilitarian, then
\[
(A,t)
\]
is Pareto efficient for every balanced transfer vector $t$.
\end{keyformula}

\begin{examtip}
The highest-yield takeaways from this chapter are:
\begin{itemize}
    \item know the difference between a conditional statement, its contrapositive, and its converse;
    \item know the difference between an allocation and a policy;
    \item know why balanced transfers are imposed;
    \item know how quasi-linear utility makes willingness to pay equal to utility differences;
    \item know the theorem:
    \[
    \text{Pareto-efficient policy with transfers}
    \iff
    \text{utilitarian allocation};
    \]
    \item know that the result depends critically on quasi-linear utility and the feasibility of implementing the right transfers.
\end{itemize}
\end{examtip}
\newpage
% chapters/ch12_part2_summary.tex

\section{Part II Summary}

Part II develops a sequence of welfare frameworks that become progressively broader. Lecture 7 begins at the level of the individual and asks how utility changes can be measured in money. Lecture 8 aggregates consumer and producer welfare in a single market and studies taxes, tax incidence, and deadweight loss. Lecture 9 then moves beyond partial equilibrium to the exchange economy, where Pareto efficiency and competitive equilibrium become the central concepts. Lecture 10 asks how society should rank alternatives or utility outcomes, introducing social decision mechanisms, social welfare functions, and the utility possibility frontier. Finally, Lecture 11 shows that once utility is quasi-linear in money and balanced transfers are feasible, Pareto-efficient policies correspond to utilitarian allocations.

This summary is designed for revision rather than first exposure. The emphasis is therefore on the main objects, the highest-yield results, the links across topics, and the standard patterns behind common exam questions.

\subsection{Part II at a Glance}

\begin{policynote}[title={Part II in one conceptual chain}]
The broad logic of Part II is:
\begin{enumerate}
    \item measure individual welfare in money;
    \item aggregate welfare in a market;
    \item generalize efficiency beyond a single market;
    \item distinguish efficiency from social choice;
    \item show when transfers make utilitarianism the efficiency criterion for policies.
\end{enumerate}
In short, Part II moves from \emph{how much one person gains or loses}, to \emph{how a market performs}, to \emph{which allocations are efficient}, to \emph{which efficient allocations society may want}, and finally to \emph{when money transfers let us separate efficiency from distribution}.
\end{policynote}

\subsection{Core Frameworks and Definitions}

\subsubsection{Monetary welfare at the individual level}

Lecture 7 asks how utility changes can be expressed in dollars.

\begin{definition}{Core monetary welfare measures}{ch12-monetary-measures}
The main individual monetary welfare measures are:
\begin{itemize}
    \item \emph{reservation price}: the maximum amount a consumer is willing to pay for a good or marginal unit;
    \item \emph{consumer surplus}:
    \[
    CS=\int_0^{x(p)} P(q)\,dq-p\,x(p);
    \]
    \item \emph{compensating variation (CV)}: the income needed at the \emph{new} prices to restore the \emph{old} utility level;
    \item \emph{equivalent variation (EV)}: the income removable at the \emph{old} prices that leaves the consumer at the \emph{new} utility level;
    \item \emph{producer surplus}:
    \[
    PS=\int_0^q \bigl(p-MC(\tilde q)\bigr)\,d\tilde q.
    \]
\end{itemize}
\end{definition}

The key distinction is that consumer surplus is usually convenient but approximate, whereas CV and EV are exact welfare measures for price changes. Consumer surplus becomes exact under quasi-linear utility.

\subsubsection{Aggregate welfare in a single market}

Lecture 8 aggregates consumer and producer welfare.

\begin{definition}{Single-market social welfare}{ch12-market-sw}
In the standard one-good market with inverse demand \(P_D(Q)\) and inverse supply \(P_S(Q)\), social welfare is
\[
SW(Q)=CS(Q)+PS(Q)=\int_0^Q \bigl(P_D(q)-P_S(q)\bigr)\,dq.
\]
The competitive quantity maximizes this expression under the standard downward-sloping demand and upward-sloping supply assumptions.
\end{definition}

This framework is still partial equilibrium: it is powerful for one-market tax questions, but it is not yet the general notion of efficiency used in exchange economies.

\subsubsection{Pareto efficiency in the exchange economy}

Lecture 9 replaces \(CS+PS\) with Pareto efficiency as the central efficiency concept.

\begin{definition}{Feasibility, Pareto improvement, and Pareto efficiency}{ch12-pe-core}
In a pure exchange economy:
\begin{itemize}
    \item an allocation is \emph{feasible} if total consumption equals total endowment good by good;
    \item a \emph{Pareto improvement} makes everyone weakly better off and at least one person strictly better off;
    \item an allocation is \emph{Pareto efficient} if no feasible Pareto improvement exists.
\end{itemize}
In the two-good, two-consumer case, the Edgeworth box is the main graphical tool, and at a smooth interior Pareto-efficient point,
\[
MRS_A=MRS_B.
\]
\end{definition}

This chapter separates efficiency from fairness: a Pareto-efficient allocation need not be envy-free or otherwise equitable.

\subsubsection{Social choice and social welfare}

Lecture 10 asks how society ranks alternatives or utility outcomes.

\begin{definition}{SDM, SWF, and social optimum}{ch12-sdm-swf}
A \emph{social decision mechanism} (SDM) takes individuals' rankings over alternatives as input and outputs a social ranking. A \emph{social welfare function} (SWF) takes utility levels as input and assigns a social value
\[
W(u_1,\dots,u_n).
\]
A \emph{social optimum} is a feasible allocation that maximizes the chosen SWF over the feasible utility set.
\end{definition}

This is the chapter where the course explicitly turns from positive efficiency analysis to normative aggregation.

\subsubsection{Policies with transfers}

Lecture 11 adds money transfers to the allocation problem.

\begin{definition}{Allocation, policy, and balanced transfers}{ch12-policy-transfers}
An \emph{allocation} specifies who gets the goods. A \emph{policy} specifies both an allocation \(A\) and transfers \(t=(t_1,\dots,t_n)\). Transfers are \emph{balanced} if
\[
\sum_{i=1}^n t_i=0.
\]
Under quasi-linear utility,
\[
U_i(A,t_i)=v_i(A)+t_i,
\]
so money can redistribute utility directly across individuals.
\end{definition}

This is the final conceptual step of Part II: once balanced transfers are feasible and utility is quasi-linear in money, one can separate the choice of the allocation from the choice of who ultimately bears the gains and losses.

\subsection{Key Results from Part II}

\begin{keyformula}[title={Key Formula: Part II Master Summary}]
\begin{align*}
CS &= \int_0^{x(p)} P(q)\,dq-p\,x(p), \\
V(p,I) &= \max_x U(x)\ \text{s.t.}\ p\cdot x\le I, \\
V(\text{new prices},I+CV) &= V(\text{old prices},I), \\
V(\text{old prices},I-EV) &= V(\text{new prices},I), \\
SW(Q) &= \int_0^Q \bigl(P_D(q)-P_S(q)\bigr)\,dq, \\
DWL &= \int_{Q'}^{Q^*} \bigl(P_D(q)-P_S(q)\bigr)\,dq, \\
\text{interior Pareto efficiency} &\Rightarrow MRS_A=MRS_B, \\
\text{competitive equilibrium} &\Rightarrow \text{Pareto efficiency}, \\
\text{social optimum} &\Rightarrow \text{Pareto efficiency}, \\
W(u_A,u_B)=\gamma_A u_A+\gamma_B u_B &\Rightarrow \text{iso-welfare slope } -\frac{\gamma_A}{\gamma_B}, \\
U_i(A,t_i)=v_i(A)+t_i,\ \sum_i t_i=0
&\Rightarrow
(A,t)\text{ is Pareto efficient} \iff A\text{ is utilitarian.}
\end{align*}
\end{keyformula}

The main examinable results can be organized as follows.

\begin{theorem}{High-yield results across Lectures 7--11}{ch12-main-results}
The central results of Part II are:
\begin{enumerate}
    \item Under quasi-linear utility,
    \[
    U(x,m)=v(x)+m,
    \]
    demand for \(x\) has no income effect and consumer surplus is an exact welfare measure.
    \item In the standard single-good market, the competitive quantity maximizes
    \[
    CS+PS.
    \]
    Taxes reduce quantity, create a wedge between buyer and seller prices, and generally generate deadweight loss.
    \item In the exchange economy, a competitive equilibrium allocation is Pareto efficient, and under the standard convexity assumptions any Pareto-efficient allocation can be decentralized after suitable redistribution of endowments.
    \item For any increasing SWF, a social optimum is Pareto efficient; conversely, under the regularity assumptions emphasized in the course, each Pareto-efficient point can be supported by some SWF.
    \item No non-dictatorial SDM can simultaneously satisfy Universal Domain, the Pareto property, and IIA while always producing a transitive social ranking.
    \item Under quasi-linear utility in money and balanced transfers, Pareto-efficient policies are exactly those built from utilitarian allocations.
\end{enumerate}
\end{theorem}

\begin{examtip}
A useful way to memorize the architecture of Part II is to attach one question to each lecture:
\begin{itemize}
    \item \textbf{Lecture 7:} how do I measure welfare?
    \item \textbf{Lecture 8:} what quantity is efficient in a market?
    \item \textbf{Lecture 9:} what allocation is Pareto efficient?
    \item \textbf{Lecture 10:} which Pareto-efficient point should society choose?
    \item \textbf{Lecture 11:} when do transfers make utilitarianism the relevant efficiency criterion?
\end{itemize}
\end{examtip}

\subsection{Cross-Topic Comparisons}

\subsubsection{Consumer surplus versus EV and CV}

Lecture 7 contains one of the most important distinctions in the course.

\begin{itemize}
    \item \(CS\) is computed from Marshallian demand and is usually easy to obtain.
    \item \(EV\) and \(CV\) are defined through indirect utility and are exact.
    \item In general,
    \[
    \min\{|CV|,|EV|\} \le |\Delta CS| \le \max\{|CV|,|EV|\}.
    \]
    \item Under quasi-linear utility, all three coincide.
\end{itemize}

This distinction matters again in Lecture 11, because quasi-linearity is exactly the assumption that turns willingness to pay into a clean monetary measure of utility difference.

\subsubsection{Market welfare versus Pareto efficiency}

Lecture 8 and Lecture 9 use different welfare languages.

\begin{itemize}
    \item In Lecture 8, welfare is
    \[
    CS+PS,
    \]
    which is especially useful in a single-market demand--supply diagram.
    \item In Lecture 9, welfare is not summarized by a single market-area formula. Instead, efficiency means the absence of feasible Pareto improvements.
\end{itemize}

The conceptual link is that both frameworks are trying to identify unrealized gains from trade. In Lecture 8 this appears as lost surplus triangles. In Lecture 9 it appears as allocations inside the Edgeworth-box improvement lens rather than on the contract curve.

\subsubsection{Efficiency versus fairness versus social choice}

Part II repeatedly warns against collapsing these concepts into one another.

\begin{itemize}
    \item A Pareto-efficient allocation need not be fair.
    \item An envy-free allocation need not be the competitive-equilibrium allocation from the given endowment.
    \item A social optimum depends on the chosen SWF.
    \item An SDM need not produce a transitive or stable ranking, as Condorcet cycles show.
\end{itemize}

So the course does \emph{not} say that efficiency settles all normative questions. Instead, it shows exactly where efficiency analysis ends and social-choice or distributive judgments begin.

\subsubsection{Redistribution before trade versus transfers after allocation}

Lecture 9 and Lecture 11 solve a similar problem in different ways.

\begin{itemize}
    \item In Lecture 9, the Second Welfare Theorem says that if society wants a particular Pareto-efficient allocation, it may be supported by first redistributing endowments and then allowing competitive trade.
    \item In Lecture 11, if utility is quasi-linear in money, society can instead choose the utilitarian allocation and then use balanced money transfers to determine the final distribution of welfare.
\end{itemize}

The common logic is the separation of \emph{efficiency} from \emph{distribution}. The technical difference is that Lecture 9 works through endowment redistribution in the exchange economy, whereas Lecture 11 works through utility-transferable money.

\subsubsection{SDMs versus SWFs}

Lecture 10 introduces two very different aggregation objects.

\begin{itemize}
    \item An SDM uses \emph{ordinal rankings} over alternatives.
    \item An SWF uses \emph{utility levels} and therefore supports geometric analysis using the utility possibility frontier.
\end{itemize}

Arrow's theorem is about SDMs on ordinal preference profiles. It does not directly invalidate the SWF analysis used later in Lecture 10 or in Lecture 11.

\subsection{Standard Problem Types and Solution Patterns}

\begin{examtip}[title={Method Pattern: Welfare-measure questions}]
If the question is about reservation price, \(CS\), \(EV\), or \(CV\):
\begin{enumerate}
    \item identify whether the object is discrete, continuous, or a price change;
    \item if exact monetary welfare is needed, derive indirect utility and solve for \(EV\) or \(CV\);
    \item check whether utility is quasi-linear, because that may let you use consumer surplus exactly;
    \item keep the reference point straight:
    \[
    CV=\text{new prices, old utility},
    \qquad
    EV=\text{old prices, new utility}.
    \]
\end{enumerate}
\end{examtip}

\begin{examtip}[title={Method Pattern: Tax and deadweight-loss questions}]
If the question is about a per-unit or ad-valorem tax:
\begin{enumerate}
    \item define buyer price \(p_b\) and seller price \(p_s\);
    \item write the wedge condition;
    \item solve demand equals supply using the correct post-tax prices;
    \item compare with the no-tax benchmark;
    \item compute tax revenue, incidence, and deadweight loss separately.
\end{enumerate}
\end{examtip}

\begin{examtip}[title={Method Pattern: Exchange-economy questions}]
If the question is about Pareto efficiency or competitive equilibrium in an exchange economy:
\begin{enumerate}
    \item write total resources and feasibility conditions;
    \item derive each consumer's demand from utility and endowment wealth;
    \item solve for the equilibrium price ratio using market clearing;
    \item if checking efficiency, use either the improvement-lens logic or
    \[
    MRS_A=MRS_B
    \]
    at smooth interior allocations;
    \item if fairness is asked, test envy-freeness separately.
\end{enumerate}
\end{examtip}

\begin{examtip}[title={Method Pattern: Social-choice questions}]
If the question is about an SDM or SWF:
\begin{enumerate}
    \item distinguish immediately whether the object is an SDM or an SWF;
    \item for an SDM, compute the social ranking methodically and then test properties such as Pareto, Universal Domain, and IIA;
    \item for an SWF over a utility possibility set, look for the highest attainable iso-welfare curve on the frontier;
    \item for weighted utilitarianism in two dimensions, match the frontier slope to
    \[
    -\frac{\gamma_A}{\gamma_B}.
    \]
\end{enumerate}
\end{examtip}

\begin{examtip}[title={Method Pattern: Lecture 11 utilitarianism questions}]
If the question involves money transfers:
\begin{enumerate}
    \item distinguish allocation from policy;
    \item check that transfers are balanced;
    \item check whether utility is quasi-linear in money:
    \[
    U_i(A,t_i)=v_i(A)+t_i;
    \]
    \item compute willingness to pay as
    \[
    v_i(B)-v_i(A);
    \]
    \item identify the allocation that maximizes
    \[
    \sum_i v_i(A).
    \]
\end{enumerate}
\end{examtip}

\subsection{Common Mistakes and Exam Traps}

\begin{commonmistake}
Do not confuse the reservation price curve with the inverse demand curve. They coincide only under special assumptions, especially quasi-linear utility.
\end{commonmistake}

\begin{commonmistake}
Do not confuse the statutory side of a quantity tax with its economic incidence. The side that formally remits the tax need not be the side that bears more of the burden.
\end{commonmistake}

\begin{commonmistake}
Do not say that a Pareto-efficient allocation is automatically fair. Pareto efficiency rules out mutually beneficial reallocation; it does not say anything by itself about envy, equality, or justice.
\end{commonmistake}

\begin{commonmistake}
Do not confuse an SDM with an SWF. Arrow's theorem concerns SDMs defined on ordinal rankings, not every possible welfare criterion used in economics.
\end{commonmistake}

\begin{commonmistake}
Do not confuse the \emph{contrapositive}
\[
A \to B
\quad \Longleftrightarrow \quad
\neg B \to \neg A
\]
with the \emph{converse}
\[
B \to A.
\]
Lecture 11 uses contraposition repeatedly, and losing this distinction can collapse an otherwise correct proof.
\end{commonmistake}

\begin{commonmistake}
Do not confuse an allocation with a policy in Lecture 11. Pareto efficiency there is applied to policies \((A,t)\), not merely to allocations \(A\) in isolation.
\end{commonmistake}

\subsection{Integrated Worked Review Examples}

\begin{examplebox}[title={Worked Review Example 1: When is consumer surplus exact?}]
Suppose a consumer has quasi-linear utility
\[
U(x,m)=2\sqrt{x}+m,
\]
income is high enough that the interior optimum is affordable, and the price of \(x\) rises from
\[
p'=1 \quad \text{to} \quad p''=2.
\]

Because utility is quasi-linear in money, the optimal quantity solves
\[
\frac{d}{dx}\bigl(2\sqrt{x}-px\bigr)=0
\quad \Longrightarrow \quad
\frac{1}{\sqrt{x}}=p.
\]
Hence
\[
x(p)=\frac{1}{p^2}.
\]
So
\[
x(1)=1,
\qquad
x(2)=\frac14.
\]

The loss in consumer surplus from the price increase is
\[
\Delta CS_{\text{loss}}
=
\int_{1}^{2} x(p)\,dp
=
\int_1^2 \frac{1}{p^2}\,dp
=
\left[-\frac{1}{p}\right]_1^2
=
\frac12.
\]

Under quasi-linear utility, consumer surplus is exact, so
\[
\boxed{\Delta CS=EV=CV=\frac12.}
\]

This is the key revisional lesson from Lecture 7: whenever the utility function is quasi-linear in money, the distinction between \(CS\), \(EV\), and \(CV\) disappears.
\end{examplebox}

\begin{examplebox}[title={Worked Review Example 2: Linear tax incidence and deadweight loss}]
Suppose
\[
Q^D=20-p_b,
\qquad
Q^S=p_s,
\]
and a per-unit tax of
\[
t=4
\]
is imposed.

\textbf{Without tax:}
\[
20-p=p
\quad \Longrightarrow \quad
p^*=10,
\qquad
Q^*=10.
\]

\textbf{With tax:}
\[
p_b=p_s+4,
\qquad
20-p_b=p_s.
\]
Substituting gives
\[
20-(p_s+4)=p_s
\quad \Longrightarrow \quad
16=2p_s
\quad \Longrightarrow \quad
p_s=8.
\]
Hence
\[
p_b=12,
\qquad
Q'=8.
\]

So buyers' price rises by
\[
12-10=2,
\]
while sellers' price falls by
\[
10-8=2.
\]
Thus the tax burden is shared equally in this case.

Deadweight loss is
\[
DWL=\frac12 \times 4 \times (10-8)=4.
\]

Therefore
\[
\boxed{p_b=12,\quad p_s=8,\quad Q'=8,\quad DWL=4.}
\]

This example consolidates the Lecture 8 pattern: write the wedge condition carefully, solve for the post-tax equilibrium, then keep incidence and deadweight loss conceptually separate.
\end{examplebox}

\begin{examplebox}[title={Worked Review Example 3: Competitive equilibrium, efficiency, and fairness}]
Consider a two-good, two-consumer exchange economy with total resources
\[
(12,12),
\]
endowments
\[
\omega^A=(4,8),
\qquad
\omega^B=(8,4),
\]
and utilities
\[
U_A(x_1,x_2)=x_1^{1/2}x_2^{1/2},
\qquad
U_B(x_1,x_2)=x_1^{1/2}x_2^{1/2}.
\]

Normalize prices so that
\[
p_2=1,
\qquad
p_1=p.
\]
Each consumer has Cobb--Douglas demand with equal expenditure shares, so
\[
x_1^i=\frac{m_i}{2p},
\qquad
x_2^i=\frac{m_i}{2},
\]
where
\[
m_A=4p+8,
\qquad
m_B=8p+4.
\]

Market clearing for good 1 requires
\[
\frac{4p+8}{2p}+\frac{8p+4}{2p}=12,
\]
so
\[
\frac{12p+12}{2p}=12
\quad \Longrightarrow \quad
12p+12=24p
\quad \Longrightarrow \quad
p=1.
\]

Thus
\[
m_A=m_B=12,
\]
and each consumer chooses
\[
x^{A*}=x^{B*}=(6,6).
\]

Hence the competitive equilibrium allocation is
\[
\boxed{x^{A*}=x^{B*}=(6,6).}
\]

Because this is a competitive equilibrium, Lecture 9 implies it is Pareto efficient. It is also envy-free here because both consumers receive the same bundle and have identical preferences. The revisional point is that these are distinct properties:
\[
\text{competitive equilibrium} \Rightarrow \text{Pareto efficiency},
\]
but envy-freeness is an additional fairness test, not part of the First Welfare Theorem itself.
\end{examplebox}

\begin{examplebox}[title={Worked Review Example 4: Condorcet cycles under pairwise majority voting}]
Suppose three voters have preferences
\[
\text{Voter 1: } X \succ Z \succ Y,
\qquad
\text{Voter 2: } Z \succ Y \succ X,
\qquad
\text{Voter 3: } Y \succ X \succ Z.
\]

Pairwise comparisons give:
\begin{itemize}
    \item \(Y\) beats \(X\), because voters 2 and 3 prefer \(Y\) to \(X\);
    \item \(X\) beats \(Z\), because voters 1 and 3 prefer \(X\) to \(Z\);
    \item \(Z\) beats \(Y\), because voters 1 and 2 prefer \(Z\) to \(Y\).
\end{itemize}

Therefore
\[
Y \succ^* X,
\qquad
X \succ^* Z,
\qquad
Z \succ^* Y.
\]

So the social ranking cycles:
\[
\boxed{Y \succ^* X \succ^* Z \succ^* Y.}
\]

This example captures the core Lecture 10 warning: a seemingly natural voting rule can fail to produce a transitive social ranking on an unrestricted domain.
\end{examplebox}

\begin{examplebox}[title={Worked Review Example 5: Utilitarian allocation with balanced transfers}]
Suppose a single indivisible ticket must be assigned to either Ann or Ben. Their willingness to pay is
\[
WTP_A=10,
\qquad
WTP_B=16.
\]

\textbf{Without transfers:} if Ann initially has the ticket, giving it to Ben harms Ann, so the initial allocation may already be Pareto efficient relative to allocations alone.

\textbf{With balanced transfers and quasi-linear utility:} giving the ticket to Ben and making Ben pay Ann an amount \(x\) with
\[
10 \le x \le 16
\]
makes Ann weakly willing to give up the ticket and Ben weakly willing to receive it.

Hence any transfer in that interval supports a Pareto improvement, and the allocation that gives the ticket to Ben is the utilitarian one because it maximizes total willingness to pay.

So the efficient policy is built from the allocation
\[
\boxed{\text{ticket to Ben}}
\]
together with a balanced transfer \(x\) from Ben to Ann.

This is the core Lecture 11 insight: once quasi-linear money transfers are feasible, the efficient allocation sends the object to the person with the highest willingness to pay, while transfers determine distribution.
\end{examplebox}

\subsection{Lecture 11 Final Synthesis}

Lecture 11 is the endpoint of Part II, so it is worth stating its conceptual role explicitly. Chapters 7--10 gradually build the machinery needed for the final theorem.

\begin{itemize}
    \item Lecture 7 develops monetary welfare measures and shows why quasi-linearity matters.
    \item Lecture 8 shows how monetary measures can be aggregated into market welfare.
    \item Lecture 9 shows that efficiency can be defined without relying on a single-market surplus diagram.
    \item Lecture 10 distinguishes efficiency from the separate problem of social ranking.
    \item Lecture 11 then proves that under quasi-linear utility in money and balanced transfers, the efficient policy problem collapses to a utilitarian allocation problem.
\end{itemize}

\begin{policynote}[title={Efficiency then distribution}]
The final message of Part II is:
\begin{enumerate}
    \item choose the allocation that maximizes
    \[
    \sum_{i=1}^n v_i(A);
    \]
    \item then use balanced transfers to determine how the resulting utility pie is divided.
\end{enumerate}
Thus, under the Lecture 11 assumptions, different allocations determine the \emph{size} of the pie, while transfers determine only its \emph{division}.
\end{policynote}

\begin{examtip}
If an exam question asks why utilitarianism appears in Lecture 11, the correct answer is not simply ``because it is morally attractive.'' The lecture's actual argument is structural:
\begin{itemize}
    \item utility is quasi-linear in money;
    \item transfers are balanced;
    \item willingness to pay equals utility difference in money terms;
    \item therefore Pareto-efficient policies are exactly those based on utilitarian allocations.
\end{itemize}
\end{examtip}

\subsection{Part II Revision Checklist}

\begin{examtip}[title={What you should be able to do before the exam}]
By the end of revision for Part II, you should be able to:
\begin{enumerate}
    \item compute reservation prices, consumer surplus, producer surplus, \(EV\), and \(CV\), and state when consumer surplus is exact;
    \item solve linear tax problems with buyer price, seller price, tax revenue, incidence, and deadweight loss all clearly distinguished;
    \item interpret an Edgeworth box, identify Pareto improvements, and solve basic two-good, two-consumer competitive-equilibrium problems;
    \item explain the First and Second Welfare Theorems and state the assumptions under which they are used in the course;
    \item distinguish SDMs from SWFs, identify Condorcet-type failures, and interpret tangency on the utility possibility frontier;
    \item state Arrow's theorem accurately without overclaiming what it proves;
    \item distinguish allocations from policies in Lecture 11, compute willingness to pay under quasi-linear utility, and identify utilitarian allocations;
    \item explain clearly how Part II separates efficiency questions from distributional or social-choice questions.
\end{enumerate}
\end{examtip}

\begin{keyformula}[title={Key Formula: Part II Final Recall Sheet}]
\begin{align*}
CS &= \int_0^{x(p)} P(q)\,dq-px(p), \\
PS &= \int_0^q \bigl(p-MC(\tilde q)\bigr)\,d\tilde q, \\
V(\text{new},I+CV) &= V(\text{old},I), \\
V(\text{old},I-EV) &= V(\text{new},I), \\
SW(Q) &= \int_0^Q \bigl(P_D(q)-P_S(q)\bigr)\,dq, \\
DWL &= \int_{Q'}^{Q^*} \bigl(P_D(q)-P_S(q)\bigr)\,dq, \\
\text{interior PE} &\Rightarrow MRS_A=MRS_B, \\
\text{CE} &\Rightarrow \text{PE}, \\
\text{PE} &\Rightarrow \text{social optimum for some SWF}, \\
\text{weighted utilitarian slope} &= -\frac{\gamma_A}{\gamma_B}, \\
U_i(A,t_i)&=v_i(A)+t_i,\qquad \sum_i t_i=0, \\
(A,t)\text{ PE policy} &\iff A\text{ utilitarian.}
\end{align*}
\end{keyformula}









\appendix

\section{Formula \& Key Results Reference}
\begin{itemize}
    \item Chapter 1: Budget sets, direct and strict revealed preference, WARP.
    \item Part I: Slutsky decomposition, endowment-budget formulas, intertemporal constraints, asset-pricing identities, and uncertainty results.
    \item Part II: Welfare measures, surplus formulas, Pareto efficiency criteria, and social choice results.
\end{itemize}

\section{Notation}
\renewcommand{\arraystretch}{1.2}
\begin{center}
\begin{tabular}{|p{2.2cm}|p{5.0cm}|p{6.0cm}|}
\hline
\textbf{symbol} & \textbf{meaning} & \textbf{context} \\
\hline
$x \in \mathbb{R}^L_+$ & consumption bundle & Consumer theory chapters \\
\hline
$p \in \mathbb{R}^L_{++}$ & price vector & Consumer theory chapters \\
\hline
$m$ & expenditure / wealth level & Budget constraints \\
\hline
$B(p,m)$ & budget set & Choice under prices and expenditure \\
\hline
$\succeq, \succ, \sim$ & weak preference, strict preference, indifference & Preference theory \\
\hline
$R, P$ & directly revealed preference, strictly directly revealed preference & Revealed preference \\
\hline
\end{tabular}
\end{center}

\end{document}
